% Leçon 43 — Révisions générales — Chinois 1ère A — Cours signé OnBuch+
% Programme officiel MINESEC. Compiler avec : tectonic ZHA43-revision-revisions.tex
\newcommand{\DOCMATIERE}{Chinois}
\newcommand{\DOCNIVEAU}{1ère A}
\newcommand{\DOCMODULE}{Module 5 : Mass médias et télécommunication}
\newcommand{\DOCLECON}{ZHA43}
\newcommand{\DOCTITRE}{Leçon 43 — Révisions générales}
% OnBuch+ — préambule « cours signés » ENRICHI (compilé avec tectonic / XeLaTeX)
% Système visuel « Pop » d'OnBuch+ : fond crème (jamais blanc), bordures encre
% épaisses, ombres « dures » décalées, titres Archivo Black en capitales.
% Version MATHÉMATIQUES : graphes (pgfplots/tikz), boîtes pédagogiques
% (définition, propriété, méthode, attention, le savais-tu, exercice résolu,
% exercice, TP, bilan), page de garde et signature OnBuch+ ancrée en bas.
%
% Macros attendues dans main.tex AVANT \input{preamble} :
%   \DOCMATIERE  \DOCNIVEAU  \DOCMODULE  \DOCLECON (numéro)  \DOCTITRE
\documentclass[11pt]{article}

\usepackage{fontspec}
\usepackage[french]{babel} % français = langue principale ; le chinois via \zh{..} / textebox (xeCJK)
\usepackage{xeCJK}
\usepackage{pifont} % \ding{51} \ding{55} utilisés par les modèles
\usepackage[a4paper,margin=2cm,top=3.4cm,bottom=2.8cm,headheight=1.4cm,footskip=1.3cm]{geometry}
\usepackage{amsmath,amssymb,amsfonts}
\usepackage{mathtools} % \xmapsto, \coloneqq... souvent utilisés par les modèles
\usepackage{cancel} % simplifications barrées (bilans de forces)
% \abs{z} (module, valeur absolue) et \norm{u} (norme), utilisés par les modèles
\providecommand{\abs}{}\let\abs\relax\DeclarePairedDelimiter{\abs}{\lvert}{\rvert}
\providecommand{\norm}{}\let\norm\relax\DeclarePairedDelimiter{\norm}{\lVert}{\rVert}
\usepackage[table]{xcolor} % [table] : fournit aussi \rowcolors (alternance de lignes)
\usepackage{pagecolor}
\usepackage{tikz}
\usetikzlibrary{arrows.meta,positioning,calc,shapes.geometric,shapes.misc,shapes.arrows,decorations.pathmorphing,decorations.pathreplacing,decorations.markings,patterns,shadows,mindmap,trees,fit,backgrounds,matrix,intersections,angles,quotes}
\usepackage{pgfplots}
\usepackage[version=4]{mhchem} % équations nucléaires \ce{^{235}_{92}U}
\usepackage[european,siunitx]{circuitikz} % schémas électriques
\pgfplotsset{compat=1.18}
\usepgfplotslibrary{fillbetween,groupplots} % aires entre courbes (calcul intégral)
\usepackage{enumitem}
\usepackage{fancyhdr}
% [most] charge toutes les bibliothèques tcolorbox (listings, minted, raster,
% documentation...) et gonfle inutilement la mémoire TeX sur un document long
% (~300 boîtes) : on ne charge que ce qui est réellement utilisé.
\usepackage[skins,breakable]{tcolorbox}
\usepackage{graphicx}
\usepackage{colortbl}
\usepackage{booktabs}
\usepackage{tabularx}
\usepackage{diagbox} % case d'en-tête diagonale des tableaux à double entrée
% Colonnes extensibles centrée / gauche / droite (C, L, R), souvent utilisées par les modèles
\newcolumntype{C}{>{\centering\arraybackslash}X}
\newcolumntype{L}{>{\raggedright\arraybackslash}X}
\newcolumntype{R}{>{\raggedleft\arraybackslash}X}
\usepackage{array}
\usepackage{multicol}
\usepackage{multirow}
\usepackage{siunitx}
% parse-numbers=false : accepte aussi \SI{2,5 \times 10^{-3}}{...}
\sisetup{locale=FR,output-decimal-marker={,},group-minimum-digits=4,parse-numbers=false}
\DeclareSIUnit\ohm{\ensuremath{\symup{\Omega}}} % U+2126 absent de la police
\DeclareSIUnit\division{div} % réglages d'oscilloscope (ms/div, V/div)
\DeclareSIUnit\tr{tr} % tours (vitesse de rotation, ex. \SI{1500}{\tr\per\minute})
\DeclareSIUnit\molar{M} % concentration molaire (ex. \SI{3}{\molar}, \SI{1}{\micro\molar})
\DeclareSIUnit\an{an} % année (le modèle confond parfois avec \year, un compteur TeX) — ex. \SI{5}{\kilo\meter\per\an}
\DeclareSIUnit\FCFA{FCFA} % franc CFA (ex. \SI{500}{\FCFA\per\kilo\gram})
\DeclareSIUnit\degreeBrix{\textdegree Brix} % teneur en sucre d'un sirop/jus (réfractométrie)
\DeclareSIUnit\month{mois} % le modèle utilise parfois \month, un compteur TeX, comme unité de durée
\usepackage{qrcode}
% \degree : commande fréquemment utilisée par les modèles pour un angle ou une
% température en dehors de \SI{}{} (non fournie par les paquets chargés).
\providecommand{\degree}{^{\circ}}
% \qed (fin de démonstration), utilisé par les modèles sans amsthm
\providecommand{\qed}{\hfill$\square$}
% \barre : double barre des valeurs interdites dans un tableau de variations (tkz-tab)
\providecommand{\barre}{\ensuremath{\|}}
% environnement proof (amsthm non chargé) utilisé par les modèles
\ifdefined\proof\else\newenvironment{proof}[1][Démonstration]{\par\noindent\textit{#1.}\ }{\qed\par}\fi
\usepackage{lastpage}
\usepackage{hyperref}
\hypersetup{hidelinks}

% ============================================================
% Palette « Pop » OnBuch+ — mêmes hex que la classe P (plus_ui.dart)
% ============================================================
\definecolor{poporange}{HTML}{F59321}
\definecolor{popdark}{HTML}{E07A0C}
\definecolor{popcream}{HTML}{FFF6E8}
\definecolor{popink}{HTML}{1C1714}
\definecolor{poppurple}{HTML}{7A5AE0}
\definecolor{popgreen}{HTML}{2E9E6A}
\definecolor{poppink}{HTML}{E0457B}
\definecolor{popblue}{HTML}{2F7DE1}
\definecolor{popgold}{HTML}{C98A12}
\definecolor{popmuted}{HTML}{8A7F76}
\definecolor{popline}{HTML}{EADFD2}
\definecolor{popgray}{HTML}{8A7F76}
\colorlet{popgrey}{popgray}
% Alias tolérants (le modèle nomme parfois les couleurs différemment)
\colorlet{popwhite}{white}
\colorlet{popblack}{popink}
\colorlet{popred}{poppink}
\colorlet{popyellow}{popgold}
% Teintes claires (fonds de tableaux/encadrés) — le modèle nomme parfois
% les couleurs avec un suffixe L (ex. popblueL) sans les avoir définies.
\colorlet{poporangeL}{poporange!20}
\colorlet{popdarkL}{popdark!20}
\colorlet{poppurpleL}{poppurple!20}
\colorlet{popgreenL}{popgreen!20}
\colorlet{poppinkL}{poppink!20}
\colorlet{popblueL}{popblue!20}
\colorlet{popgoldL}{popgold!20}
\colorlet{popredL}{popred!20}
\colorlet{popgrayL}{popgray!20}
% Teintes claires pour fonds de tableaux / zones
\colorlet{popblueL}{popblue!10!white}
\colorlet{poporangeL}{poporange!14!white}
\colorlet{popgreenL}{popgreen!12!white}
\colorlet{poppurpleL}{poppurple!12!white}
\colorlet{poppinkL}{poppink!12!white}
\colorlet{popgoldL}{popgold!14!white}

\pagecolor{popcream}

% ============================================================
% Typographie
% ============================================================
\defaultfontfeatures{Ligatures=TeX}
\setmainfont{PlusJakartaSans-Medium.ttf}[Path=../../../fonts/,BoldFont=PlusJakartaSans-Bold.ttf]
% Symboles, API et emojis absents de la police principale : repli sur DejaVu Sans (caractères actifs)
\newfontface\symfont{DejaVuSans.ttf}[Path=../../../fonts/]
\makeatletter
\newcommand\symactive[1]{\catcode#1=13 \begingroup\lccode`\~=#1 \lowercase{\endgroup\def~}{{\symfont\char#1}}}
\newcommand\symblank[1]{\catcode#1=13 \begingroup\lccode`\~=#1 \lowercase{\endgroup\def~}{}}
\newcommand\symhyph[1]{\catcode#1=13 \begingroup\lccode`\~=#1 \lowercase{\endgroup\def~}{\mbox{-}}}
\newcount\symc
\newcommand\symrange[2]{\symc=#1 \@whilenum\symc<#2 \do{\expandafter\symactive\expandafter{\the\symc}\advance\symc by 1 }}
\newcommand\symblankrange[2]{\symc=#1 \@whilenum\symc<#2 \do{\expandafter\symblank\expandafter{\the\symc}\advance\symc by 1 }}
\makeatother
\symrange{"0250}{"02FF}\symactive{"2713}\symactive{"2714}\symactive{"2717}\symactive{"2718}\symactive{"2610}\symactive{"2611}\symactive{"2612}
\symrange{"2600}{"27BF}
\catcode"2705=13 \begingroup\lccode`\~="2705 \lowercase{\endgroup\def~}{{\symfont\char"2713}}\symblank{"26BD}\symblank{"26A1}
\symrange{"2460}{"257F}\symactive{"223C}\symactive{"2248}\symactive{"2260}
\symactive{"01F8}\symactive{"01F9}\symactive{"1E3E}\symactive{"1E3F}
\symhyph{"2011}\symhyph{"2E1A}\symblankrange{"1F300}{"1FAFF}\symblank{"FE0F}
% Caractères chinois : WenQuanYi Zen Hei (sous-ensemble GB2312 + pinyin), gras/italique simulés
\setCJKmainfont{WenQuanYiZenHei-subset.ttf}[Path=../../../fonts/,AutoFakeBold=3,AutoFakeSlant=0.2]
\xeCJKsetup{CJKspace=false}
% Pinyin : voyelles à caron (ǎ ǐ ǒ ǔ ǖ ǘ ǚ ǜ) absentes de la police principale -> caractères actifs
\newfontface\pyfont{WenQuanYiZenHei-subset.ttf}[Path=../../../fonts/]
\newcommand\pyactive[1]{\catcode#1=13 \begingroup\lccode`\~=#1 \lowercase{\endgroup\def~}{{\pyfont\char#1}}}
\pyactive{"01CD}\pyactive{"01CE}\pyactive{"01CF}\pyactive{"01D0}\pyactive{"01D1}\pyactive{"01D2}\pyactive{"01D3}\pyactive{"01D4}\pyactive{"01D5}\pyactive{"01D6}\pyactive{"01D7}\pyactive{"01D8}\pyactive{"01D9}\pyactive{"01DA}\pyactive{"01DB}\pyactive{"01DC}
% Maths en sans-serif (Fira Math) avec chiffres et lettres droites de Plus
% Jakarta, pour que les formules se fondent dans le texte.
\usepackage[math-style=ISO,bold-style=ISO]{unicode-math}
\setmathfont{FiraMath-Regular.otf}[Path=../../../fonts/]
\setmathfont{PlusJakartaSans-Medium.ttf}[Path=../../../fonts/,range={up/num,up/{Latin,latin},\mathup}]
\usepackage{icomma} % virgule décimale sans espace en mode math
% Caractères mathématiques Unicode absents des polices de texte (Plus Jakarta,
% Archivo Black) : rendus par leur équivalent mathématique, en texte comme en
% formule — sinon ils s'affichent en carré vide (ex. « Arithmétique dans ℤ »).
\usepackage{newunicodechar}
\newunicodechar{ℝ}{\ensuremath{\mathbb{R}}}
\newunicodechar{ℤ}{\ensuremath{\mathbb{Z}}}
\newunicodechar{ℂ}{\ensuremath{\mathbb{C}}}
\newunicodechar{ℕ}{\ensuremath{\mathbb{N}}}
\newunicodechar{ℚ}{\ensuremath{\mathbb{Q}}}
\newunicodechar{𝔻}{\ensuremath{\mathbb{D}}}
\newunicodechar{α}{\ensuremath{\alpha}}
\newunicodechar{β}{\ensuremath{\beta}}
\newunicodechar{γ}{\ensuremath{\gamma}}
\newunicodechar{δ}{\ensuremath{\delta}}
\newunicodechar{ε}{\ensuremath{\varepsilon}}
\newunicodechar{θ}{\ensuremath{\theta}}
\newunicodechar{λ}{\ensuremath{\lambda}}
\newunicodechar{μ}{\ensuremath{\mu}}
\newunicodechar{σ}{\ensuremath{\sigma}}
\newunicodechar{τ}{\ensuremath{\tau}}
\newunicodechar{φ}{\ensuremath{\varphi}}
\newunicodechar{ψ}{\ensuremath{\psi}}
\newunicodechar{ω}{\ensuremath{\omega}}
\newunicodechar{Σ}{\ensuremath{\Sigma}}
% majuscules produites par \MakeUppercase dans les titres (α-aminés -> Α)
\newunicodechar{Α}{\ensuremath{\alpha}}
\newunicodechar{Β}{\ensuremath{\beta}}
\newunicodechar{Γ}{\ensuremath{\Gamma}}
\newunicodechar{Δ}{\ensuremath{\Delta}}
\newunicodechar{Ε}{\ensuremath{\varepsilon}}
\newunicodechar{Θ}{\ensuremath{\Theta}}
\newunicodechar{Λ}{\ensuremath{\Lambda}}
\newunicodechar{Μ}{\ensuremath{\mu}}
\newunicodechar{Τ}{\ensuremath{\tau}}
\newunicodechar{Φ}{\ensuremath{\Phi}}
\newunicodechar{Ψ}{\ensuremath{\Psi}}
\newunicodechar{Ω}{\ensuremath{\Omega}}
\newunicodechar{↦}{\ensuremath{\mapsto}}
\newunicodechar{∈}{\ensuremath{\in}}
\newunicodechar{∉}{\ensuremath{\notin}}
\newunicodechar{∧}{\ensuremath{\wedge}}
\newunicodechar{⇔}{\ensuremath{\Leftrightarrow}}
\newunicodechar{⟺}{\ensuremath{\Longleftrightarrow}}
\newunicodechar{⇒}{\ensuremath{\Rightarrow}}
\newunicodechar{⟹}{\ensuremath{\Longrightarrow}}
\newunicodechar{∘}{\ensuremath{\circ}}
\newunicodechar{∪}{\ensuremath{\cup}}
\newunicodechar{∩}{\ensuremath{\cap}}
\newunicodechar{≡}{\ensuremath{\equiv}}
\newunicodechar{⊂}{\ensuremath{\subset}}
\newunicodechar{⊥}{\ensuremath{\perp}}
\newunicodechar{∃}{\ensuremath{\exists}}
\newunicodechar{∀}{\ensuremath{\forall}}
\newunicodechar{∛}{\ensuremath{\sqrt[3]{\,}}}
\newunicodechar{∜}{\ensuremath{\sqrt[4]{\,}}}
\newunicodechar{⃗}{\ensuremath{{}^{\to}}}
\newunicodechar{ⁿ}{\ifmmode^{n}\else\textsuperscript{\ensuremath{n}}\fi}
\newunicodechar{ˣ}{\ifmmode^{x}\else\textsuperscript{\ensuremath{x}}\fi}
\newunicodechar{ᵇ}{\ifmmode^{b}\else\textsuperscript{\ensuremath{b}}\fi}
\newunicodechar{ʳ}{\ifmmode^{r}\else\textsuperscript{\ensuremath{r}}\fi}
\newunicodechar{ᵘ}{\ifmmode^{u}\else\textsuperscript{\ensuremath{u}}\fi}
\newunicodechar{ʸ}{\ifmmode^{y}\else\textsuperscript{\ensuremath{y}}\fi}
\newunicodechar{ᵃ}{\ifmmode^{a}\else\textsuperscript{\ensuremath{a}}\fi}
\newunicodechar{ᵉ}{\ifmmode^{e}\else\textsuperscript{\ensuremath{e}}\fi}
\newunicodechar{ᵏ}{\ifmmode^{k}\else\textsuperscript{\ensuremath{k}}\fi}
\newunicodechar{ᵖ}{\ifmmode^{p}\else\textsuperscript{\ensuremath{p}}\fi}
\newunicodechar{ᴺ}{\ifmmode^{N}\else\textsuperscript{\ensuremath{N}}\fi}
\newunicodechar{ᶜ}{\ifmmode^{c}\else\textsuperscript{\ensuremath{c}}\fi}
\newunicodechar{ʰ}{\ifmmode^{h}\else\textsuperscript{\ensuremath{h}}\fi}
\newunicodechar{ᵠ}{\ifmmode^{q}\else\textsuperscript{\ensuremath{q}}\fi}
\newunicodechar{ₙ}{\ifmmode_{n}\else\textsubscript{\ensuremath{n}}\fi}
\newunicodechar{ₐ}{\ifmmode_{a}\else\textsubscript{\ensuremath{a}}\fi}
\newunicodechar{ᵢ}{\ifmmode_{i}\else\textsubscript{\ensuremath{i}}\fi}
\newunicodechar{ᵣ}{\ifmmode_{r}\else\textsubscript{\ensuremath{r}}\fi}
\newunicodechar{ₖ}{\ifmmode_{k}\else\textsubscript{\ensuremath{k}}\fi}
\newunicodechar{ᵦ}{\ifmmode_{\beta}\else\textsubscript{\ensuremath{\beta}}\fi}
\newunicodechar{ₓ}{\ifmmode_{x}\else\textsubscript{\ensuremath{x}}\fi}
\newfontfamily\popdisplay{ArchivoBlack-Regular.ttf}[Path=../../../fonts/]
\newfontfamily\poplogo{SpaceGrotesk-Bold.ttf}[Path=../../../fonts/]
\newfontfamily\popsemi{PlusJakartaSans-SemiBold.ttf}[Path=../../../fonts/]
\newfontfamily\popxbold{PlusJakartaSans-ExtraBold.ttf}[Path=../../../fonts/]
\newfontfamily\popxboldls{PlusJakartaSans-ExtraBold.ttf}[Path=../../../fonts/,LetterSpace=3.0]

\setlist[enumerate]{leftmargin=1.7em,itemsep=5pt,topsep=4pt}
\setlist[itemize]{leftmargin=1.7em,itemsep=4pt,topsep=4pt,label={\color{poporange}\textbullet}}
\setlength{\parindent}{0pt}
\setlength{\parskip}{5pt}
\renewcommand{\arraystretch}{1.25}

% pgfplots : style Pop par défaut
\pgfplotsset{
  every axis/.append style={axis line style={popink,line width=1pt},tick style={popink},
    grid style={popline,line width=0.5pt},label style={font=\small},tick label style={font=\footnotesize}},
  popcurve/.style={poporange,line width=1.8pt,no markers,smooth},
}
% Même style disponible dans un \draw TikZ simple (hors pgfplots)
\tikzset{popcurve/.style={poporange,line width=1.8pt}}
\tikzset{popgrid/.style={popline,very thin}}
\colorlet{popcurve}{poporange} % le nom du style est parfois utilisé comme couleur

% ============================================================
% Pill / squiggle
% ============================================================
\newcommand{\poppill}[3]{%
  \tikz[baseline=-0.55ex]\node[draw=popink,line width=1.2pt,rounded corners=2.6pt,%
    fill=#2,inner xsep=6.5pt,inner ysep=3pt]{%
    \popxboldls\fontsize{7}{7}\selectfont\color{#3}\MakeUppercase{#1}};%
}
\newcommand{\popsquiggle}[1]{%
  \tikz[baseline=-0.4ex]{
    \draw[#1,line width=2.6pt,line cap=round]
      plot [smooth] coordinates {(0,0.09) (0.34,0.01) (0.68,0.21) (1.02,0.01) (1.36,0.10)};
  }%
}
% Mot-clé surligné (marqueur orange sous le texte)
\newcommand{\cle}[1]{{\popxbold\color{popdark}#1}}

% ============================================================
% En-tête / pied
% ============================================================
\pagestyle{fancy}
\fancyhf{}
\renewcommand{\headrulewidth}{0pt}
\renewcommand{\footrulewidth}{0pt}
\lhead{\raisebox{-0.15ex}{\poplogo\fontsize{13}{13}\selectfont\color{popink}OnBuch\color{poporange}+}}
\rhead{\poppill{\DOCMATIERE\ \textbullet\ \DOCNIVEAU\ \textbullet\ Leçon \DOCLECON}{white}{popink}}
\lfoot{\popxbold\fontsize{7.6}{7.6}\selectfont\color{popmuted}\MakeUppercase{Cours signé OnBuch+}\ \textbullet\ {\popsemi\DOCTITRE}}
\rfoot{\tikz[baseline=-0.6ex]\node[rounded corners=8.5pt,fill=popink,minimum height=17pt,minimum width=17pt,inner xsep=5pt,inner sep=0]{\popxbold\fontsize{8}{8}\selectfont\color{popcream}\thepage\,/\,\pageref*{LastPage}};}

% ============================================================
% Titres
% ============================================================
\newcommand{\chaptitre}[2]{%
  \begin{tcolorbox}[enhanced,colback=poporange,colframe=popink,boxrule=2.6pt,arc=11pt,
      drop shadow=popink,left=15pt,right=15pt,top=12pt,bottom=11pt,before skip=3pt,after skip=18pt]
    \poppill{#2}{popink}{popcream}\par\vspace{9pt}
    {\raggedright\hyphenpenalty=10000\exhyphenpenalty=10000\popdisplay\fontsize{22}{24}\selectfont\color{popcream}\MakeUppercase{#1}\par}
  \end{tcolorbox}
}
% Section numérotée : pastille orange avec le numéro + titre Archivo
\newcounter{popsec}
\newcommand{\coursec}[1]{%
  \stepcounter{popsec}\setcounter{popsub}{0}%
  \par\vspace{16pt}\noindent
  \begin{minipage}{\linewidth}
  \tikz[baseline=-0.7ex]\node[draw=popink,line width=1.6pt,fill=poporange,rounded corners=4pt,minimum size=22pt,inner sep=0]{\popdisplay\fontsize{12}{12}\selectfont\color{popcream}\thepopsec};%
  \hspace{8pt}{\popdisplay\fontsize{15}{17}\selectfont\color{popink}\MakeUppercase{#1}}\par
  \vspace{4pt}\noindent{\color{poporange}\rule{36pt}{3.6pt}}
  \end{minipage}\par\nopagebreak\vspace{8pt}
}
\newcounter{popsub}
\newcommand{\courssub}[1]{%
  \stepcounter{popsub}%
  \par\vspace{9pt}\noindent{\popxbold\fontsize{11.5}{13}\selectfont\color{popink}{\color{poporange}\thepopsec.\thepopsub}\hspace{6pt}#1}\par\nopagebreak\vspace{4pt}
}

% ============================================================
% Boîtes pédagogiques — même recette : fond blanc, bordure épaisse colorée,
% ombre dure encre, badge plein. Titre optionnel [#1].
% ============================================================
\tcbset{popbase/.style={breakable,enhanced,colback=white,boxrule=2.3pt,arc=9pt,
  shadow={3pt}{-3pt}{0pt}{popink},left=14pt,right=14pt,top=11pt,bottom=11pt,before skip=11pt,after skip=15pt,
  fontupper=\popsemi\fontsize{10.6}{14.5}\selectfont\color{popink},
  fontlower=\popsemi\fontsize{10.6}{14.5}\selectfont\color{popink}}}
% Coche dessinée (le glyphe ✓ n'existe pas dans les polices du document)
\newcommand{\popcheck}[1]{\tikz[baseline=-0.5ex]\draw[#1,line width=1.6pt,line cap=round,line join=round](-0.13,0.0)--(-0.04,-0.09)--(0.13,0.1);}
\providecommand{\coche}{\popcheck{popgreen}}
\providecommand{\resultat}[1]{\par\smallskip\noindent\fcolorbox{popgreen}{popgreenL}{\parbox{\dimexpr\linewidth-2\fboxsep-2\fboxrule\relax}{\textbf{Résultat.} #1}}\par\smallskip}
\newcommand{\popboxhead}[3]{\poppill{#1}{#2}{white}\ifx\relax#3\relax\else\hspace{6pt}{\popxbold\fontsize{10.6}{12}\selectfont\color{popink}#3}\fi\par\vspace{7pt}}

\newtcolorbox{aretenir}[1][]{popbase,colframe=popgreen,before upper={\popboxhead{À retenir}{popgreen}{#1}}}
% Texte en chinois (document de lecture, dialogue, extrait) : cadre
% d'étude. Un titre optionnel (source, auteur) s'affiche dans l'en-tête.
\newtcolorbox{textebox}[1][]{popbase,colframe=popblue,before upper={\popboxhead{文本 · Texte}{popblue}{#1}},after upper={}}
% Marques ✓ / ✗ (forme correcte / fautive) souvent utilisées par les modèles
\providecommand{\cross}{\textcolor{poppink}{\ensuremath{\times}}}
\providecommand{\xmark}{\cross}\providecommand{\crossmark}{\cross}\providecommand{\cmark}{\ensuremath{\checkmark}}
% Ligne de réponse / justification à compléter (inventée par les modèles)
\providecommand{\justification}[1]{\par\noindent\textit{Justification~:} \dotfill\par}
% \textipa (package tipa non chargé) : la police ne couvre pas l'API -> simple passage
\providecommand{\textipa}[1]{#1}
% Soulignement (ulem non chargé) : \uline, \ul
\providecommand{\uline}[1]{\underline{#1}}\providecommand{\ul}[1]{\underline{#1}}
% \glossaire{\item ...} inventée par les modèles : liste de glossaire
\providecommand{\glossaire}[1]{\begin{itemize}#1\end{itemize}}
% \alert (beamer) inventée par les modèles : texte mis en évidence
\providecommand{\alert}[1]{\textbf{\textcolor{poppink}{#1}}}
% variantes ASCII d'ordinaux écrites en macro
\providecommand{\iere}{\textsuperscript{re}}\providecommand{\ieme}{\textsuperscript{e}}\providecommand{\ere}{\textsuperscript{re}}\providecommand{\eme}{\textsuperscript{e}}\providecommand{\er}{\textsuperscript{er}}\providecommand{\ie}{\textsuperscript{e}}
% Ordinaux français écrits en macro par les modèles : 1\ière, 3\ième
\newcommand{\ière}{\textsuperscript{re}}\newcommand{\ième}{\textsuperscript{e}}
% \exemple (inventée par les modèles) : étiquette « Exemple : »
\providecommand{\exemple}{\textbf{Exemple~:}\ }
% \square utilisable aussi hors mode math (cases à cocher)
\AtBeginDocument{\let\squareMath\square\renewcommand{\square}{\ensuremath{\squareMath}}}
% Mot ou phrase en chinois à l'intérieur d'un paragraphe français
\newcommand{\zh}[1]{\ifmmode\text{#1}\else{#1}\fi}
\let\eng\zh \let\esp\zh \let\all\zh % alias tolérés
% flèches d'intonation inventées par les modèles
\providecommand{\textsearrow}{}\renewcommand{\textsearrow}{\ensuremath{\searrow}}\providecommand{\textnearrow}{}\renewcommand{\textnearrow}{\ensuremath{\nearrow}}
% \fr{..} : mot français (inventé par les modèles)
\providecommand{\fr}[1]{\textit{#1}}
% zhbox : encadré de texte chinois inventé par les modèles
\newenvironment{zhbox}{\begin{quote}}{\end{quote}}
% \courssubsub : titre de niveau 3 inventé par les modèles
\providecommand{\courssubsub}[1]{\par\medskip\noindent\textbf{#1}\par\smallskip}
% \zhbf : texte chinois en gras inventé par les modèles
\providecommand{\zhbf}[1]{\textbf{#1}}
% \vs : « versus » inventé par les modèles
\providecommand{\vs}{\textit{vs}\ }
% \blank : case à compléter inventée par les modèles
\providecommand{\blank}[1][]{\underline{\hspace{1.6em}}}
\newtcolorbox{exemplebox}[1][]{popbase,colframe=poporange,before upper={\popboxhead{Exemple}{poporange}{#1}}}
\newtcolorbox{definition}[1][]{popbase,colframe=popblue,before upper={\popboxhead{Définition}{popblue}{#1}}}
\newtcolorbox{propriete}[1][]{popbase,colframe=poppurple,before upper={\popboxhead{Propriété}{poppurple}{#1}}}
\newtcolorbox{methode}[1][]{popbase,colframe=popgold,before upper={\popboxhead{Méthode}{popgold}{#1}}}
\newtcolorbox{attention}[1][]{popbase,colframe=poppink,before upper={\popboxhead{Attention}{poppink}{#1}}}
\newtcolorbox{experience}[1][]{popbase,colframe=popblue,colback=popblueL,before upper={\popboxhead{Expérience}{popblue}{#1}}}
% « Le savais-tu ? » : carte violette pleine (contexte, culture, Cameroun)
\newtcolorbox{savaistu}[1][]{popbase,colback=poppurple,colframe=popink,
  fontupper=\popsemi\fontsize{10.4}{14.2}\selectfont\color{white},
  before upper={\poppill{Le savais-tu\,?}{white}{poppurple}\ifx\relax#1\relax\else\hspace{6pt}{\popxbold\color{white}#1}\fi\par\vspace{7pt}}}
% Exercice résolu : énoncé en haut, \tcblower puis la solution sur fond crème
\newcounter{popexo}\newcounter{popexr}
\newtcolorbox{exoresolu}[1][]{popbase,colframe=poporange,colbacklower=popcream,
  segmentation style={popink,line width=1.2pt,dashed},
  before upper={\stepcounter{popexr}\popboxhead{Exercice résolu \thepopexr}{poporange}{#1}},
  before lower={\poppill{Solution}{popink}{popcream}\par\vspace{6pt}}}
% Exercice (non résolu en place) — la correction est dans la partie Corrigés
\newtcolorbox{exercice}[1][]{popbase,colframe=popink,
  before upper={\stepcounter{popexo}\popboxhead{Exercice \thepopexo}{popink}{#1}}}
% Corrigé
\newtcolorbox{corrige}[1][]{popbase,colframe=popgreen,colback=popgreenL,
  before upper={\popboxhead{Corrigé}{popgreen}{#1}}}
% Objectifs / prérequis / bilan
\newtcolorbox{objectifs}{popbase,colframe=popink,colback=poporangeL,
  before upper={\popboxhead{Objectifs de la leçon}{popink}{}}}
\newtcolorbox{prerequis}{popbase,colframe=popink,colback=popblueL,
  before upper={\popboxhead{Prérequis}{popblue}{}}}
\newtcolorbox{bilan}{popbase,colframe=popink,colback=white,boxrule=2.8pt,
  before upper={\popboxhead{Fiche bilan}{poporange}{L'essentiel à connaître pour l'examen}}}
% Figure encadrée avec légende
\newtcolorbox{popfigure}[1][]{enhanced,breakable=false,colback=white,colframe=popline,boxrule=1.4pt,arc=8pt,
  left=10pt,right=10pt,top=10pt,bottom=8pt,before skip=10pt,after skip=12pt,halign=center,
  lower separated=false,halign lower=center,
  fontlower=\popsemi\fontsize{9}{11.5}\selectfont\color{popmuted},
  #1}
\newcommand{\legende}[1]{\par\vspace{6pt}{\popsemi\fontsize{9}{11.5}\selectfont\color{popmuted}#1}}

% ============================================================
% Signature OnBuch+
% ============================================================
\newcommand{\poplogoinline}[1]{{\poplogo\fontsize{#1}{#1}\selectfont\color{popink}OnBuch\color{poporange}+}}
\newcommand{\signatureonbuch}{%
  \begin{tcolorbox}[enhanced,colback=popink,colframe=popink,boxrule=2.2pt,arc=10pt,
      drop shadow=poporange,left=14pt,right=14pt,top=10pt,bottom=10pt,before skip=0pt,after skip=0pt]
    \tikz[baseline=-0.7ex]\node[circle,fill=popgreen,draw=popcream,line width=1.3pt,minimum size=22pt,inner sep=0]{\popcheck{white}};%
    \hspace{9pt}\begin{minipage}[c]{0.62\linewidth}
      {\popxbold\fontsize{12}{13}\selectfont\color{popcream}Signé\ \textbullet\ {\poplogo OnBuch\color{poporange}+}}\par\vspace{2pt}
      {\raggedright\popsemi\fontsize{8}{10.5}\selectfont\color{popline}\MakeUppercase{Équipe pédagogique OnBuch+}\\\MakeUppercase{Conforme au programme officiel MINESEC}\par}
    \end{minipage}\hfill
    {\poplogo\fontsize{22}{22}\selectfont\color{popcream}OnBuch\color{poporange}+}
  \end{tcolorbox}%
}

% ============================================================
% Page de garde
% #1 titre · #2 matière · #3 niveau · #4 module · #5 n° leçon · #6 durée ·
% #7 description / objectifs (texte ou itemize)
% ============================================================
\newcommand{\pagegarde}[7]{%
  \thispagestyle{empty}
  \newgeometry{a4paper,margin=1.7cm,top=1.6cm,bottom=1.4cm}%
  \begin{tikzpicture}[remember picture,overlay]
    \fill[poporange!18] ([xshift=-3.2cm,yshift=-3.6cm]current page.north east) circle (4.6cm);
    \fill[poppurple!14] ([xshift=2.4cm,yshift=7.6cm]current page.south west) circle (3.8cm);
  \end{tikzpicture}%
  {\poplogo\fontsize{30}{30}\selectfont\color{popink}OnBuch\color{poporange}+}\hfill
  \raisebox{0.6ex}{\poppill{Cours signé \textbullet\ Premium}{popink}{popcream}}\par
  \vspace{1pt}\popsquiggle{poppurple}\par
  \vspace{8pt}
  \begin{tcolorbox}[enhanced,colback=poporange,colframe=popink,boxrule=2.8pt,arc=13pt,
      drop shadow=popink,left=18pt,right=18pt,top=12pt,bottom=13pt,after skip=10pt]
    \poppill{#2 \textbullet\ #3}{popink}{popcream}\hspace{5pt}\poppill{#4}{white}{popink}\par\vspace{8pt}
    {\popdisplay\fontsize{14}{14}\selectfont\color{popink}LEÇON #5}\par\vspace{4pt}
    {\raggedright\hyphenpenalty=10000\exhyphenpenalty=10000\popdisplay\fontsize{25}{27}\selectfont\color{popcream}\MakeUppercase{#1}\par}
    \vspace{8pt}
    {\popxbold\fontsize{9.5}{9.5}\selectfont\color{popink}\MakeUppercase{Durée conseillée : #6}}
  \end{tcolorbox}
  \begin{tcolorbox}[enhanced,colback=white,colframe=popink,boxrule=2.2pt,arc=10pt,
      drop shadow=popink,left=16pt,right=16pt,top=10pt,bottom=10pt,after skip=8pt]
    \poppill{Dans cette leçon}{poppurple}{white}\par\vspace{5pt}
    \popsemi\fontsize{9.3}{12.6}\selectfont\color{popink}#7
  \end{tcolorbox}
  \noindent\begin{minipage}[t]{0.487\linewidth}
    \begin{tcolorbox}[enhanced,colback=popgreen,colframe=popink,boxrule=2.2pt,arc=10pt,
        drop shadow=popink,left=11pt,right=11pt,top=10pt,bottom=10pt,height=3.1cm,valign=center]
      \tcbox[colback=white,colframe=popink,boxrule=1.3pt,arc=5pt,boxsep=0pt,nobeforeafter,
        left=3pt,right=3pt,top=3pt,bottom=3pt,valign=center]{\qrcode[height=1.9cm]{https://play.google.com/store/apps/details?id=com.luvvix.onbuch&hl=fr}}%
      \hspace{8pt}\begin{minipage}[c]{0.5\linewidth}
        {\popdisplay\fontsize{11}{12.5}\selectfont\color{white}\MakeUppercase{Télécharge OnBuch}}\par\vspace{4pt}
        {\raggedright\popsemi\fontsize{8.4}{11}\selectfont\color{white}Gratuit \textbullet\ Google Play\\Tuteur IA, annales, concours, résultats\par}
      \end{minipage}
    \end{tcolorbox}
  \end{minipage}\hfill
  \begin{minipage}[t]{0.487\linewidth}
    \begin{tcolorbox}[enhanced,colback=poppurple,colframe=popink,boxrule=2.2pt,arc=10pt,
        drop shadow=popink,left=11pt,right=11pt,top=10pt,bottom=10pt,height=3.1cm,valign=center]
      \tikz[baseline=-0.7ex]\node[circle,fill=white,draw=popink,line width=1.3pt,minimum size=1.6cm,inner sep=0]{\popdisplay\fontsize{24}{24}\selectfont\color{poppurple}+};%
      \hspace{8pt}\begin{minipage}[c]{0.6\linewidth}
        {\popdisplay\fontsize{11}{12.5}\selectfont\color{white}\MakeUppercase{Passe à OnBuch+}}\par\vspace{4pt}
        {\raggedright\popsemi\fontsize{8.4}{11}\selectfont\color{white}Tous les cours signés, Tuteur IA illimité, fiches PDF — onglet OnBuch+ de l'app\par}
      \end{minipage}
    \end{tcolorbox}
  \end{minipage}
  \vspace{8pt}
  \popcontenu
  \vfill
  \signatureonbuch
  \restoregeometry
  \newpage
}

% Rangée « Ce cours contient » (page de garde)
\newcommand{\poptile}[3]{% #1 couleur #2 gros texte #3 légende
  \begin{tcolorbox}[enhanced,colback=#1,colframe=popink,boxrule=2pt,arc=9pt,shadow={2.5pt}{-2.5pt}{0pt}{popink},
      left=6pt,right=6pt,top=9pt,bottom=9pt,halign=center,valign=center,height=1.75cm,width=\linewidth,nobeforeafter]
    {\popdisplay\fontsize{12.5}{13}\selectfont\color{white}\MakeUppercase{#2}}\par\vspace{3pt}
    {\centering\popsemi\fontsize{7.8}{9.5}\selectfont\color{white}#3\par}
  \end{tcolorbox}}
\newcommand{\popcontenu}{%
  \par\noindent{\popxboldls\fontsize{7.5}{7.5}\selectfont\color{popmuted}CE COURS SIGNÉ CONTIENT}\par\vspace{7pt}
  \noindent\begin{minipage}[t]{0.235\linewidth}\poptile{popblue}{Cours}{complet, illustré, pas à pas}\end{minipage}\hfill
  \begin{minipage}[t]{0.235\linewidth}\poptile{poporange}{Méthodes}{et exercices résolus}\end{minipage}\hfill
  \begin{minipage}[t]{0.235\linewidth}\poptile{poppink}{Exercices}{type examen + corrigés}\end{minipage}\hfill
  \begin{minipage}[t]{0.235\linewidth}\poptile{popgreen}{Fiche bilan}{l'essentiel à retenir}\end{minipage}\par}

% Sommaire
\newtcolorbox{sommaire}{enhanced,colback=white,colframe=popink,boxrule=2.2pt,arc=10pt,
  drop shadow=popink,left=18pt,right=18pt,top=13pt,bottom=13pt,before skip=6pt,after skip=20pt,
  before upper={\poppill{Sommaire}{poppurple}{white}\par\vspace{9pt}\popsemi\fontsize{10.6}{14.5}\selectfont\color{popink}}}

% Signature de fin de cours — poussée tout en bas de la dernière page
\newcommand{\findecours}{%
  \par\vfill\nopagebreak
  \begin{center}\popsquiggle{poporange}\end{center}\vspace{4pt}
  \signatureonbuch
}
\AtBeginDocument{\def\llbracket{\mathopen{[\mkern-3.5mu[}}\def\rrbracket{\mathclose{]\mkern-3.5mu]}}} % intervalles d'entiers (glyphe ⟦ absent de la police)
% Glyphes absents de Fira Math : redéfinis à partir de glyphes disponibles
\AtBeginDocument{%
  \def\setminus{\mathbin{\backslash}}\let\smallsetminus\setminus
  \def\vdots{\mathord{\vbox{\baselineskip3.2pt\lineskiplimit0pt\kern4pt\hbox{.}\hbox{.}\hbox{.}}}}%
  \def\ddots{\mathinner{\mkern1mu\raise6.5pt\hbox{.}\mkern2mu\raise3.6pt\hbox{.}\mkern2mu\raise0.7pt\hbox{.}\mkern1mu}}%
  \def\triangle{\mathord{\scalebox{0.9}{$\symup{\Delta}$}}}%
  \def\nabla{\mathord{\raisebox{\depth}{\scalebox{1}[-1]{$\symup{\Delta}$}}}}%
  \def\checkmark{\popcheck{popdark}}%
  \def\star{\mathbin{\ast}}\def\bigstar{\mathord{\ast}}%
  \def\diamond{\mathbin{\scalebox{0.6}{$\Diamond$}}}%
  \def\bigcup{\mathop{\mathchoice{\raisebox{-0.3em}{\scalebox{1.7}{$\cup$}}}{\scalebox{1.25}{$\cup$}}{\cup}{\cup}}\displaylimits}%
  \def\bigcap{\mathop{\mathchoice{\raisebox{-0.3em}{\scalebox{1.7}{$\cap$}}}{\scalebox{1.25}{$\cap$}}{\cap}{\cap}}\displaylimits}%
  \def\lesseqqgtr{\mathrel{\lessgtr}}\def\bigoplus{\mathop{\oplus}\displaylimits}%
}
\providecommand{\ssi}{si et seulement si} % abréviation fréquente
\AtBeginDocument{\providecommand{\wideparen}{\overparen}} % arc de cercle

%% FIGFIX-BEGIN v3 — correctifs globaux des figures (docs/onbuch-generation/tools/figfix)
\usepackage{adjustbox}\usepackage{etoolbox}
% toute figure TikZ/pgfplots plus large que la ligne est réduite à la largeur disponible
\BeforeBeginEnvironment{tikzpicture}{\begin{adjustbox}{max width=\linewidth}}
\AfterEndEnvironment{tikzpicture}{\end{adjustbox}}
% légendes pgfplots lisibles par-dessus les courbes
\pgfplotsset{every axis/.append style={legend style={fill=white,fill opacity=0.92,text opacity=1,draw=black!15,inner sep=3pt}}}
% étiquettes d'arêtes (syntaxe edge["..."]) avec fond blanc
\tikzset{every edge quotes/.append style={fill=white,inner sep=1.2pt}}
% bulles de cartes mentales : texte à la ligne dans la bulle, bulle un peu plus grande
\tikzset{every concept/.append style={text width=2.0cm,align=center,minimum size=2.3cm,inner sep=2pt}}
%% FIGFIX-END
\begin{document}

\pagegarde{Leçon 43 — Révisions générales}{Chinois}{1ère A}{Module 5 : Mass médias et télécommunication}{ZHA43}{2 h}{Leçon de révisions générales clôturant l'année de 1ère A : synthèse du vocabulaire, des structures grammaticales et des caractères étudiés, suivie d'exercices progressifs de type examen couvrant compréhension, grammaire, traduction et expression écrite.
\vspace{4pt}
\begin{multicols}{2}\small
\begin{itemize}
\item À la fin de cette leçon, je sais mobiliser le vocabulaire essentiel de l'année (salutations, famille, école, temps, loisirs, médias)
\item Je sais utiliser correctement les structures grammaticales clés : 是, 有, 在, 了, 的, 吗, 呢, 不/没, les mots interrogatifs et les compléments de mesure
\item Je sais reconnaître et écrire les caractères fréquents de l'année en respectant radicaux et ordre des traits
\item Je sais comprendre un texte court de type examen et répondre à des questions de compréhension
\item Je sais traduire des phrases du français vers le chinois et du chinois vers le français
\item Je sais rédiger un court texte cohérent (présentation, journée, message) avec des connecteurs simples
\end{itemize}\end{multicols}}

\begin{sommaire}
\begin{multicols}{2}
\begin{enumerate}
\item Bilan du vocabulaire de l'année
\item Synthèse des structures grammaticales
\item Caractères : radicaux, traits, pièges
\item Entraînement à la compréhension écrite
\item Traduction et expression écrite guidées
\item Simulation d'examen et stratégies
\item Activité d'intégration
\item Méthodes et exercices résolus
\item Exercices
\item Corrigés des exercices
\item Fiche bilan
\end{enumerate}
\end{multicols}
\end{sommaire}

\begin{objectifs}
\begin{itemize}
\item À la fin de cette leçon, je sais mobiliser le vocabulaire essentiel de l'année (salutations, famille, école, temps, loisirs, médias)
\item Je sais utiliser correctement les structures grammaticales clés : 是, 有, 在, 了, 的, 吗, 呢, 不/没, les mots interrogatifs et les compléments de mesure
\item Je sais reconnaître et écrire les caractères fréquents de l'année en respectant radicaux et ordre des traits
\item Je sais comprendre un texte court de type examen et répondre à des questions de compréhension
\item Je sais traduire des phrases du français vers le chinois et du chinois vers le français
\item Je sais rédiger un court texte cohérent (présentation, journée, message) avec des connecteurs simples
\item Je sais repérer et corriger les erreurs typiques des francophones (ordre des mots, omission de 是, confusion 不/没)
\item Je sais gérer mon temps et ma stratégie face à une épreuve de type examen
\end{itemize}
\end{objectifs}

\begin{prerequis}
\begin{itemize}
\item Vocabulaire et structures des leçons 1 à 42 de 1ère A
\item Pinyin, tons et règles de prononciation de base
\item Radicaux et ordre des traits vus depuis la Seconde
\item Structures fondamentales : 是/有/在, phrase verbale SVO, particule 的
\end{itemize}
\end{prerequis}

\newpage

\coursec{Bilan du vocabulaire de l'année}

C'est la fin de l'année à Yaoundé : les élèves de 1ère A du lycée bilingue préparent leur examen de chinois. Amina feuillette ses cahiers : salutations, famille, école, hobbies, mass médias... Comment tout réviser efficacement en deux heures ? Comme un athlète avant la compétition, il faut un entraînement complet : vocabulaire, grammaire, caractères, traduction et expression. Cette leçon est votre séance d'entraînement général avant le grand jour.

\courssub{Carte mentale des champs lexicaux de l'année}

Pour structurer votre révision, rien ne vaut une vision d'ensemble. Le vocabulaire de l'année s'organise autour de cinq grands pôles thématiques. Observez la carte mentale ci-dessous : elle relie le nœud central « 一年词汇 » (le vocabulaire de l'année) aux cinq domaines majeurs, chacun illustré par trois mots-clés à connaître parfaitement (caractère, pinyin, sens).

\begin{popfigure}
\begin{tikzpicture}[
  mindmap,
  grow cyclic,
  every node/.style={concept, execute at begin node=\hskip0pt},
  concept color=popblue,
  level 1/.append style={level distance=4.5cm, sibling angle=72},
  level 2/.append style={level distance=3cm, sibling angle=40},
  identity/.style={concept color=poporange},
  famille/.style={concept color=popgreen},
  ecole/.style={concept color=poppurple},
  temps/.style={concept color=poppink},
  medias/.style={concept color=popgold}
]
\node[concept, align=center]{一年词汇\\Vocabulaire de l'année}
  child[identity] { node[concept] {Identité}
    child { node[concept, align=center]{你\\nǐ\\toi} }
    child { node[concept, align=center]{好\\hǎo\\bien/bon} }
    child { node[concept, align=center]{名字\\míngzi\\nom} }
  }
  child[famille] { node[concept] {Famille}
    child { node[concept, align=center]{家\\jiā\\famille/foyer} }
    child { node[concept, align=center]{爸爸\\bàba\\papa} }
    child { node[concept, align=center]{妈妈\\māma\\maman} }
  }
  child[ecole] { node[concept] {École}
    child { node[concept, align=center]{学校\\xuéxiào\\école} }
    child { node[concept, align=center]{老师\\lǎoshī\\professeur} }
    child { node[concept, align=center]{学生\\xuésheng\\élève} }
  }
  child[temps] { node[concept] {Temps}
    child { node[concept, align=center]{年\\nián\\année} }
    child { node[concept, align=center]{月\\yuè\\mois} }
    child { node[concept, align=center]{星期\\xīngqī\\semaine} }
  }
  child[medias] { node[concept] {Médias}
    child { node[concept, align=center]{电视\\diànshì\\télévision} }
    child { node[concept, align=center]{手机\\shǒujī\\téléphone} }
    child { node[concept, align=center]{新闻\\xīnwén\\nouvelles} }
  };
\end{tikzpicture}
\legende{Carte mentale des 5 champs lexicaux majeurs de l'année (HSK 1-2). Chaque branche contient 3 mots « socles » à maîtriser en priorité : caractères, pinyin, sens et classificateur associé.}
\end{popfigure}

\begin{methode}[Comment réviser cette carte mentale ?]
\begin{enumerate}
\item \textbf{Parcours thématique :} cachez les branches, récitez les 3 mots de chaque domaine en donnant le pinyin (tons exacts) et le sens.
\item \textbf{Parcours transversal :} prenez un mot (ex. \zh{看} kàn), dites à quel domaine il appartient (Médias / Loisirs) et construisez une phrase : \zh{我看电视。} (Wǒ kàn diànshì.) « Je regarde la télévision. »
\item \textbf{Test croisé :} à deux, l'un donne le français, l'autre écrit le caractère + pinyin sur ardoise. Inversez les rôles.
\end{enumerate}
\end{methode}

\courssub{Mots-outils incontournables : pronoms, démonstratifs, interrogatifs}

Ces petits mots sont le « ciment » de la phrase. Sans eux, impossible de construire un énoncé correct. Nous les passons en revue par catégorie.

\begin{definition}[Pronoms personnels et possessifs]
Les pronoms personnels chinois ne distinguent pas le genre à l'oral (\zh{tā} se prononce pareil pour « il » et « elle »), mais seulement à l'écrit (\zh{他} / \zh{她} / \zh{它}). Le possessif se forme simplement en ajoutant \zh{的} (de) après le pronom.
\begin{center}
\begin{tabularx}{\linewidth}{|X|X|X|X|}
\hline
\rowcolor{popblueL} \textbf{Personne} & \textbf{Sujet / Objet} & \textbf{Pinyin} & \textbf{Possessif (avec \zh{的})} \\
\hline
1ère sg. & 我 & wǒ & 我的 (wǒ de) — mon / ma / mes \\
2ème sg. & 你 / 您 & nǐ / nín & 你的 / 您的 (nǐ de / nín de) — ton, tes / votre \\
3ème sg. & 他 / 她 / 它 & tā & 他的 / 她的 / 它的 (tā de) — son / sa / ses \\
1ère pl. & 我们 & wǒmen & 我们的 (wǒmen de) — notre / nos \\
2ème pl. & 你们 & nǐmen & 你们的 (nǐmen de) — votre / vos \\
3ème pl. & 他们 / 她们 / 它们 & tāmen & 他们的 / 她们的 / 它们的 (tāmen de) — leur / leurs \\
\hline
\end{tabularx}
\end{center}
\emph{Note :} \zh{您} (nín) est la forme polie pour « vous » (singulier), utilisée envers un aîné, un professeur, un inconnu. Au Cameroun, on l'utilise naturellement avec les enseignants.
\end{definition}

\begin{definition}[Démonstratifs]
Ils s'emploient obligatoirement avec un classificateur (souvent \zh{个} gè) devant un nom.
\begin{center}
\begin{tabularx}{\linewidth}{|X|X|X|X|}
\hline
\rowcolor{popblueL} \textbf{Proche} & \textbf{Pinyin} & \textbf{Lointain} & \textbf{Pinyin} \\
\hline
这 / 这个 & zhè / zhè ge & 那 / 那个 & nà / nà ge \\
这些 & zhèxiē & 那些 & nàxiē \\
这里 / 这儿 & zhèlǐ / zhèr & 那里 / 那儿 & nàlǐ / nàr \\
\hline
\end{tabularx}
\end{center}
\zh{这本书} (zhè běn shū) — ce livre-ci. \quad \zh{那个学生} (nà ge xuésheng) — cet élève-là.
\end{definition}

\begin{propriete}[Les mots interrogatifs ne changent pas l'ordre de la phrase]
En chinois, les mots interrogatifs ne déclenchent \textbf{ni inversion du sujet ni auxiliaire} (pas de « est-ce que »). Le mot interrogatif occupe simplement la place du mot qu'il remplace : \textbf{Sujet + Verbe + Interrogatif}. Pour une question oui/non, on ajoute la particule finale \zh{吗} (ma) : \zh{你是学生吗？} (Nǐ shì xuésheng ma ?) — Es-tu élève ?
\end{propriete}

\begin{center}
\begin{tabularx}{\linewidth}{|X|X|X|X|}
\hline
\rowcolor{popblueL} \textbf{Caractères} & \textbf{Pinyin} & \textbf{Traduction} & \textbf{Exemple (Caractères + Pinyin + Traduction)} \\
\hline
什么 & shénme & quoi / quel(le) & 你叫什么名字？(Nǐ jiào shénme míngzi ?) — Comment t'appelles-tu ? \\
谁 & shéi & qui & 谁是老师？(Shéi shì lǎoshī ?) — Qui est le professeur ? \\
哪儿 / 哪里 & nǎr / nǎlǐ & où & 你住在哪儿？(Nǐ zhù zài nǎr ?) — Où habites-tu ? \\
怎么 & zěnme & comment & 怎么去学校？(Zěnme qù xuéxiào ?) — Comment aller à l'école ? \\
几 & jǐ & combien (petit nombre) & 你有几个弟弟？(Nǐ yǒu jǐ ge dìdi ?) — Combien as-tu de petits frères ? \\
多少 & duōshao & combien (grand nombre) / quel prix & 你们学校有多少学生？(Nǐmen xuéxiào yǒu duōshao xuésheng ?) — Combien d'élèves y a-t-il dans votre école ? \\
哪 & nǎ & lequel / laquelle (+ classif.) & 哪本书是你的？(Nǎ běn shū shì nǐ de ?) — Quel livre est le tien ? \\
几岁 & jǐ suì & quel âge (enfant) & 你几岁？(Nǐ jǐ suì ?) — Quel âge as-tu ? \\
\hline
\end{tabularx}
\end{center}

\begin{attention}[Piège majeur : \cle{几} vs \cle{多少}]
\begin{itemize}
\item \textbf{几 (jǐ)} : s'utilise pour des quantités \textbf{petites (en général inférieures à 10)}. Il est \textbf{obligatoirement suivi d'un classificateur} : \zh{几个人} (jǐ ge rén) — combien de personnes.
\item \textbf{多少 (duōshao)} : s'utilise pour des quantités \textbf{grandes ou indéterminées} (prix, population). Le classificateur est \textbf{facultatif} : \zh{多少钱？} (Duōshao qián ?) — Combien ça coûte ? \zh{多少(个)学生？} (Duōshao (ge) xuésheng ?)
\end{itemize}
\emph{Erreur fréquente :} demander \zh{你几岁？} à un adulte. \zh{几岁} se réserve aux enfants ; pour un adulte, on dit \zh{您多大？} (Nín duō dà ?) ou, plus formel, \zh{您多大年纪？} (Nín duō dà niánjì ?).
\end{attention}

\courssub{Classificateurs fréquents : 个, 本, 张, 位, 家}

En chinois, \textbf{tout nom compté ou désigné} nécessite un classificateur entre le nombre (ou le démonstratif) et le nom. Structure : \textbf{Nombre / Démonstratif + Classificateur + Nom}.

\begin{center}
\begin{tabularx}{\linewidth}{|X|X|X|X|}
\hline
\rowcolor{popblueL} \textbf{Classif.} & \textbf{Pinyin} & \textbf{Usage principal / Noms typiques} & \textbf{Exemples} \\
\hline
个 & gè & Classificateur universel (personnes, objets divers, notions abstraites) & 一个人 (yí ge rén), 三个苹果 (sān ge píngguǒ), 这个问题 (zhè ge wèntí) \\
本 & běn & Livres, cahiers, revues, dictionnaires & 一本书 (yì běn shū), 两本词典 (liǎng běn cídiǎn) \\
张 & zhāng & Objets plats : papier, table, lit, billet, carte, photo & 一张纸 (yì zhāng zhǐ), 三张桌子 (sān zhāng zhuōzi), 这张照片 (zhè zhāng zhàopiàn) \\
位 & wèi & \textbf{Politesse} : personnes (professeurs, invités, clients) & 一位老师 (yí wèi lǎoshī), 两位同学 (liǎng wèi tóngxué) \\
家 & jiā & Établissements : école, hôpital, restaurant, entreprise ; famille & 一家学校 (yì jiā xuéxiào), 这家医院 (zhè jiā yīyuàn), 我家 (wǒ jiā) \\
\hline
\end{tabularx}
\end{center}

\begin{savaistu}[Le classificateur \cle{口} (kǒu) — « les bouches à nourrir »]
Le caractère \zh{口} (kǒu) signifie « bouche ». Utilisé comme classificateur, il compte les \textbf{membres d'une famille} ou d'un foyer : \zh{五口人} (wǔ kǒu rén) — « cinq bouches », c'est-à-dire une famille de cinq personnes. Image très vivante : on compte ceux qui mangent sous le même toit ! On ne l'utilise \textbf{pas} hors du contexte familial : \zh{*五口学生} est faux, on dit \zh{五个学生} (wǔ ge xuésheng) ou, poliment, \zh{五位学生} (wǔ wèi xuésheng).
\end{savaistu}

\courssub{Verbes d'action fréquents — le « noyau dur » HSK 1-2}

Ces douze verbes couvrent l'essentiel des besoins communicationnels de base. Maîtrisez leur sens, leur pinyin (les tons !) et leur construction.

\begin{center}
\begin{tabularx}{\linewidth}{|X|X|X|X|}
\hline
\rowcolor{popblueL} \textbf{Verbe} & \textbf{Pinyin} & \textbf{Sens principal} & \textbf{Construction type + Exemple} \\
\hline
是 & shì & être (identité, définition) & S + 是 + Nom : 我是学生。(Wǒ shì xuésheng.) \\
有 & yǒu & avoir / il y a (existence) & S + 有 + O : 我有书。(Wǒ yǒu shū.) / 这儿有医院。(Zhèr yǒu yīyuàn.) \\
在 & zài & être situé, se trouver & S + 在 + Lieu : 我在雅温得。(Wǒ zài Yǎwēndé.) \\
去 & qù & aller (vers) & S + 去 + Lieu : 我去学校。(Wǒ qù xuéxiào.) \\
来 & lái & venir (vers le locuteur) & S + 来 + Lieu : 他来喀麦隆。(Tā lái Kāmàilóng.) \\
做 & zuò & faire / fabriquer & S + 做 + O : 做作业 (zuò zuòyè), 做饭 (zuò fàn) \\
学 & xué & apprendre / étudier & S + 学 + O : 学中文 (xué Zhōngwén) \\
说 & shuō & parler / dire & S + 说 + Langue / Paroles : 说法语 (shuō Fǎyǔ), 说汉语 (shuō Hànyǔ) \\
看 & kàn & regarder / lire / voir & S + 看 + O : 看电视 (kàn diànshì), 看书 (kàn shū) \\
听 & tīng & écouter & S + 听 + O : 听音乐 (tīng yīnyuè), 听课 (tīng kè) \\
写 & xiě & écrire & S + 写 + O : 写字 (xiě zì), 写作业 (xiě zuòyè) \\
买 & mǎi & acheter & S + 买 + O : 买手机 (mǎi shǒujī), 买票 (mǎi piào) \\
\hline
\end{tabularx}
\end{center}

\begin{exemplebox}[Verbes en contexte — phrases types]
\begin{itemize}
\item \zh{我每天七点起床，八点上课。} (Wǒ měitiān qī diǎn qǐchuáng, bā diǎn shàngkè.) — Je me lève à 7 h, j'ai cours à 8 h. (\zh{起床} qǐchuáng « se lever », \zh{上课} shàngkè « avoir cours ».)
\item \zh{晚上我喜欢看电视，也喜欢上网看新闻。} (Wǎnshang wǒ xǐhuan kàn diànshì, yě xǐhuan shàngwǎng kàn xīnwén.) — Le soir, j'aime regarder la télévision et aussi aller sur Internet lire les nouvelles.
\item \zh{我想买一本新书。} (Wǒ xiǎng mǎi yì běn xīn shū.) — Je voudrais acheter un nouveau livre. (\zh{想} xiǎng = vouloir, avoir l'intention de.)
\item \zh{请听老师的话。} (Qǐng tīng lǎoshī de huà.) — Écoutez ce que dit le professeur. (\zh{请} qǐng = s'il vous plaît ; introduit l'impératif poli.)
\end{itemize}
\end{exemplebox}

\courssub{Antonymes et paires utiles — enrichir le lexique}

Travailler par paires d'antonymes double l'efficacité de la mémorisation. Voici les dix paires « socles » du programme.

\begin{center}
\begin{tabularx}{\linewidth}{|X|X|X|X|X|X|}
\hline
\rowcolor{popblueL} \textbf{Chinois} & \textbf{Pinyin} & \textbf{Français} & \textbf{Antonyme} & \textbf{Pinyin} & \textbf{Français} \\
\hline
大 & dà & grand & 小 & xiǎo & petit \\
多 & duō & beaucoup & 少 & shǎo & peu \\
好 & hǎo & bon / bien & 坏 & huài & mauvais \\
快 & kuài & rapide / vite & 慢 & màn & lent \\
新 & xīn & nouveau / neuf & 旧 & jiù & vieux / ancien \\
高 & gāo & haut / grand (taille) & 低 & dī & bas \\
长 & cháng & long & 短 & duǎn & court \\
贵 & guì & cher & 便宜 & piányi & bon marché \\
热 & rè & chaud & 冷 & lěng & froid \\
远 & yuǎn & loin & 近 & jìn & près \\
\hline
\end{tabularx}
\end{center}

\begin{aretenir}[Astuce de révision : « la phrase miroir »]
Pour chaque paire, écrivez \textbf{deux phrases miroir} : une avec l'adjectif, une avec son antonyme. Exemple : \zh{这本书很贵。} (Zhè běn shū hěn guì.) « Ce livre est cher. » / \zh{那本书很便宜。} (Nà běn shū hěn piányi.) « Ce livre-là est bon marché. » Cela ancre le vocabulaire dans la syntaxe. Rappel : un adjectif seul ne peut pas être prédicat ; on ajoute \zh{很} hěn (sans valeur intensive réelle) ou \zh{不} bù pour la négation.
\end{aretenir}

\courssub{Texte support de révision : « Mon quotidien à Yaoundé »}

Le texte ci-dessous synthétise les champs lexicaux de l'année : identité, famille, temps, école, médias. Lisez-le à voix haute \textbf{trois fois} en respectant scrupuleusement les tons (marqués sur le pinyin). Ensuite, répondez aux questions de compréhension.

\begin{textebox}[Texte de révision global — 我叫李明]
\zh{我叫李明，是雅温得一所中学的学生。} \\
\emph{Wǒ jiào Lǐ Míng, shì Yǎwēndé yì suǒ zhōngxué de xuésheng.} \\
\zh{我家有五口人：爸爸、妈妈、哥哥、妹妹和我。} \\
\emph{Wǒ jiā yǒu wǔ kǒu rén : bàba, māma, gēge, mèimei hé wǒ.} \\
\zh{我每天六点半起床，七点吃早饭，七点半去学校。} \\
\emph{Wǒ měitiān liù diǎn bàn qǐchuáng, qī diǎn chī zǎofàn, qī diǎn bàn qù xuéxiào.} \\
\zh{上午有四节课，中午在食堂吃饭。} \\
\emph{Shàngwǔ yǒu sì jié kè, zhōngwǔ zài shítáng chī fàn.} \\
\zh{下午三点放学，我和同学一起踢足球。} \\
\emph{Xiàwǔ sān diǎn fàngxué, wǒ hé tóngxué yìqǐ tī zúqiú.} \\
\zh{晚上七点做作业，九点看电视或者上网看新闻。} \\
\emph{Wǎnshang qī diǎn zuò zuòyè, jiǔ diǎn kàn diànshì huòzhě shàngwǎng kàn xīnwén.} \\
\zh{周末我喜欢去公园，也帮妈妈买菜做饭。} \\
\emph{Zhōumò wǒ xǐhuan qù gōngyuán, yě bāng māma mǎi cài zuò fàn.}
\end{textebox}

\begin{center}
\begin{tabularx}{\linewidth}{|X|X|X|}
\hline
\rowcolor{popblueL} \textbf{Mots nouveaux / Rappels} & \textbf{Pinyin} & \textbf{Traduction} \\
\hline
中学 & zhōngxué & collège / lycée \\
食堂 & shítáng & cantine / réfectoire \\
放学 & fàngxué & finir les cours / sortir de l'école \\
踢足球 & tī zúqiú & jouer au football \\
作业 & zuòyè & devoirs \\
或者 & huòzhě & ou bien (dans une affirmation) \\
周末 & zhōumò & week-end \\
公园 & gōngyuán & parc \\
买菜 & mǎi cài & faire le marché / acheter des légumes \\
\hline
\end{tabularx}
\end{center}

\begin{exoresolu}[Compréhension du texte — questions type examen]
\textbf{Énoncé :} lisez le texte ci-dessus et répondez en français aux questions suivantes.
\begin{enumerate}
\item Qui parle dans ce texte ? Quel est son statut ?
\item Combien de personnes composent sa famille ? Citez-les.
\item À quelle heure se lève-t-il ? À quelle heure va-t-il à l'école ?
\item Que fait-il l'après-midi après les cours ?
\item Quelles sont ses deux activités du soir après les devoirs ?
\item Que fait-il le week-end ? (Deux activités.)
\end{enumerate}

\tcblower
\textbf{Solution détaillée :}
\begin{enumerate}
\item C'est \textbf{Li Ming} (\zh{李明}) qui parle. Il est \textbf{élève} (\zh{学生}) dans un \textbf{collège/lycée} (\zh{中学}) à \textbf{Yaoundé} (\zh{雅温得}).
\item Sa famille compte \textbf{cinq personnes} (\zh{五口人}) : son \textbf{père} (\zh{爸爸}), sa \textbf{mère} (\zh{妈妈}), son \textbf{grand frère} (\zh{哥哥}), sa \textbf{petite sœur} (\zh{妹妹}) et \textbf{lui-même} (\zh{我}).
\item Il se lève à \textbf{6 h 30} (\zh{六点半}) et va à l'école à \textbf{7 h 30} (\zh{七点半}). (Il petit-déjeune à 7 h.)
\item L'après-midi, les cours finissent à \textbf{15 h} (\zh{三点}). Il \textbf{joue au football} (\zh{踢足球}) avec ses \textbf{camarades de classe} (\zh{同学}).
\item Après les devoirs (19 h), il \textbf{regarde la télévision} (\zh{看电视}) \textbf{ou bien} (\zh{或者}) \textbf{va sur Internet lire les nouvelles} (\zh{上网看新闻}).
\item Le week-end, il \textbf{aime aller au parc} (\zh{去公园}) et \textbf{aide sa mère à faire le marché et à cuisiner} (\zh{帮妈妈买菜做饭}).
\end{enumerate}
\end{exoresolu}

\courssub{Entraînement guidé : traduction et réemploi}

\begin{exercice}[Mise en pratique immédiate — Thème (français → chinois)]
Traduisez les phrases suivantes en chinois (caractères + pinyin). Utilisez le vocabulaire et les structures de la leçon.
\begin{enumerate}
\item Ma famille compte quatre personnes : papa, maman, ma grande sœur et moi.
\item Je vais à l'école à 7 h 20 tous les jours.
\item Le soir, mon grand frère regarde la télévision, moi je lis un livre.
\item Ce week-end, je voudrais aller au parc avec mes camarades de classe.
\item Combien d'élèves y a-t-il dans votre école ?
\end{enumerate}
\end{exercice}

\begin{corrige}[Exercice 1 — Thème]
\begin{enumerate}
\item \zh{我家有四口人：爸爸、妈妈、姐姐和我。} (Wǒ jiā yǒu sì kǒu rén : bàba, māma, jiějie hé wǒ.)
\item \zh{我每天七点二十去学校。} (Wǒ měitiān qī diǎn èrshí qù xuéxiào.) — \emph{Note :} 20 = \zh{二十} èrshí ; on peut aussi dire \zh{七点二十分} (qī diǎn èrshí fēn).
\item \zh{晚上，哥哥看电视，我看书。} (Wǎnshang, gēge kàn diànshì, wǒ kàn shū.) — \emph{Structure :} deux propositions juxtaposées, sans mot de liaison.
\item \zh{这个周末，我想和同学一起去公园。} (Zhè ge zhōumò, wǒ xiǎng hé tóngxué yìqǐ qù gōngyuán.)
\item \zh{你们学校有多少学生？} (Nǐmen xuéxiào yǒu duōshao xuésheng ?) — \emph{Rappel :} \zh{多少} pour un grand nombre ; \zh{你们} pour « votre » (pluriel).
\end{enumerate}
\end{corrige}

\begin{exercice}[Version (chinois → français) — précision lexicale]
Traduisez en français en soignant le registre (familier / standard / soutenu).
\begin{enumerate}
\item \zh{请问，那位老师姓什么？}
\item \zh{我没有哥哥，只有一个妹妹。}
\item \zh{这本中文书很贵，那本很便宜。}
\item \zh{明天上午我们不上课，去踢足球吧！}
\item \zh{您贵姓？}
\end{enumerate}
\end{exercice}

\begin{corrige}[Exercice 2 — Version]
\begin{enumerate}
\item « Excusez-moi, quel est le nom de famille de ce professeur ? » — \zh{姓} xìng = avoir pour nom de famille ; \zh{请问} qǐngwèn = excusez-moi de vous demander ; \zh{位} wèi = classificateur de politesse.
\item « Je n'ai pas de grand frère, je n'ai qu'une petite sœur. » — \zh{只有} zhǐyǒu = n'avoir que, seulement.
\item « Ce livre de chinois est cher, celui-là est bon marché. » — \zh{那本} nà běn = celui-là (le classificateur \zh{本} suffit, le nom \zh{书} est sous-entendu) ; \zh{便宜} piányi = bon marché.
\item « Demain matin, nous n'avons pas cours, allons jouer au foot ! » — \zh{吧} ba = particule finale de suggestion ; \zh{不上课} bù shàngkè = ne pas avoir cours.
\item « Quel est votre nom de famille ? » — Formule \textbf{très polie} (\zh{贵} guì = « honorable »), réservée aux aînés et aux inconnus. Ne dites pas \zh{你叫什么名字？} à un supérieur.
\end{enumerate}
\end{corrige}

\begin{attention}[Erreurs fréquentes des francophones — bilan vocabulaire]
\begin{itemize}
\item \textbf{Confondre \zh{二} (èr) et \zh{两} (liǎng) :} \zh{二} sert à compter (un, deux, trois...), dans les numéros de téléphone et les nombres composés (\zh{十二} shí'èr = 12). \zh{两} est \textbf{obligatoire devant un classificateur} : \zh{两个人} (liǎng ge rén), \zh{两本书} (liǎng běn shū), \zh{两点} (liǎng diǎn — 2 h).
\item \textbf{Oublier le classificateur :} \zh{*三书} est faux → \zh{三本书} (sān běn shū). \zh{*一老师} est faux → \zh{一位老师} (yí wèi lǎoshī).
\item \textbf{Mélanger \zh{在} (zài) et \zh{有} (yǒu) pour la localisation :} \zh{我在学校。} (Wǒ zài xuéxiào.) « Je suis à l'école » (position du sujet). \zh{学校有学生。} (Xuéxiào yǒu xuésheng.) « Il y a des élèves dans l'école » (existence). Le sujet et le sens sont différents !
\item \textbf{Les tons de \zh{一} :} seul, dans un nombre ou un ordinal, il garde le \textbf{1er ton} (\zh{第一} dìyī). Devant un 4ème ton, il prend le \textbf{2ème ton} : \zh{一个} (yí ge), \zh{一样} (yíyàng). Devant un 1er, 2ème ou 3ème ton, il prend le \textbf{4ème ton} : \zh{一本} (yì běn), \zh{一年} (yì nián), \zh{一起} (yìqǐ).
\item \textbf{\zh{和} (hé) ne relie que des noms :} \zh{我和同学} (wǒ hé tóngxué) « mes camarades et moi ». Pour relier des \textbf{phrases}, n'utilisez pas \zh{和} : employez \zh{也} (yě, « aussi »), \zh{还} (hái, « de plus ») ou une simple virgule : \zh{我喜欢看电视，也喜欢看书。}
\end{itemize}
\end{attention}

\begin{experience}[Activité orale : « le portrait chinois » — 3 minutes par binôme]
\begin{enumerate}
\item L'élève A pioche une « carte d'identité » fictive (nom, âge, ville, famille, hobby, heure de lever et de coucher).
\item L'élève B l'interroge en utilisant \textbf{au moins cinq mots interrogatifs} différents (\zh{什么}, \zh{谁}, \zh{哪儿}, \zh{几}, \zh{怎么}, \zh{多少}, \zh{几岁}...).
\item L'élève A répond par des phrases complètes, avec classificateurs.
\item Inversion des rôles. Le professeur évalue : justesse des tons, présence des classificateurs, ordre des mots dans les questions.
\end{enumerate}
\emph{Exemple de carte :} \zh{王小红，16岁，杜阿拉，4口人，喜欢听音乐，6:30起床，10:00睡觉。} (Wáng Xiǎohóng, shíliù suì, Dù'ālā, sì kǒu rén, xǐhuan tīng yīnyuè, liù diǎn bàn qǐchuáng, shí diǎn shuìjiào.) — Wang Xiaohong, 16 ans, Douala, famille de 4 personnes, aime écouter de la musique, se lève à 6 h 30, se couche à 22 h.
\end{experience}

\coursec{Synthèse des structures grammaticales}

Après avoir fait le bilan du vocabulaire de l'année, passons maintenant à l'architecture de la phrase chinoise. Cette synthèse reprend, une à une, toutes les structures grammaticales étudiées depuis septembre : la phrase d'identité avec \zh{是}, la possession avec \zh{有}, la localisation avec \zh{在}, la particule \zh{的}, les questions, les négations \zh{不} et \zh{没}, l'accompli avec \zh{了}, l'ordre des mots et l'expression de l'heure et des dates. L'objectif : être capable de construire, de corriger et de traduire n'importe quelle phrase simple du programme de Première, comme on vous le demandera à l'examen.

\courssub{Texte d'observation : la journée de Marthe}

Lisons d'abord ce court texte. Toutes les structures de l'année y apparaissent. Soulignez mentalement les mots-outils \zh{是}, \zh{有}, \zh{在}, \zh{的}, \zh{了}, \zh{没}.

\begin{textebox}[La journée de Marthe]
玛尔特是雅温得一所中学的学生。\\
\textit{Mǎ'ěrtè shì Yǎwēndé yì suǒ zhōngxué de xuéshēng.}\\
« Marthe est élève d'un collège de Yaoundé. »\\[0.3em]
她家有四口人：爸爸、妈妈、哥哥和她。\\
\textit{Tā jiā yǒu sì kǒu rén : bàba, māma, gēge hé tā.}\\
« Sa famille compte quatre personnes : papa, maman, son grand frère et elle. »\\[0.3em]
她没有弟弟，也没有妹妹。\\
\textit{Tā méiyǒu dìdi, yě méiyǒu mèimei.}\\
« Elle n'a pas de petit frère, ni de petite sœur. »\\[0.3em]
她每天在学校学习汉语。\\
\textit{Tā měitiān zài xuéxiào xuéxí Hànyǔ.}\\
« Chaque jour, elle apprend le chinois à l'école. »\\[0.3em]
昨天她看了一个中国电影，没看电视。\\
\textit{Zuótiān tā kàn le yí ge Zhōngguó diànyǐng, méi kàn diànshì.}\\
« Hier, elle a vu un film chinois, elle n'a pas regardé la télé. »\\[0.3em]
晚上八点，她在家里做作业。\\
\textit{Wǎnshang bā diǎn, tā zài jiā lǐ zuò zuòyè.}\\
« À huit heures du soir, elle fait ses devoirs à la maison. »\\[0.3em]
你呢？你昨天做什么了？\\
\textit{Nǐ ne ? Nǐ zuótiān zuò shénme le ?}\\
« Et toi ? Qu'as-tu fait hier ? »
\end{textebox}

\textbf{Questions d'observation}
\begin{enumerate}
\item Dans la première phrase, quel mot relie « Marthe » et « élève » ?
\item Comment dit-on « elle n'a pas » ? Quel mot précède \zh{有} ?
\item Dans « elle apprend le chinois à l'école », où se place le complément de lieu \zh{在学校} dans la phrase chinoise ?
\item Quel mot marque que l'action « regarder un film » est accomplie ?
\end{enumerate}

\textbf{Réponses attendues} : 1) \zh{是} ; 2) \zh{没有}, le mot \zh{没} précède \zh{有} ; 3) avant le verbe \zh{学习} ; 4) la particule \zh{了}, placée après le verbe \zh{看}. Vérifions maintenant chaque structure en détail.

\courssub{La phrase d'identité avec 是}

\begin{propriete}[La phrase d'identité : Sujet + 是 + nom]
Le verbe \zh{是} (shì) relie un sujet à un \textbf{nom} (identité, profession, nationalité). Structure : \cle{Sujet + 是 + nom}. Négation : \zh{不是} (bú shì).
\begin{itemize}
\item \zh{我是学生。} \textit{Wǒ shì xuéshēng.} « Je suis élève. »
\item \zh{她是喀麦隆人。} \textit{Tā shì Kāmàilóng rén.} « Elle est Camerounaise. »
\item \zh{他不是老师，他是医生。} \textit{Tā bú shì lǎoshī, tā shì yīshēng.} « Il n'est pas professeur, il est médecin. »
\end{itemize}
\textbf{Exception majeure} : on n'utilise jamais \zh{是} devant un adjectif. On dit \zh{我很好} (Wǒ hěn hǎo, « je vais bien »), jamais \zh{我是很好}. Devant un adjectif, c'est l'adverbe \zh{很} (hěn) qui fait le lien, sans valeur intensive réelle.
\end{propriete}

\courssub{La possession et l'existence avec 有}

\begin{propriete}[La négation de 有 est toujours 没有]
\zh{有} (yǒu) exprime la possession (« avoir ») et l'existence (« il y a »). Sa négation est \textbf{toujours} \zh{没有} (méiyǒu), jamais \zh{不有}.
\begin{itemize}
\item \zh{我家有五口人。} \textit{Wǒ jiā yǒu wǔ kǒu rén.} « Ma famille compte cinq personnes. »
\item \zh{我没有手机。} \textit{Wǒ méiyǒu shǒujī.} « Je n'ai pas de téléphone. »
\item \zh{我没有弟弟。} \textit{Wǒ méiyǒu dìdi.} « Je n'ai pas de petit frère. »
\item \zh{桌子上有一本书。} \textit{Zhuōzi shang yǒu yì běn shū.} « Il y a un livre sur la table. »
\end{itemize}
\textbf{Rappel} : le classificateur des personnes d'une famille est \zh{口} (kǒu) ; celui des livres est \zh{本} (běn). Devant un classificateur, « deux » se dit \zh{两} (liǎng) : \zh{两本书} « deux livres », jamais \zh{二本书}.
\end{propriete}

\courssub{La localisation avec 在}

\zh{在} (zài) est le verbe « être (quelque part) » ou la préposition « à, dans, sur ». Il introduit le lieu, placé \textbf{avant} le verbe d'action.

\begin{itemize}
\item \zh{我在学校。} \textit{Wǒ zài xuéxiào.} « Je suis à l'école. »
\item \zh{书在桌子上。} \textit{Shū zài zhuōzi shang.} « Le livre est sur la table. »
\item \zh{我们在杜阿拉学习。} \textit{Wǒmen zài Dù'ālā xuéxí.} « Nous étudions à Douala. »
\end{itemize}

\begin{attention}[是 ou 在 ? Le piège du francophone]
Erreur fréquente : traduire « Je suis à l'école » par \zh{我是学校}, ce qui signifierait littéralement « je suis (le bâtiment) école » ! Règle : \zh{是} relie \textbf{deux noms} (identité) ; \zh{在} indique \textbf{le lieu}. « Je suis élève » → \zh{我是学生。} ; « Je suis à l'école » → \zh{我在学校。}
\end{attention}

\courssub{La particule 的 : possession et qualification}

\begin{propriete}[La particule 的]
\zh{的} (de) se place entre un déterminant (possesseur, qualificatif) et le nom : \cle{X + 的 + nom}.
\begin{itemize}
\item Possession : \zh{我的老师} \textit{wǒ de lǎoshī} « mon professeur » ; \zh{玛尔特的书} \textit{Mǎ'ěrtè de shū} « le livre de Marthe ».
\item Qualification : \zh{雅温得的中学} \textit{Yǎwēndé de zhōngxué} « un collège de Yaoundé » ; \zh{中国的电影} \textit{Zhōngguó de diànyǐng} « un film chinois ».
\end{itemize}
\textbf{Cas particulier} : après un pronom personnel suivi d'un lien de parenté ou d'un lieu très proche, \zh{的} peut tomber : \zh{我妈妈} (wǒ māma, « ma maman »), \zh{我们家} (wǒmen jiā, « notre maison »).
\end{propriete}

\courssub{Les trois types de questions}

\begin{propriete}[Former une question en chinois]
\begin{enumerate}
\item \textbf{Question oui/non} : phrase affirmative + \zh{吗} (ma). \zh{你是学生吗？} \textit{Nǐ shì xuéshēng ma?} « Es-tu élève ? » Réponse : \zh{是。/ 不是。}
\item \textbf{Question avec mot interrogatif} : le mot interrogatif (\zh{什么}, \zh{谁}, \zh{哪儿}, \zh{几}, \zh{怎么}) \textbf{reste à la place exacte de la réponse}, sans inversion. \zh{你叫什么？} \textit{Nǐ jiào shénme?} « Comment t'appelles-tu ? » → \zh{我叫李明。} \textit{Wǒ jiào Lǐ Míng.}
\item \textbf{Question renvoyée avec 呢} : \zh{呢} (ne) retourne la question : \zh{我很好，你呢？} \textit{Wǒ hěn hǎo, nǐ ne?} « Je vais bien, et toi ? »
\end{enumerate}
\end{propriete}

\begin{savaistu}[呢, la question en un seul mot]
En chinois, un seul petit mot suffit pour retourner une question : \zh{呢} (ne). Après avoir répondu « Je vais bien », il suffit de dire \zh{你呢？} (nǐ ne ?) pour demander « et toi ? ». Économique et très utilisé dans les dialogues : \zh{我去学校，你呢？} « Moi je vais à l'école, et toi ? »
\end{savaistu}

\courssub{Les négations 不 et 没}

\begin{methode}[Choisir entre 不 et 没]
\begin{enumerate}
\item La phrase contient-elle \zh{有} (possession) ? → toujours \zh{没} : \zh{我没有弟弟。}
\item L'action est-elle passée et non accomplie ? → \zh{没} : \zh{我没去。} \textit{Wǒ méi qù.} « Je n'y suis pas allé. »
\item Présent, futur, goûts, habitudes, états ? → \zh{不} : \zh{我不去。} \textit{Wǒ bú qù.} « Je n'y vais pas. » ; \zh{我不喜欢咖啡。} \textit{Wǒ bù xǐhuan kāfēi.} « Je n'aime pas le café. »
\end{enumerate}
\end{methode}

\begin{center}
\begin{tabularx}{\linewidth}{|X|X|X|X|}\hline
\rowcolor{popblueL} Négation & Emploi & Exemple chinois + pinyin & Traduction \\ \hline
\zh{不} (bù) & présent, futur, goûts, états & \zh{我不去。} Wǒ bú qù. & Je n'y vais pas. \\ \hline
\zh{不} (bù) & négation de \zh{是} & \zh{他不是老师。} Tā bú shì lǎoshī. & Il n'est pas professeur. \\ \hline
\zh{没} (méi) & action passée non faite & \zh{我没看电视。} Wǒ méi kàn diànshì. & Je n'ai pas regardé la télé. \\ \hline
\zh{没有} (méiyǒu) & possession, existence & \zh{我没有手机。} Wǒ méiyǒu shǒujī. & Je n'ai pas de téléphone. \\ \hline
\end{tabularx}
\end{center}

\courssub{L'accompli avec 了}

\begin{propriete}[La particule 了 d'accompli]
\zh{了} (le) placé \textbf{après le verbe} marque une action accomplie (souvent au passé). Structure : \cle{Sujet + Verbe + 了 + complément}.
\begin{itemize}
\item \zh{我昨天看了一个电影。} \textit{Wǒ zuótiān kàn le yí ge diànyǐng.} « Hier j'ai vu un film. »
\item \zh{你昨天做什么了？} \textit{Nǐ zuótiān zuò shénme le?} « Qu'as-tu fait hier ? »
\item \zh{她买了两本书。} \textit{Tā mǎi le liǎng běn shū.} « Elle a acheté deux livres. »
\end{itemize}
\textbf{Négation} : une action non accomplie se nie avec \zh{没} (ou \zh{没有}) et \textbf{sans} \zh{了} : \zh{我没看电影。} « Je n'ai pas vu de film. »
\end{propriete}

\begin{attention}[了 n'est pas un passé automatique]
\zh{了} marque l'accomplissement, pas simplement le passé. On ne le met pas avec les verbes d'état ni les habitudes : on dit \zh{我每天学习汉语。} (sans \zh{了}) pour « chaque jour j'étudie le chinois ». Et à la négative, \zh{了} disparaît : \zh{我没去。}, jamais \zh{我没去了。}
\end{attention}

\courssub{L'ordre des mots : la grande règle de l'année}

En français, on dit « J'apprends le chinois \textbf{à l'école} \textbf{tous les jours} » : les compléments de temps et de lieu viennent après. En chinois, c'est l'inverse : le cadre (temps, lieu) se place \textbf{avant le verbe}. C'est la règle la plus rentable de l'année : elle conditionne la réussite de toutes les traductions de l'examen.

\begin{propriete}[Ordre canonique de la phrase chinoise]
\cle{Sujet + Temps + Lieu + Verbe + Complément}. Le temps et le lieu se placent entre le sujet et le verbe, dans cet ordre : d'abord le temps, ensuite le lieu. L'adverbe \zh{也} (yě, « aussi ») se place juste après le sujet, avant la négation ou le verbe : \zh{我也没有妹妹。} « Moi non plus je n'ai pas de petite sœur. »
\end{propriete}

\begin{popfigure}
\begin{tikzpicture}[node distance=0.35cm, every node/.style={font=\small}]
\node[draw=popblue, fill=popblueL, rounded corners, minimum width=1.7cm, minimum height=0.9cm, align=center] (s) {Sujet\\ \zh{我}};
\node[draw=poporange, fill=poporangeL, rounded corners, minimum width=1.7cm, minimum height=0.9cm, align=center, right=of s] (t) {Temps\\ \zh{每天}};
\node[draw=popgreen, fill=popgreenL, rounded corners, minimum width=1.9cm, minimum height=0.9cm, align=center, right=of t] (l) {Lieu\\ \zh{在学校}};
\node[draw=poppurple, fill=poppurpleL, rounded corners, minimum width=1.7cm, minimum height=0.9cm, align=center, right=of l] (v) {Verbe\\ \zh{学习}};
\node[draw=poppink, fill=poppinkL, rounded corners, minimum width=1.7cm, minimum height=0.9cm, align=center, right=of v] (c) {Complément\\ \zh{汉语}};
\draw[-{Stealth}, thick, popdark] (s) -- (t);
\draw[-{Stealth}, thick, popdark] (t) -- (l);
\draw[-{Stealth}, thick, popdark] (l) -- (v);
\draw[-{Stealth}, thick, popdark] (v) -- (c);
\node[below=0.5cm of l, align=center, font=\itshape] {Wǒ měitiān zài xuéxiào xuéxí Hànyǔ.\\ « Chaque jour, j'apprends le chinois à l'école. »};
\end{tikzpicture}
\legende{Schéma de la phrase chinoise : le temps et le lieu précèdent le verbe.}
\end{popfigure}

\begin{center}
\begin{tabularx}{\linewidth}{|X|X|X|}\hline
\rowcolor{popblueL} Structure & Exemple chinois + pinyin & Traduction \\ \hline
Sujet + Verbe + Complément & \zh{我学习汉语。} Wǒ xuéxí Hànyǔ. & J'apprends le chinois. \\ \hline
Sujet + Temps + Verbe & \zh{我每天学习。} Wǒ měitiān xuéxí. & J'étudie chaque jour. \\ \hline
Sujet + Temps + Lieu + Verbe + Complément & \zh{我每天在学校学习汉语。} Wǒ měitiān zài xuéxiào xuéxí Hànyǔ. & Chaque jour, j'apprends le chinois à l'école. \\ \hline
\end{tabularx}
\end{center}

\courssub{L'heure et les dates : du plus grand au plus petit}

Le chinois procède toujours \textbf{du plus grand au plus petit}, à l'inverse du français : année → mois → jour ; moment de la journée → heure.

\begin{itemize}
\item \zh{2024年5月10日} \textit{èr líng èr sì nián wǔ yuè shí rì} « le 10 mai 2024 » (litt. « année 2024, mois 5, jour 10 »).
\item \zh{晚上八点} \textit{wǎnshang bā diǎn} « huit heures du soir » (litt. « soir, huit heures »).
\item \zh{我晚上八点做作业。} \textit{Wǒ wǎnshang bā diǎn zuò zuòyè.} « Je fais mes devoirs à huit heures du soir. »
\end{itemize}

\begin{attention}[L'heure : le moment de la journée d'abord]
En français on dit « à huit heures du soir » ; en chinois, le moment de la journée (\zh{早上} matin, \zh{下午} après-midi, \zh{晚上} soir) précède toujours l'heure : \zh{下午三点} « trois heures de l'après-midi ». Ne dites jamais \zh{三点下午}.
\end{attention}

\courssub{Exercice résolu et entraînement}

\begin{exoresolu}[Traduction et correction]
\textbf{Énoncé.} 1) Traduisez en chinois : « Hier, je n'ai pas regardé la télé à la maison. » 2) Corrigez la phrase fautive : \zh{我是学校。} (sens voulu : « Je suis à l'école. ») 3) Traduisez : \zh{你昨天做什么了？}
\tcblower
\textbf{Solution détaillée.} 1) Action passée non faite → négation \zh{没} ; le lieu \zh{在家} se place avant le verbe ; le temps \zh{昨天} avant le lieu : \zh{我昨天在家没看电视。} \textit{Wǒ zuótiān zài jiā méi kàn diànshì.} 2) \zh{是} relie deux noms ; pour un lieu, il faut \zh{在} : \zh{我在学校。} \textit{Wǒ zài xuéxiào.} 3) \zh{你昨天做什么了？} \textit{Nǐ zuótiān zuò shénme le?} « Qu'as-tu fait hier ? » — le mot interrogatif \zh{什么} occupe la place du complément, et \zh{了} marque l'accompli.
\end{exoresolu}

\begin{exercice}[À vous de jouer]
\begin{enumerate}
\item Traduisez en français : \zh{我家有四口人。} / \zh{我没有妹妹。} / \zh{她每天在学校学习汉语。}
\item Traduisez en chinois : « Mon professeur est Chinois. » ; « Je n'ai pas de téléphone. » ; « Hier, j'ai vu un film à Douala. »
\item Choisissez \zh{不} ou \zh{没} : 我昨天\_\_\_去市场。/ 他\_\_\_喜欢咖啡。/ 我\_\_\_有哥哥。
\item Remettez dans l'ordre : \zh{学习 / 我 / 在学校 / 汉语 / 每天}.
\end{enumerate}
\end{exercice}

\begin{corrige}[Exercice « À vous de jouer »]
1) « Ma famille compte quatre personnes. » / « Je n'ai pas de petite sœur. » / « Chaque jour, elle apprend le chinois à l'école. » 2) \zh{我的老师是中国人。} \textit{Wǒ de lǎoshī shì Zhōngguó rén.} ; \zh{我没有手机。} \textit{Wǒ méiyǒu shǒujī.} ; \zh{我昨天在杜阿拉看了一个电影。} \textit{Wǒ zuótiān zài Dù'ālā kàn le yí ge diànyǐng.} 3) \zh{没} (passé) ; \zh{不} (goût) ; \zh{没} (possession → \zh{没有}). 4) \zh{我每天在学校学习汉语。} — Sujet + Temps + Lieu + Verbe + Complément.
\end{corrige}

\begin{aretenir}[L'essentiel de la synthèse grammaticale]
\begin{itemize}
\item Identité : \zh{是} + nom (négation \zh{不是}) ; jamais \zh{是} devant un adjectif (on utilise \zh{很}).
\item Possession : \zh{有}, négation obligatoire \zh{没有} ; « deux » + classificateur = \zh{两}.
\item Lieu : \zh{在}, placé avant le verbe ; ne pas confondre avec \zh{是}.
\item \zh{的} relie possesseur ou qualificatif au nom ; il peut tomber après un pronom devant un proche.
\item Questions : \zh{吗} pour oui/non ; mot interrogatif à la place de la réponse ; \zh{呢} pour renvoyer la question.
\item \zh{不} = présent, futur, goûts ; \zh{没} = passé et possession.
\item \zh{了} après le verbe = action accomplie ; il disparaît à la négative.
\item Ordre d'or : Sujet + Temps + Lieu + Verbe + Complément ; dates et heures du plus grand au plus petit.
\end{itemize}
\end{aretenir}

\coursec{Caractères : radicaux, traits, pièges}

Après avoir synthétisé les structures grammaticales qui font la phrase chinoise, nous devons consolider l'ossature même de l'écriture : le caractère. Maîtriser les \cle{radicaux} (部首, bùshǒu), respecter l'\cle{ordre des traits} (笔顺, bǐshùn) et discriminer les \cle{caractères proches} (易混淆字, yì hùnxiáo zì) sont les trois piliers qui permettent de passer de la reconnaissance passive à la production active et correcte, indispensable pour l'examen écrit.

\courssub{Les radicaux clés de l'année : révision et classification}

\begin{definition}[Le radical (部首, bùshǒu)]
Un \textbf{radical} (ou \cle{clé de dictionnaire}) est un composant graphique récurrent qui sert à classer les caractères dans les dictionnaires. Il donne souvent — mais pas toujours — une indication sur le \textbf{sens} (clé sémantique) ou sur la \textbf{prononciation} (clé phonétique) du caractère. Connaître les radicaux les plus fréquents permet de deviner le sens d'un caractère inconnu et de le retrouver rapidement.
\end{definition}

\begin{textebox}[Révision des 9 radicaux majeurs de l'année]
Voici les neuf radicaux étudiés en Première, classés par domaine sémantique, avec pour chacun deux caractères exemples vus cette année.

\begin{center}
\begin{tabularx}{\linewidth}{|X|X|X|X|}
\hline
\rowcolor{popblueL} \textbf{Radical (forme)} & \textbf{Nom / Sens} & \textbf{Exemple 1} & \textbf{Exemple 2} \\
\hline
人 / 亻 & \zh{人} (rén) : l'homme, la personne & \zh{你} (nǐ) toi & \zh{他} (tā) lui \\
\hline
女 & \zh{女} (nǚ) : la femme, le féminin & \zh{妈} (mā) maman & \zh{好} (hǎo) bon \\
\hline
口 & \zh{口} (kǒu) : la bouche, l'ouverture & \zh{吃} (chī) manger & \zh{听} (tīng) écouter \\
\hline
日 & \zh{日} (rì) : le soleil, le jour & \zh{早} (zǎo) matin & \zh{明} (míng) clair / lendemain \\
\hline
月 & \zh{月} (yuè) : la lune, le mois & \zh{有} (yǒu) avoir & \zh{期} (qī) période, semaine \\
\hline
氵 (水) & \zh{水} (shuǐ) : l'eau, liquide & \zh{河} (hé) fleuve & \zh{喝} (hē) boire \\
\hline
讠 (言) & \zh{言} (yán) : la parole, le langage & \zh{语} (yǔ) langue & \zh{说} (shuō) parler \\
\hline
心 / 忄 & \zh{心} (xīn) : le cœur, le sentiment & \zh{想} (xiǎng) penser & \zh{怕} (pà) avoir peur \\
\hline
宀 & \zh{宀} (mián) : le toit, la maison & \zh{家} (jiā) famille/maison & \zh{安} (ān) paisible \\
\hline
\end{tabularx}
\end{center}
\end{textebox}

\begin{savaistu}[Le pouvoir des radicaux]
Connaître une vingtaine de radicaux « productifs » (comme ceux ci-dessus) permet de \textbf{deviner le sens approximatif de plusieurs centaines de caractères}. Par exemple, si vous voyez le radical \zh{氵} (eau) dans un caractère inconnu comme \zh{湖} (hú), \zh{海} (hǎi) ou \zh{汤} (tāng), vous pouvez inférer qu'il s'agit d'un liquide, d'un plan d'eau ou d'une action liée à l'eau. C'est la clé pour l'autonomie en lecture.
\end{savaistu}

\courssub{Règles d'ordre des traits : rappels et application}

\begin{propriete}[Les 5 règles d'or de l'ordre des traits (笔顺规则)]
Respecter l'ordre des traits n'est pas une contrainte calligraphique : il garantit l'\textbf{équilibre} du caractère, la \textbf{vitesse} d'écriture et la \textbf{reconnaissance} par les logiciels OCR (reconnaissance d'écriture manuscrite).
\begin{enumerate}
    \item \textbf{De haut en bas} (从上到下) : \zh{三} (sān), \zh{王} (wáng).
    \item \textbf{De gauche à droite} (从左到右) : \zh{川} (chuān), \zh{好} (hǎo).
    \item \textbf{Horizontal avant vertical} (先横后竖) : \zh{十} (shí), \zh{干} (gàn).
    \item \textbf{Extérieur avant intérieur} (先外后内) : \zh{月} (yuè), \zh{同} (tóng).
    \item \textbf{On ferme en dernier} (最后封口) : \zh{国} (guó), \zh{田} (tián) — le trait de fermeture horizontal s'écrit à la fin.
\end{enumerate}
\end{propriete}

\begin{methode}[Écrire 家 (jiā) : « Un cochon sous un toit »]
Le caractère \zh{家} (famille, maison, suffixe de spécialiste) est un idéogramme composé : \zh{宀} (toit) + \zh{豕} (shǐ, cochon, porc). Dans la Chine ancienne, la présence d'un cochon sous le toit symbolisait la prospérité du foyer.
\begin{enumerate}
    \item \textbf{Le radical \zh{宀} (toit) — 3 traits :} Point (丶) → Trait horizontal court (一) → Trait horizontal long avec crochet final (一乛). \emph{Astuce : le toit protège, il se pose en premier.}
    \item \textbf{Le composant \zh{豕} (cochon) — 7 traits :} Trait horizontal (一) → Trait vertical (丨) → Horizontal (一) → Horizontal (一) → Vertical (丨) → Horizontal (一) → Trait oblique gauche (丿) + point final (丶) \emph{(souvent écrits liés)}.
    \item \textbf{Total : 10 traits.} Vérifiez que le « cochon » ne dépasse pas du « toit » sur les côtés.
\end{enumerate}
\end{methode}

\begin{exemplebox}[Étude détaillée de quatre caractères fondamentaux]
\begin{center}
\begin{tabularx}{\linewidth}{|X|X|X|X|}
\hline
\rowcolor{popblueL} \textbf{Caractère} & \textbf{Nb traits / Radical} & \textbf{Décomposition / Ordre clé} & \textbf{Exemple contextuel} \\
\hline
\zh{我} (wǒ) & 7 traits \\ Radical : \zh{戈} (gē, hallebarde) & 1. \zh{手} (main, 3 traits) :横、竖钩、提 \\ 2. \zh{戈} (4 traits) :横、竖、撇、捺 & \zh{我是喀麦隆学生。} (Wǒ shì Kāmàilóng xuésheng.) — Je suis élève au Cameroun. \\
\hline
\zh{学} (xué) & 8 traits (simplifié) \\ Radical : \zh{子} (zǐ, enfant) & 1. \zh{宀} (toit, 3 traits) \\ 2. \zh{子} (enfant, 3 traits) :横、竖钩、横 \\ 3. \zh{夂} (pied/arrêt, 2 traits) :撇、捺 & \zh{我喜欢学中文。} (Wǒ xǐhuan xué Zhōngwén.) — J'aime apprendre le chinois. \\
\hline
\zh{语} (yǔ) & 9 traits \\ Radical : \zh{讠} (yán, parole) & 1. \zh{讠} (2 traits) :点、竖折 \\ 2. \zh{吾} (wú, 7 traits) :横、竖、横、横、撇、横、捺 & \zh{中文是一门美丽的语言。} (Zhōngwén shì yī mén měilì de yǔyán.) — Le chinois est une belle langue. \\
\hline
\zh{家} (jiā) & 10 traits \\ Radical : \zh{宀} (mián, toit) & Voir \textbf{Méthode} ci-dessus (宀 + 豕) & \zh{我的家在雅温得。} (Wǒ de jiā zài Yǎwēndé.) — Ma maison est à Yaoundé. \\
\hline
\end{tabularx}
\end{center}
\end{exemplebox}

\begin{exoresolu}[Analyse de l'ordre des traits : \zh{国} (guó) vs \zh{园} (yuán)]
\textbf{Énoncé :} Écrivez dans l'ordre les numéros des traits pour \zh{国} et \zh{园}. Expliquez la différence du composant intérieur.
\tcblower
\textbf{Solution détaillée :}
\begin{itemize}
    \item \zh{国} (guó, pays) : 8 traits. Structure \textbf{全包围} (enfermement total). Ordre : 1-2-3 (外 \zh{囗} :横、竖折、横) → 4-5-6-7 (内 \zh{玉} :横、横、竖、横) → 8 (封口 :横). \emph{Règle : extérieur → intérieur → fermeture.}
    \item \zh{园} (yuán, jardin) : 7 traits. Structure \textbf{全包围}. Ordre : 1-2-3 (外 \zh{囗} : identique) → 4-5-6-7 (内 \zh{元} :横、竖折、撇、捺). Pas de trait de fermeture final car \zh{元} touche le bas de l'enceinte.
    \item \textbf{Piège :} Ne pas écrire le \zh{口} (kǒu) carré parfait pour l'enceinte \zh{囗} (wéi) : le trait du bas (fermeture) doit être plus long que le trait du haut pour « porter » le contenu.
\end{itemize}
\end{exoresolu}

\courssub{Caractères proches : pièges visuels et discrimination}

\begin{attention}[Les paires « jumelles » à ne jamais confondre]
Les francophones confondent souvent ces paires par manque d'attention au \textbf{détail distinctif} (un trait manquant, une proportion changée). En examen, écrire \zh{大} pour \zh{太} ou \zh{未} pour \zh{末} change le sens de la phrase et coûte des points.

\begin{popfigure}
\legende{Tableau des caractères confondables : discrimination visuelle}
\begin{center}
\begin{tabularx}{\linewidth}{|X|X|X|X|X|}
\hline
\rowcolor{poporangeL} \textbf{Paire} & \textbf{Caractère A} & \textbf{Sens A} & \textbf{Caractère B} & \textbf{Sens B} \\
\hline
1 & \zh{人} (rén) & Personne, homme & \zh{入} (rù) & Entrer \\
\hline
2 & \zh{大} (dà) & Grand & \zh{太} (tài) & Trop, très (point en plus) \\
\hline
3 & \zh{日} (rì) & Soleil, jour (rectangle plat) & \zh{目} (mù) & Œil (rectangle haut, deux traits intérieurs) \\
\hline
4 & \zh{未} (wèi) & Pas encore (trait court en haut) & \zh{末} (mò) & Fin, bout (trait long en haut) \\
\hline
5 & \zh{我} (wǒ) & Je, moi & \zh{找} (zhǎo) & Chercher (main \zh{扌} à gauche) \\
\hline
6 & \zh{问} (wèn) & Demander (porte \zh{门}) & \zh{间} (jiān) & Intervalle, entre (porte \zh{门} + soleil \zh{日}) \\
\hline
\end{tabularx}
\end{center}
\end{popfigure}

\end{attention}

\begin{exemplebox}[Discrimination en contexte]
\begin{enumerate}
    \item \zh{今天\underline{太}热了！} (Jīntiān \underline{tài} rè le !) — Il fait \textbf{trop} chaud aujourd'hui ! (Pas \zh{大}).
    \item \zh{我\underline{未}做完作业。} (Wǒ \underline{wèi} zuò wán zuòyè.) — Je n'ai \textbf{pas encore} fini mes devoirs. (Pas \zh{末}).
    \item \zh{请\underline{问}，图书馆在哪儿？} (Qǐng \underline{wèn}, túshūguǎn zài nǎr ?) — Excusez-moi (\textbf{demander}), où est la bibliothèque ?
    \item \zh{中\underline{间}休息十分钟。} (Zhōng \underline{jiān} xiūxi shí fēnzhōng.) — Pause de 10 minutes \textbf{entre} les deux / au milieu.
\end{enumerate}
\end{exemplebox}

\courssub{Stratégies de mémorisation : du pinyin au caractère}

\begin{propriete}[Analyse par composants : Clé sémantique + Clé phonétique]
Plus de 80\% des caractères sont des \textbf{phono-idéogrammes} (形声字, xíngshēngzì) : \cle{Clé de sens (radical)} + \cle{Clé de son (phonétique)}.
\begin{center}
\textbf{Sens (Radical)} + \textbf{Son (Phonétique)} = \textbf{Caractère}
\end{center}
Exemple : \zh{河} (hé, fleuve) = \zh{氵} (eau, sens) + \zh{可} (kě, son approximatif).
\end{propriete}

\begin{methode}[Comment mémoriser un nouveau caractère en 3 étapes]
\begin{enumerate}
    \item \textbf{Décomposer :} Identifier le radical (sens) et le composant phonétique (son). Ex: \zh{妈} = \zh{女} + \zh{马}.
    \item \textbf{Créer une histoire (mnémotechnique) :} Lier le sens du radical au sens du caractère via le son. \emph{« La \zh{女} (femme) qui fait \zh{马} (ma) est \zh{妈} (maman). »}
    \item \textbf{Écrire en dictant l'ordre :} Dire les noms des traits à voix haute (横、竖、撇、捺...) pour ancrer la motricité fine.
\end{enumerate}
\end{methode}

\begin{exemplebox}[Le cas modèle : 好 (hǎo) = 女 + 子]
\begin{popfigure}
\legende{Décomposition idéographique de 好 (hǎo) : Femme + Enfant = Bon / Aimer}
\begin{tikzpicture}[node distance=1.2cm, every node/.style={font=\small, align=center}]
    % Nœuds
    \node[draw=poppink, thick, fill=poppinkL, rounded corners, minimum width=2.5cm, minimum height=1.5cm, align=center] (nv){\zh{女}\\(nǚ)\\Femme};
    \node[draw=popblue, thick, fill=popblueL, rounded corners, minimum width=2.5cm, minimum height=1.5cm, right=of nv, align=center] (zi){\zh{子}\\(zǐ)\\Enfant};
    \node[draw=popgreen, very thick, fill=popgreenL, rounded corners, minimum width=3cm, minimum height=1.8cm, below=of $(nv.south)!0.5!(zi.south)$, align=center] (hao){\zh{好}\\(hǎo)\\Bon, aimer};
    
    % Flèches
    \draw[-Stealth, thick, popdark] (nv.south) -- ++(0,-0.3) -| (hao.north west) node[midway, left, font=\tiny] {Clé de sens};
    \draw[-Stealth, thick, popdark] (zi.south) -- ++(0,-0.3) -| (hao.north east) node[midway, right, font=\tiny] {Clé de son (ancien)};
    
    % Annotation
    \node[below=0.2cm of hao, font=\footnotesize, text=popmuted] {Idée : Une femme avec son enfant = harmonie, bonté.};
\end{tikzpicture}
\end{popfigure}
\end{exemplebox}

\begin{aretenir}[Autres associations utiles de l'année]
\begin{itemize}
    \item \zh{休} (xiū, se reposer) = \zh{人} (homme) + \zh{木} (arbre) → L'homme appuyé à l'arbre se repose.
    \item \zh{林} (lín, bois) = \zh{木} + \zh{木} ; \zh{森} (sēn, forêt) = \zh{木} × 3.
    \item \zh{明} (míng, clair/lumineux) = \zh{日} (soleil) + \zh{月} (lune) → Les deux astres ensemble éclairent.
    \item \zh{尖} (jiān, pointu) = \zh{小} (petit) + \zh{大} (grand) → Le petit sur le grand fait une pointe.
\end{itemize}
\end{aretenir}

\courssub{Dictée blanche : 10 caractères fréquents de l'année}

\begin{exercice}[Dictée blanche — Épreuve d'entraînement]
\textbf{Consigne :} Écrivez les \textbf{caractères chinois} (simplifiés) correspondant aux définitions ou pinyin ci-dessous. Respectez l'ordre des traits. Pour chaque caractère, indiquez son \textbf{radical} principal.
\begin{enumerate}
    \item [wǒ] — Pronom « je, moi ». Radical : \_\_\_\_\_\_
    \item [xié] — Verbe « écrire ». Radical : \_\_\_\_\_\_
    \item [tīng] — Verbe « écouter ». Radical : \_\_\_\_\_\_
    \item [jiā] — Nom « famille, maison » ; suffixe de métier. Radical : \_\_\_\_\_\_
    \item [huì] — Verbe « savoir (une compétence) ; réunion ». Radical : \_\_\_\_\_\_
    \item [yǔ] — Nom « langue, langage ». Radical : \_\_\_\_\_\_
    \item [páng] — Adjectif « à côté, à côté de ». Radical : \_\_\_\_\_\_
    \item [qīng] — Adjectif « clair, léger ; vert/bleu ». Radical : \_\_\_\_\_\_
    \item [xiǎng] — Verbe « penser, vouloir, manquer à ». Radical : \_\_\_\_\_\_
    \item [guó] — Nom « pays, nation ». Radical : \_\_\_\_\_\_
\end{enumerate}
\end{exercice}

\begin{corrige}[Exercice : Dictée blanche]
\begin{center}
\begin{tabularx}{\linewidth}{|c|X|X|X|}
\hline
\rowcolor{popblueL} \textbf{N°} & \textbf{Caractère} & \textbf{Pinyin / Sens} & \textbf{Radical (Clé)} \\
\hline
1 & \zh{我} & wǒ — je, moi & \zh{戈} (gē) \\
\hline
2 & \zh{写} & xiě — écrire & \zh{冖} (mì) \emph{ou} \zh{与} (yǔ) \emph{(classement moderne souvent sous 与)} \\
\hline
3 & \zh{听} & tīng — écouter & \zh{口} (kǒu) \\
\hline
4 & \zh{家} & jiā — famille, maison & \zh{宀} (mián) \\
\hline
5 & \zh{会} & huì — savoir / réunion & \zh{人} / \zh{亻} (rén) \emph{(simplifié : 云 sous 人)} \\
\hline
6 & \zh{语} & yǔ — langue & \zh{讠} (yán) \\
\hline
7 & \zh{旁} & páng — côté & \zh{方} (fāng) \\
\hline
8 & \zh{清} & qīng — clair, pur & \zh{氵} (shuǐ) \\
\hline
9 & \zh{想} & xiǎng — penser & \zh{心} (xīn) / \zh{忄} \\
\hline
10 & \zh{国} & guó — pays & \zh{囗} (wéi) \\
\hline
\end{tabularx}
\end{center}
\textbf{Note pédagogique :} Pour \zh{写}, le radical officiel dans les dictionnaires récents est souvent \zh{冖} (mì, couvert) ou \zh{与}. Pour \zh{会}, le radical est \zh{人} (la partie supérieure). Vérifiez toujours l'index radical de votre dictionnaire de référence.
\end{corrige}

\coursec{Entraînement à la compréhension écrite}

Après avoir consolidé nos connaissances sur les \textbf{caractères : radicaux, traits et pièges visuels}, nous abordons maintenant une compétence cruciale pour l'examen : la \textbf{compréhension écrite}. En examen, vous n'aurez pas de dictionnaire. La clé est d'activer vos acquis (vocabulaire HSK 1-3, structures grammaticales, mots-outils) et d'appliquer une \textbf{méthode de lecture active} pour extraire l'information essentielle d'un texte authentique.

\courssub{Texte support d'entraînement : Une élève à Douala}

Commençons par la lecture d'un texte de niveau HSK 2-3, typique des sujets d'examen en classe de Première. Il met en scène \zh{王芳} (Wáng Fāng), une lycéenne de Douala. Lisez-le une première fois \textbf{sans vous arrêter} sur les mots inconnus, pour saisir l'idée générale.

\begin{textebox}[Texte d'examen — La journée de Wang Fang]
王芳是杜阿拉的中学生。她今年十六岁，学习汉语两年了。她每天早上六点起床，七点半去学校。中午她在学校食堂吃饭。下午三点放学，她常常和同学一起在图书馆学习汉字，复习课文。周末她喜欢听音乐，有时候和朋友上网聊天，看中国电影。她说汉语很有意思，虽然汉字有点儿难。以后她想去北京大学学习中文，也想去长城旅游。\\
Wáng Fāng shì Dù'ālā de zhōngxuéshēng. Tā jīnnián shíliù suì, xuéxí Hànyǔ liǎng nián le. Tā měitiān zǎoshang liù diǎn qǐchuáng, qī diǎn bàn qù xuéxiào. Zhōngwǔ tā zài xuéxiào shítáng chīfàn. Xiàwǔ sān diǎn fàngxué, tā chángcháng hé tóngxué yìqǐ zài túshūguǎn xuéxí Hànzì, fùxí kèwén. Zhōumò tā xǐhuan tīng yīnyuè, yǒu shíhou hé péngyou shàngwǎng liáotiān, kàn Zhōngguó diànyǐng. Tā shuō Hànyǔ hěn yǒu yìsi, suīrán Hànzì yǒudiǎnr nán. Yǐhòu tā xiǎng qù Běijīng Dàxué xuéxí Zhōngwén, yě xiǎng qù Chángchéng lǚyóu.\\
\emph{Wang Fang est lycéenne à Douala. Elle a seize ans cette année ; cela fait deux ans qu'elle apprend le chinois. Chaque jour, elle se lève à six heures du matin et va à l'école à sept heures et demie. Le midi, elle mange à la cantine de l'école. Les cours finissent à quinze heures ; elle étudie souvent les caractères avec ses camarades à la bibliothèque et révise les textes de leçon. Le week-end, elle aime écouter de la musique ; parfois elle bavarde en ligne avec ses amis et regarde des films chinois. Elle dit que le chinois est très intéressant, bien que les caractères soient un peu difficiles. Plus tard, elle voudrait aller étudier le chinois à l'Université de Pékin et aussi visiter la Grande Muraille.}
\end{textebox}

\begin{center}
\begin{tabularx}{\linewidth}{|X|X|X|X|}
\hline
\rowcolor{popblueL} \textbf{Caractères} & \textbf{Pinyin} & \textbf{Nature} & \textbf{Traduction} \\
\hline
杜阿拉 & Dù'ālā & n.p. & Douala (ville) \\
中学生 & zhōngxuéshēng & n. & collégien / lycéen \\
学习 & xuéxí & v. & étudier, apprendre \\
两年 & liǎng nián & num. + n. mesure & deux ans \\
每天 & měitiān & adv. de temps & chaque jour, tous les jours \\
起床 & qǐchuáng & v. & se lever \\
食堂 & shítáng & n. & cantine, réfectoire \\
放学 & fàngxué & v. & finir les cours, sortir de l'école \\
图书馆 & túshūguǎn & n. & bibliothèque \\
复习 & fùxí & v. & réviser \\
课文 & kèwén & n. & texte de leçon \\
周末 & zhōumò & n. & week-end \\
上网 & shàngwǎng & v. & aller sur Internet, surfer \\
聊天 & liáotiān & v. & bavarder, discuter \\
虽然 & suīrán & conj. & bien que, quoique \\
有点儿 & yǒudiǎnr & adv. & un peu (nuance souvent négative) \\
以后 & yǐhòu & n. de temps & plus tard, à l'avenir, après \\
北京大学 & Běijīng Dàxué & n.p. & Université de Pékin (Beida) \\
中文 & Zhōngwén & n. & chinois (langue écrite et parlée) \\
长城 & Chángchéng & n.p. & la Grande Muraille \\
旅游 & lǚyóu & v. & voyager, faire du tourisme \\
\hline
\end{tabularx}
\end{center}

\courssub{Méthode : Réussir l'épreuve de compréhension écrite}

La compréhension écrite ne teste pas seulement votre vocabulaire, mais votre \textbf{capacité à traiter l'information}. Voici la démarche en 4 étapes, conforme aux attendus des épreuves de classe de Première.

\begin{methode}[Stratégie de lecture active en 4 étapes]
\begin{enumerate}
\item \textbf{Lisez les questions AVANT le texte.} Cela oriente votre attention : cherchez-vous un lieu (\emph{Où ?}), une durée (\emph{Depuis quand ?}), une activité (\emph{Que fait-elle ?}) ou une opinion (\emph{Pourquoi ?}) ?
\item \textbf{Lecture rapide (survol) :} identifiez le \cle{sujet} (Qui ?), les \cle{marqueurs de temps} (Quand ? : \zh{每天}、\zh{周末}、\zh{以后}) et la \cle{structure globale} du texte (ici : routine quotidienne $\rightarrow$ week-end $\rightarrow$ projet futur).
\item \textbf{Lecture détaillée (repérage) :} soulignez les \textbf{mots connus} (vocabulaire HSK 1-3). Pour les mots inconnus, déduisez le sens par :
      \begin{itemize}
      \item le \textbf{contexte} (ex. : \zh{食堂} apparaît avec \zh{吃饭} « manger » $\rightarrow$ un lieu où l'on mange) ;
      \item les \textbf{radicaux} (ex. : \zh{图书馆} contient \zh{书} « livre » $\rightarrow$ un lieu avec des livres) ;
      \item la \textbf{composition du mot} (ex. : \zh{放学} = \zh{放} « libérer » + \zh{学} « études » $\rightarrow$ finir les cours).
      \end{itemize}
\item \textbf{Répondez en français} (sauf consigne contraire) par des phrases complètes, en citant le \textbf{segment du texte} qui justifie votre réponse.
\end{enumerate}
\end{methode}

\courssub{Analyse linguistique du texte : points de grammaire clés}

Le texte est une mine d'or pour réviser les structures de l'année. Observons trois points critiques.

\begin{attention}[Piège majeur : \zh{两年了} = durée qui continue jusqu'à maintenant]
En français, on dit « \emph{Elle apprend le chinois \textbf{depuis} deux ans} » (présent + \emph{depuis}).
En chinois, la structure est : \textbf{Sujet + Verbe + (Objet) + Durée + \zh{了}}.
\zh{学习汉语两年了} (xuéxí Hànyǔ liǎng nián le) ne signifie \textbf{pas} « elle a étudié deux ans (c'est fini) ».
Le \zh{了} final (particule de changement d'état) indique que l'action a commencé dans le passé \textbf{et continue encore au moment où l'on parle}.
\begin{itemize}
\item \zh{我住在这里三年了。} (Wǒ zhù zài zhèlǐ sān nián le.) — J'habite ici \emph{depuis} trois ans (et j'y suis encore).
\item \zh{他学中文一年了。} (Tā xué Zhōngwén yī nián le.) — Il apprend le chinois \emph{depuis} un an.
\end{itemize}
\emph{Erreur fréquente :} ne traduisez pas par le passé composé (« elle a étudié ») sauf si le contexte indique clairement que l'action est terminée.
\end{attention}

\begin{exemplebox}[Structure \zh{想} + Verbe : intention / désir futur]
\zh{以后她想去北京大学学习中文。}
(Yǐhòu tā xiǎng qù Běijīng Dàxué xuéxí Zhōngwén.)
\emph{Plus tard, elle voudrait aller étudier le chinois à l'Université de Pékin.}

\begin{center}
\begin{tabularx}{\linewidth}{|X|X|X|}
\hline
\rowcolor{popblueL} \textbf{Structure} & \textbf{Exemple (Caractères + Pinyin)} & \textbf{Traduction} \\
\hline
Sujet + \zh{想} + Verbe (déplacement) & \zh{我想去喀麦隆。} (Wǒ xiǎng qù Kāmàilóng.) & Je voudrais aller au Cameroun. \\
Sujet + \zh{想} + Verbe (action) & \zh{我们想吃喀麦隆菜。} (Wǒmen xiǎng chī Kāmàilóng cài.) & Nous avons envie de manger des plats camerounais. \\
Négation : Sujet + \zh{不想} + V & \zh{他不想上网。} (Tā bù xiǎng shàngwǎng.) & Il ne veut pas aller sur Internet. \\
\hline
\end{tabularx}
\end{center}
\emph{Note :} \zh{想} (xiǎng) exprime un \textbf{désir}, une intention ; il est plus nuancé que \zh{要} (yào, « vouloir fermement / devoir ») et plus concret que \zh{希望} (xīwàng, « espérer »).
\end{exemplebox}

\begin{propriete}[Marqueurs temporels et ordre des compléments]
Le texte respecte scrupuleusement la règle d'or de l'ordre des mots chinois : \textbf{Sujet + Temps + Lieu + Compagnon/Manière + Verbe + Complément}.
\begin{itemize}
\item \zh{她\textbf{每天早上六点}起床。} (Temps $\rightarrow$ Verbe)
\item \zh{她\textbf{常常}\textbf{和同学一起}\textbf{在图书馆}学习。} (Fréquence $\rightarrow$ Compagnon $\rightarrow$ Lieu $\rightarrow$ Verbe)
\item \zh{\textbf{周末}她喜欢听音乐。} (Temps placé en tête de phrase = thème)
\end{itemize}
\emph{Comparaison avec le français :} « \emph{Le week-end, elle aime écouter de la musique} » (mise en tête possible) ou « \emph{Elle aime écouter de la musique le week-end} ». En chinois, le complément de temps \textbf{ne se place jamais après le verbe} : il se met avant le verbe, après le sujet, ou en tête de phrase.
\end{propriete}

\courssub{Exercices type examen : compréhension et analyse}

Appliquons la méthode sur le texte de Wang Fang.

\begin{exoresolu}[Questions de compréhension directe]
\textbf{Énoncé :} Répondez aux questions en français en vous basant \textbf{uniquement} sur le texte.
\begin{enumerate}
\item Où habite Wang Fang ?
\item Depuis combien de temps apprend-elle le chinois ? Justifiez par la structure grammaticale.
\item Citez \textbf{trois} activités qu'elle fait le week-end.
\item Quel est son projet pour l'avenir (deux points) ?
\end{enumerate}

\tcblower
\textbf{Solution détaillée :}
\begin{enumerate}
\item \textbf{Elle habite à Douala.} (Texte : \zh{王芳是杜阿拉的中学生。} — « Wang Fang est lycéenne à Douala ».)
\item \textbf{Depuis deux ans.} Le texte dit : \zh{学习汉语两年了}. La structure \textbf{Verbe + Objet + Durée + \zh{了}} indique une action commencée dans le passé et qui se poursuit au présent : « cela fait deux ans qu'elle apprend le chinois ».
\item \textbf{1. Écouter de la musique} (\zh{听音乐}). \textbf{2. Bavarder en ligne avec des amis} (\zh{和朋友上网聊天}). \textbf{3. Regarder des films chinois} (\zh{看中国电影}).
\item \textbf{1. Étudier le chinois à l'Université de Pékin} (\zh{去北京大学学习中文}). \textbf{2. Visiter la Grande Muraille} (\zh{去长城旅游}).
\end{enumerate}
\end{exoresolu}

\begin{exoresolu}[Vrai / Faux + justification textuelle]
\textbf{Énoncé :} Les affirmations sont-elles vraies (V) ou fausses (F) ? Citez le segment exact du texte qui justifie votre réponse.
\begin{enumerate}
\item Wang Fang se lève à 7 h 30 le matin.
\item Elle mange le midi à la maison.
\item Elle étudie seule à la bibliothèque l'après-midi.
\item Elle trouve les caractères chinois faciles.
\end{enumerate}

\tcblower
\textbf{Solution détaillée :}
\begin{enumerate}
\item \textbf{Faux.} Texte : \zh{她每天早上六点起床，七点半去学校。} (Elle se lève à 6 h ; c'est à 7 h 30 qu'elle \emph{va à l'école}.)
\item \textbf{Faux.} Texte : \zh{中午她在学校食堂吃饭。} (Elle mange à la \textbf{cantine de l'école}, \zh{食堂}, pas à la maison.)
\item \textbf{Faux.} Texte : \zh{她常常和同学一起在图书馆学习。} (Elle étudie \textbf{avec ses camarades}, \zh{和同学一起}, pas seule.)
\item \textbf{Faux.} Texte : \zh{虽然汉字有点儿难。} (Les caractères sont \textbf{un peu difficiles}, \zh{有点儿难}, pas faciles.)
\end{enumerate}
\end{exoresolu}

\courssub{Recherche linguistique dans le texte : « chasse au trésor » grammaticale}

Un classique des épreuves : retrouver des éléments grammaticaux précis dans le corps du texte. Entraînez-vous à les repérer visuellement.

\begin{exercice}[Retrouver dans le texte]
Relisez le texte de Wang Fang et identifiez :
\begin{enumerate}
\item \textbf{Deux mots de mesure (量词)} utilisés avec des nombres. (Indice : l'un sert pour les années, l'autre pour les heures.)
\item \textbf{Une négation} : le texte n'en contient pas ; transformez la phrase « Elle aime écouter de la musique » à la forme négative en vous servant du vocabulaire du texte.
\item \textbf{Trois marqueurs temporels} (mots indiquant \emph{quand}).
\end{enumerate}
\end{exercice}

\begin{corrige}[Exercice : chasse au trésor grammaticale]
\begin{enumerate}
\item \textbf{Mots de mesure :}
      \begin{itemize}
      \item \zh{年} (nián) dans \zh{两年} (deux ans) — mot de mesure pour les années (on dit \zh{两年}, jamais \zh{二年} pour une quantité).
      \item \zh{点} (diǎn) dans \zh{六点}、\zh{七点半}、\zh{三点} (6 h, 7 h 30, 15 h) — mot de mesure pour les heures.
      \end{itemize}
\item \textbf{Négation (production) :} \zh{她不喜欢听音乐。} (Tā bù xǐhuan tīng yīnyuè.) — Elle n'aime pas écouter de la musique. (Règle : \zh{不} + verbe/adjectif pour la négation au présent ou pour une habitude.)
\item \textbf{Marqueurs temporels (au moins trois) :}
      \begin{itemize}
      \item \zh{每天} (měitiān) — chaque jour (fréquence) ;
      \item \zh{早上六点}、\zh{七点半}、\zh{中午}、\zh{下午三点} — moments précis de la journée ;
      \item \zh{周末} (zhōumò) — le week-end ;
      \item \zh{以后} (yǐhòu) — plus tard, à l'avenir.
      \end{itemize}
\end{enumerate}
\end{corrige}

\courssub{Visualisation : la journée type de Wang Fang}

Pour ancrer la chronologie et les marqueurs temporels, voici une frise synthétique. En examen, esquissez-en une rapidement sur votre brouillon pour structurer votre compréhension.

\begin{popfigure}
\begin{tikzpicture}[
    node distance=0.8cm and 1.2cm,
    event/.style={rectangle, rounded corners, draw=popblue, thick, fill=popblueL, text width=3.4cm, align=center, font=\small},
    time/.style={circle, draw=poporange, thick, fill=poporangeL, text width=1.6cm, align=center, font=\small\bfseries},
    arrow/.style={-Stealth, popblue, thick}
]
\node[time] (t1) {6h00};
\node[event, right=of t1, align=center] (e1){起床\\qǐchuáng\\(Se lever)};
\node[time, below=of t1] (t2) {7h30};
\node[event, right=of t2, align=center] (e2){去学校\\qù xuéxiào\\(Aller à l'école)};
\node[time, below=of t2] (t3) {12h00};
\node[event, right=of t3, align=center] (e3){在食堂吃饭\\zài shítáng chīfàn\\(Manger à la cantine)};
\node[time, below=of t3] (t4) {15h00};
\node[event, right=of t4, align=center] (e4){放学\\fàngxué\\(Fin des cours)};
\node[event, right=of e4, text width=3.6cm] (e5) {在图书馆复习课文、学汉字\\zài túshūguǎn fùxí kèwén, xué Hànzì\\(Réviser à la bibliothèque)};
\node[time, below=of t4, yshift=-0.5cm] (t5) {周末};
\node[event, right=of t5, text width=3.8cm] (e6) {听音乐、上网聊天\\tīng yīnyuè, shàngwǎng liáotiān\\(Musique, Internet)};
\node[event, below=of e6, text width=3.8cm, fill=popgreenL, draw=popgreen] (e7) {看中国电影\\kàn Zhōngguó diànyǐng\\(Films chinois)};

\draw[arrow] (t1) -- (e1);
\draw[arrow] (t2) -- (e2);
\draw[arrow] (t3) -- (e3);
\draw[arrow] (t4) -- (e4);
\draw[arrow] (e4) -- (e5);
\draw[arrow] (t5) -- (e6);
\draw[arrow] (e6) -- (e7);

\node[event, below=of e7, yshift=-0.8cm, fill=poppinkL, draw=poppink, text width=7.5cm] (future) {PROJET FUTUR (以后 yǐhòu) : 去北京大学学习中文 + 去长城旅游\\(Étudier le chinois à Beida + visiter la Grande Muraille)};
\draw[arrow, poppink] (e7) -- (future);
\end{tikzpicture}
\legende{Frise chronologique de la journée de Wang Fang : routine quotidienne (colonne de gauche) $\rightarrow$ week-end (en bas) $\rightarrow$ projet futur (encadré rose).}
\end{popfigure}

\courssub{Le saviez-vous ? Étudier en Chine depuis le Cameroun}

\begin{savaistu}[Bourses et Institut Confucius : un pont vers la Chine]
Chaque année, de nombreux lycéens et étudiants camerounais (de Yaoundé, Douala, Buea, Ngaoundéré...) partent étudier en Chine grâce à deux dispositifs principaux :
\begin{itemize}
\item \textbf{Les bourses du gouvernement chinois (CSC) :} elles couvrent en général les frais de scolarité, le logement et une allocation mensuelle. La sélection se fait sur dossier (résultats scolaires, niveau HSK, projet d'études), via l'ambassade de Chine au Cameroun et le MINESUP.
\item \textbf{L'Institut Confucius (Yaoundé, Douala) :} il organise les tests \textbf{HSK/HSKK} (indispensables pour les bourses), propose des cours de chinois et des séjours linguistiques d'été en Chine.
\end{itemize}
\textbf{Conseil :} visez le \textbf{HSK 3} en Première, puis le \textbf{HSK 4} en Terminale. Le projet de Wang Fang (\zh{去北京大学学习中文}) est tout à fait réaliste : les universités chinoises accueillent de plus en plus d'étudiants africains francophones. Maîtriser la compréhension écrite aujourd'hui, c'est préparer votre dossier de demain !
\end{savaistu}

\courssub{Entraînement final : production écrite guidée (thème)}

Pour clore cette section, passez de la compréhension à l'expression. Réutilisez \textbf{exactement} les structures du texte (marqueurs de temps, \zh{...两年了}, \zh{想} + verbe, \zh{虽然}...).

\begin{exercice}[Expression écrite : mon profil d'apprenant]
Rédigez un court paragraphe (60 à 80 caractères chinois) pour vous présenter sur le modèle de Wang Fang. Incluez obligatoirement :
\begin{enumerate}
\item votre ville (Yaoundé, Douala, Bafoussam...) ;
\item depuis combien de temps vous apprenez le chinois (structure \zh{...了}) ;
\item votre routine : heure de lever, école, révisions ;
\item une activité du week-end (\zh{周末}...) ;
\item un projet futur avec \zh{以后我想}...
\end{enumerate}
\emph{Aide :} utilisez le vocabulaire du tableau ci-dessus, ainsi que \zh{雅温得} (Yǎwēndé, Yaoundé), \zh{足球} (zúqiú, football), \zh{做饭} (zuòfàn, faire la cuisine).
\end{exercice}

\begin{corrige}[Exercice : mon profil d'apprenant — modèle de réponse]
\textbf{Modèle (élève à Yaoundé, un an de chinois, aime le football) :}

\zh{我是雅温得的中学生。我今年十七岁，学习汉语一年了。我每天早上六点半起床，七点二十去学校。中午我在学校食堂吃饭。下午放学后，我在家复习汉字。周末我喜欢踢足球，也看中国电影。虽然汉字有点儿难，但是汉语很有意思。以后我想去中国学习中文，也想去长城旅游。}

(Wǒ shì Yǎwēndé de zhōngxuéshēng. Wǒ jīnnián shíqī suì, xuéxí Hànyǔ yī nián le. Wǒ měitiān zǎoshang liù diǎn bàn qǐchuáng, qī diǎn èrshí qù xuéxiào. Zhōngwǔ wǒ zài xuéxiào shítáng chīfàn. Xiàwǔ fàngxué hòu, wǒ zài jiā fùxí Hànzì. Zhōumò wǒ xǐhuan tī zúqiú, yě kàn Zhōngguó diànyǐng. Suīrán Hànzì yǒudiǎnr nán, dànshì Hànyǔ hěn yǒu yìsi. Yǐhòu wǒ xiǎng qù Zhōngguó xuéxí Zhōngwén, yě xiǎng qù Chángchéng lǚyóu.)

\emph{« Je suis lycéen(ne) à Yaoundé. J'ai dix-sept ans ; cela fait un an que j'apprends le chinois. Chaque jour, je me lève à six heures et demie et je vais à l'école à sept heures vingt. Le midi, je mange à la cantine de l'école. L'après-midi, après les cours, je révise les caractères à la maison. Le week-end, j'aime jouer au football et je regarde aussi des films chinois. Bien que les caractères soient un peu difficiles, le chinois est très intéressant. Plus tard, je voudrais aller étudier le chinois en Chine et visiter la Grande Muraille. »}

\textbf{Analyse du modèle :}
\begin{itemize}
\item \zh{学习汉语一年了} : structure de durée continue correcte (Verbe + Objet + Durée + \zh{了}).
\item \zh{六点半}、\zh{七点二十} : heures précises (\zh{点} + \zh{半} « et demie » ; \zh{点} + minutes).
\item \zh{虽然...但是...} (suīrán... dànshì...) : structure de concession de l'année ; en chinois, on peut utiliser les deux termes ensemble, contrairement au français (« bien que... mais » est fautif en français, correct en chinois).
\item \zh{以后我想去...学习} : projet futur clair avec \zh{想} + verbe.
\end{itemize}
\end{corrige}

\coursec{Traduction et expression écrite guidées}

Dans la section précédente, nous avons entraîné notre œil à \emph{comprendre} un texte chinois : repérer le sujet, suivre l'ordre des mots, deviner le sens des mots inconnus grâce au contexte. Nous allons maintenant franchir une étape supplémentaire : passer d'une langue à l'autre de façon active, en \emph{traduisant} dans les deux sens (thème et version), puis \emph{produire} notre propre texte. C'est exactement ce que l'on vous demandera à l'examen : cette section est donc un entraînement en conditions réelles.

\courssub{Le texte de référence : la lettre de Wang Fang}

Pour travailler la version et la production écrite, nous nous appuyons sur un texte modèle. Lisez-le d'abord en silence, puis à voix haute en suivant le pinyin.

\begin{textebox}[Une lettre de Wang Fang]
你好！我叫王芳，我是中国人。我住在上海。\\
Nǐ hǎo ! Wǒ jiào Wáng Fāng, wǒ shì Zhōngguó rén. Wǒ zhù zài Shànghǎi.\\
« Bonjour ! Je m'appelle Wang Fang, je suis chinoise. J'habite à Shanghai. »\\
我家有四口人：爸爸、妈妈、哥哥和我。我爸爸是老师，我妈妈在医院工作。\\
Wǒ jiā yǒu sì kǒu rén ： bàba, māma, gēge hé wǒ. Wǒ bàba shì lǎoshī, wǒ māma zài yīyuàn gōngzuò.\\
« Ma famille compte quatre personnes : papa, maman, mon grand frère et moi. Mon père est professeur, ma mère travaille à l'hôpital. »\\
我每天早上七点起床，八点去学校。我学习汉语、英语和数学。\\
Wǒ měitiān zǎoshang qī diǎn qǐchuáng, bā diǎn qù xuéxiào. Wǒ xuéxí Hànyǔ, Yīngyǔ hé shùxué.\\
« Chaque matin, je me lève à sept heures et je vais à l'école à huit heures. J'étudie le chinois, l'anglais et les mathématiques. »\\
我也喜欢运动。我常常和朋友们一起踢足球。\\
Wǒ yě xǐhuan yùndòng. Wǒ chángcháng hé péngyoumen yìqǐ tī zúqiú.\\
« J'aime aussi le sport. Je joue souvent au football avec mes amis. »\\
因为汉语很有意思，所以我想去中国……不对，我已经是中国人了！我想去喀麦隆，因为喀麦隆的足球很有名！\\
Yīnwèi Hànyǔ hěn yǒuyìsi, suǒyǐ wǒ xiǎng qù Zhōngguó…… bú duì, wǒ yǐjīng shì Zhōngguó rén le ! Wǒ xiǎng qù Kāmàilóng, yīnwèi Kāmàilóng de zúqiú hěn yǒumíng !\\
« Parce que le chinois est très intéressant, je voudrais aller en Chine... non, je suis déjà chinoise ! Je voudrais aller au Cameroun, parce que le football camerounais est très célèbre ! »\\
以后我想当翻译。你呢？\\
Yǐhòu wǒ xiǎng dāng fānyì. Nǐ ne ?\\
« Plus tard, je voudrais devenir traductrice. Et toi ? »
\end{textebox}

\begin{center}
\begin{tabularx}{\linewidth}{|X|X|X|X|}
\hline
\rowcolor{popblueL} Caractères & Pinyin & Nature & Traduction \\
\hline
住 & zhù & verbe & habiter, vivre \\
\hline
口 & kǒu & classificateur & (bouche) personnes d'une famille \\
\hline
医院 & yīyuàn & nom & hôpital \\
\hline
起床 & qǐchuáng & verbe & se lever \\
\hline
一起 & yìqǐ & adverbe & ensemble \\
\hline
踢足球 & tī zúqiú & locution verbale & jouer au football \\
\hline
有名 & yǒumíng & adjectif & célèbre \\
\hline
当 & dāng & verbe & devenir, exercer comme \\
\hline
翻译 & fānyì & nom / verbe & traducteur ; traduire \\
\hline
已经 & yǐjīng & adverbe & déjà \\
\hline
\end{tabularx}
\end{center}

\courssub{La méthode de traduction en cinq étapes}

Traduire, ce n'est pas remplacer mot à mot : c'est \emph{reconstruire} la phrase selon les lois de l'autre langue. Voici la méthode officielle à appliquer à chaque phrase, au brouillon, le jour de l'examen.

\begin{methode}[Traduire une phrase en cinq étapes]
\begin{enumerate}
\item \textbf{Lire} la phrase entière deux fois, sans écrire. Identifier le sens global.
\item \textbf{Repérer} le sujet (S), le verbe (V) et le complément (C). Qui fait quoi ?
\item \textbf{Ordonner} selon le schéma chinois : Sujet + Temps + Lieu + Verbe + Complément. En chinois, le temps et le lieu se placent \emph{avant} le verbe, jamais à la fin comme en français.
\item \textbf{Traduire} mot à mot à l'intérieur de ce cadre, en choisissant le bon classificateur s'il y a un nom compté.
\item \textbf{Vérifier} : ordre des mots, présence de 是 devant un nom, 了 pour une action accomplie, tons du pinyin, traits des caractères.
\end{enumerate}
\end{methode}

\begin{popfigure}
\begin{tikzpicture}[
node distance=0.55cm,
etape/.style={rectangle, rounded corners=4pt, draw=popblue, fill=popblueL, text width=3.1cm, align=center, font=\small, minimum height=0.9cm},
fleche/.style={-{Stealth[length=2.5mm]}, thick, popdark}]
\node[etape, align=center] (lire){\textbf{1. Lire}\\ comprendre le sens global};
\node[etape, right=of lire, align=center] (reperer){\textbf{2. Repérer}\\ Sujet / Verbe / Complément};
\node[etape, right=of reperer, align=center] (ordonner){\textbf{3. Ordonner}\\ S + Temps + Lieu + V + C};
\node[etape, below=0.7cm of ordonner, align=center] (traduire){\textbf{4. Traduire}\\ mots et classificateurs};
\node[etape, left=of traduire, align=center] (verifier){\textbf{5. Vérifier}\\ ordre, 了, tons, caractères};
\draw[fleche] (lire) -- (reperer);
\draw[fleche] (reperer) -- (ordonner);
\draw[fleche] (ordonner) -- (traduire);
\draw[fleche] (traduire) -- (verifier);
\end{tikzpicture}
\legende{Organigramme de la méthode de traduction : de la lecture à la vérification.}
\end{popfigure}

\begin{exemplebox}[Application de la méthode, étape par étape]
Phrase à traduire : « Hier, j'ai joué au football avec mes amis à l'école. »
\begin{itemize}
\item Étape 2 : S = je ; V = jouer ; C = au football ; temps = hier ; lieu = à l'école ; avec = mes amis.
\item Étape 3 : Je + hier + à l'école + avec mes amis + jouer + football.
\item Étape 4 : \zh{我昨天在学校和朋友们踢了足球。}
\end{itemize}
\emph{Wǒ zuótiān zài xuéxiào hé péngyoumen tī le zúqiú.} « Hier, j'ai joué au football avec mes amis à l'école. »
\end{exemplebox}

\begin{attention}[Le piège de l'ordre des mots]
Le francophone dit : « Je joue au football \emph{à l'école} \emph{hier} » — le lieu et le temps à la fin. En chinois, c'est impossible : le temps et le lieu précèdent le verbe. On ne dit jamais \zh{我踢足球在学校昨天}. La bonne phrase : \zh{我昨天在学校踢足球。} \emph{Wǒ zuótiān zài xuéxiào tī zúqiú.} Retenez le réflexe : \textbf{d'abord le cadre (quand ? où ?), ensuite l'action}.
\end{attention}

\courssub{Thème guidé : du français vers le chinois}

Voici cinq phrases progressives qui mobilisent les verbes et mots-outils essentiels de l'année : 是 (être), 有 (avoir), 在 (être à), 了 (accompli) et un mot interrogatif. Traduisez-les d'abord seul, puis comparez avec le corrigé.

\begin{exoresolu}[Thème : cinq phrases progressives]
Énoncé. Traduisez en chinois (caractères et pinyin) :
\begin{enumerate}
\item Je suis camerounais.
\item J'ai deux frères.
\item Mon père travaille à Douala.
\item Hier, je suis allé au marché.
\item Où habites-tu ?
\end{enumerate}
\tcblower
Solution détaillée :
\begin{enumerate}
\item \zh{我是喀麦隆人。} \emph{Wǒ shì Kāmàilóng rén.} — Verbe d'état 是 devant un nom ; pas d'article en chinois.
\item \zh{我有两个哥哥。} \emph{Wǒ yǒu liǎng gè gēge.} — « Deux » devant classificateur = 两 (pas 二) ; classificateur général 个. \emph{Note :} \zh{哥哥} = frères aînés ; pour des cadets : \zh{两个弟弟} (liǎng gè dìdi).
\item \zh{我爸爸在杜阿拉工作。} \emph{Wǒ bàba zài Dù'ālā gōngzuò.} — Le lieu (在 + lieu) se place avant le verbe 工作.
\item \zh{昨天我去了市场。} \emph{Zuótiān wǒ qù le shìchǎng.} — Le temps 昨天 ouvre la phrase ; 了 marque l'action accomplie.
\item \zh{你住在哪儿？} \emph{Nǐ zhù zài nǎr ?} — Le mot interrogatif 哪儿 reste à la place de la réponse ; on ne le déplace pas comme en français.
\end{enumerate}
\end{exoresolu}

\courssub{Version guidée : du chinois vers le français}

Pour la version, on extrait des phrases du texte de Wang Fang. La difficulté est inverse : il faut restituer en français un ordre des mots naturel (souvent temps et lieu à la fin) et ne pas calquer le chinois.

\begin{exoresolu}[Version : cinq phrases du texte de Wang Fang]
Énoncé. Traduisez en français :
\begin{enumerate}
\item \zh{我住在上海。}
\item \zh{我家有四口人。}
\item \zh{我每天早上七点起床。}
\item \zh{我学习汉语、英语和数学。}
\item \zh{以后我想当翻译。}
\end{enumerate}
\tcblower
Solution détaillée :
\begin{enumerate}
\item \emph{Wǒ zhù zài Shànghǎi.} = « J'habite à Shanghai. » — 在 + ville après 住.
\item \emph{Wǒ jiā yǒu sì kǒu rén.} = « Ma famille compte quatre personnes. » — 口 est le classificateur des membres de la famille ; on ne traduit pas « bouches » !
\item \emph{Wǒ měitiān zǎoshang qī diǎn qǐchuáng.} = « Je me lève à sept heures tous les matins. » — En français, le temps passe à la fin.
\item \emph{Wǒ xuéxí Hànyǔ, Yīngyǔ hé shùxué.} = « J'étudie le chinois, l'anglais et les mathématiques. » — 和 ne relie que les deux derniers éléments, comme « et ».
\item \emph{Yǐhòu wǒ xiǎng dāng fānyì.} = « Plus tard, je voudrais devenir traductrice. » — 以后 se traduit ici « plus tard » en tête de phrase.
\end{enumerate}
\end{exoresolu}

\courssub{Les connecteurs : le squelette de votre texte}

Un bon texte d'examen n'est pas une suite de phrases isolées : c'est un enchaînement logique. Voici les cinq connecteurs à réutiliser absolument.

\begin{center}
\begin{tabularx}{\linewidth}{|X|X|X|X|}
\hline
\rowcolor{popblueL} Connecteur & Pinyin & Sens & Exemple (traduction) \\
\hline
和 & hé & et (entre noms) & \zh{我喜欢汉语和英语。} J'aime le chinois et l'anglais. \\
\hline
也 & yě & aussi (avant le verbe) & \zh{我也喜欢音乐。} Moi aussi j'aime la musique. \\
\hline
但是 & dànshì & mais & \zh{汉语很难，但是很有意思。} Le chinois est difficile, mais intéressant. \\
\hline
因为...所以... & yīnwèi... suǒyǐ... & parce que... donc... & \zh{因为汉语很有意思，所以我想去中国。} \\
\hline
以后 & yǐhòu & ensuite, plus tard & \zh{以后我想当翻译。} Plus tard, je veux devenir traducteur. \\
\hline
\end{tabularx}
\end{center}

\begin{propriete}[La place de 也 (aussi)]
L'adverbe 也 se place \textbf{avant le verbe} (et après le sujet), jamais en fin de phrase comme « aussi » en français.\\
Schéma : Sujet + 也 + Verbe (+ Complément).\\
\zh{我也喜欢音乐。} \emph{Wǒ yě xǐhuan yīnyuè.} « Moi aussi j'aime la musique. »\\
\zh{他也是学生。} \emph{Tā yě shì xuésheng.} « Lui aussi est élève. »
\end{propriete}

\begin{exemplebox}[La paire 因为...所以... en action]
\zh{因为汉语很有意思，所以我想去中国。}\\
\emph{Yīnwèi Hànyǔ hěn yǒuyìsi, suǒyǐ wǒ xiǎng qù Zhōngguó.}\\
« Parce que le chinois est très intéressant, je voudrais aller en Chine. »\\
Autre exemple : \zh{因为今天下雨，所以我不去学校。} \emph{Yīnwèi jīntiān xiàyǔ, suǒyǐ wǒ bú qù xuéxiào.} « Parce qu'il pleut aujourd'hui, je ne vais pas à l'école. »
\end{exemplebox}

\begin{attention}[Ne supprimez pas 所以 !]
En français, « parce que... donc... » ensemble est une faute de style ; on choisit l'un ou l'autre. En chinois, c'est le contraire : 因为 et 所以 apparaissent naturellement \emph{ensemble} dans la même phrase. Par réflexe francophone, beaucoup d'élèves effacent 所以 : erreur ! Gardez les deux. À l'oral comme à l'écrit, \zh{因为...，所以...} est la structure attendue et correcte.
\end{attention}

\courssub{Production écrite guidée : « Ma journée et mes loisirs »}

\begin{methode}[Rédiger en cinq étapes avec le plan fourni]
\begin{enumerate}
\item \textbf{Présentation} : nom, nationalité, ville. (\zh{我叫……，我是……人，我住在……。})
\item \textbf{Famille} : nombre de personnes, métiers. (\zh{我家有……口人。})
\item \textbf{Journée} : heures, école, matières. (\zh{我每天……点起床，……点去学校。})
\item \textbf{Loisirs} : sport, musique, amis — utilisez 也, 和, 但是. (\zh{我也喜欢……。})
\item \textbf{Projet} : utilisez 因为...所以... et 以后. (\zh{因为……，所以以后我想……。})
\end{enumerate}
Règle d'or : \textbf{une idée par phrase}, et réutilisez les structures apprises plutôt que d'improviser des phrases compliquées.
\end{methode}

\begin{exemplebox}[Texte modèle (7 phrases)]
\zh{我叫恩东，我是喀麦隆人，我住在雅温得。我家有五口人：爸爸、妈妈、两个姐姐和我。我每天早上六点起床，七点去学校。我学习汉语、法语和数学。我也喜欢运动，我常常和朋友们一起踢足球。汉语很难，但是很有意思。因为汉语很有意思，所以以后我想去中国学习。}\\
\emph{Wǒ jiào Ēndōng, wǒ shì Kāmàilóng rén, wǒ zhù zài Yǎwēndé. Wǒ jiā yǒu wǔ kǒu rén : bàba, māma, liǎng gè jiějie hé wǒ. Wǒ měitiān zǎoshang liù diǎn qǐchuáng, qī diǎn qù xuéxiào. Wǒ xuéxí Hànyǔ, Fǎyǔ hé shùxué. Wǒ yě xǐhuan yùndòng, wǒ chángcháng hé péngyoumen yìqǐ tī zúqiú. Hànyǔ hěn nán, dànshì hěn yǒuyìsi. Yīnwèi Hànyǔ hěn yǒuyìsi, suǒyǐ yǐhòu wǒ xiǎng qù Zhōngguó xuéxí.}\\
« Je m'appelle Ndong, je suis camerounais, j'habite à Yaoundé. Ma famille compte cinq personnes : papa, maman, mes deux grandes sœurs et moi. Chaque matin je me lève à six heures et je vais à l'école à sept heures. J'étudie le chinois, le français et les mathématiques. J'aime aussi le sport, je joue souvent au football avec mes amis. Le chinois est difficile, mais très intéressant. Parce que le chinois est très intéressant, plus tard je voudrais aller étudier en Chine. »
\end{exemplebox}

\begin{exercice}[À vous d'écrire]
Rédigez un texte de 5 à 8 phrases intitulé « Ma journée et mes loisirs » en suivant le plan en cinq étapes. Contraintes obligatoires : employer au moins une fois 也, 但是, 因为...所以... et 以后 ; utiliser 了 pour au moins une action passée ; écrire en caractères, avec le pinyin au-dessus ou en dessous.
\end{exercice}

\begin{corrige}[Exercice « À vous d'écrire » : éléments de correction]
Il n'existe pas une seule bonne réponse, mais tout bon texte doit contenir : une présentation avec 是 ; une phrase de famille avec 有 + 口 ; des horaires avec 点 et le temps placé avant le verbe ; un loisir avec 也 ou 和 ; une phrase de projet avec 因为...所以... et 以后. Vérifiez chaque point de la checklist ci-dessous avant de rendre votre copie. Exemple de faute fréquente à corriger : \zh{我喜欢也音乐。} est faux — corrigez en \zh{我也喜欢音乐。}
\end{corrige}

\courssub{Relecture : la checklist de l'examen}

Avant de rendre votre copie, relisez chaque phrase en cochant mentalement ces cinq points :

\begin{aretenir}[Checklist de relecture]
\begin{enumerate}
\item \textbf{Ordre des mots} : Sujet + Temps + Lieu + Verbe + Complément. Rien après le verbe sauf le complément d'objet.
\item \textbf{是 ou 在 ?} : 是 devant un nom ou une identité (\zh{我是学生。}) ; 在 devant un lieu (\zh{我在学校。}).
\item \textbf{不 ou 没 ?} : 不 pour le présent et les habitudes (\zh{我不去。}) ; 没 pour nier une action passée (\zh{我昨天没去。}) — et jamais 了 avec 没.
\item \textbf{Tons du pinyin} : chaque syllabe porte sa marque ; vérifiez les paires pièges (mā/mà, shì/shí, yǒu/yòu).
\item \textbf{Traits des caractères} : horizontal avant vertical, gauche avant droite, haut avant bas ; attention aux caractères proches (人/入, 大/太, 未/末).
\end{enumerate}
\end{aretenir}

\courssub{Grille d'évaluation type examen}

Voici la grille utilisée pour noter votre production écrite, sur 20 points. Étudiez-la : elle vous dit exactement où gagner des points.

\begin{center}
\begin{tabularx}{\linewidth}{|X|X|X|}
\hline
\rowcolor{popblueL} Critère & Points & Ce que l'examinateur regarde \\
\hline
Exactitude grammaticale & /4 & ordre des mots, 是/在, 不/没, place de 了, de 也 \\
\hline
Vocabulaire & /4 & mots justes et variés, bons classificateurs (个, 口) \\
\hline
Caractères & /4 & traits corrects, caractères reconnaissables, pas de confusion \\
\hline
Cohérence & /4 & connecteurs (和, 也, 但是, 因为...所以..., 以后), enchaînement logique \\
\hline
Respect de la consigne & /4 & 5 à 8 phrases, plan suivi, thème respecté \\
\hline
\end{tabularx}
\end{center}

\begin{savaistu}[La clarté d'abord !]
À l'examen, une phrase simple et correcte rapporte plus de points qu'une phrase longue et fautive. Les correcteurs de chinois — au Cameroun comme en Chine — valorisent d'abord la justesse de la structure. Cinq phrases courtes et parfaites avec 是, 有, 在, 也 et 因为...所以... valent mieux qu'une phrase ambitieuse où tout s'effondre. La stratégie du bon élève : écrire ce que l'on sait, relire, et ne jamais improviser une structure jamais vue en classe.
\end{savaistu}

\coursec{Simulation d'examen et stratégies}

Dans la section précédente, nous avons révisé la traduction et l'expression écrite guidées : vous savez maintenant traduire une phrase en respectant l'ordre des mots chinois et rédiger un court texte cohérent. Il reste une dernière étape, la plus importante : mettre tous ces acquis en situation réelle d'examen. Cette section vous propose une épreuve blanche complète, des stratégies de gestion du temps et du stress, puis une méthode d'auto-correction pour repérer vos dernières faiblesses avant le jour J.

\courssub{Observer d'abord : à quoi ressemble une bonne copie ?}

Avant de parler stratégie, observons deux productions d'élèves à la même question : « Présente-toi en trois phrases. »

\begin{exemplebox}[Deux copies comparées]
\textbf{Copie A (faible) :} \zh{我学生。我喜欢。汉语有意思。}\\
\emph{Wǒ xuésheng. Wǒ xǐhuan. Hànyǔ yǒu yìsi.}\\
« Je étudiant. J'aime. Le chinois intéressant. » — phrases incomplètes : le verbe \zh{是} manque, l'objet de \zh{喜欢} manque, l'adverbe \zh{很} manque devant l'adjectif.

\textbf{Copie B (réussie) :} \zh{我叫玛丽，我是雅温得的学生。我喜欢汉语，因为汉语很有意思。}\\
\emph{Wǒ jiào Mǎlì, wǒ shì Yǎwēndé de xuésheng. Wǒ xǐhuan Hànyǔ, yīnwèi Hànyǔ hěn yǒu yìsi.}\\
« Je m'appelle Marie, je suis élève à Yaoundé. J'aime le chinois parce que le chinois est très intéressant. » — verbes \zh{叫} et \zh{是} présents, phrase avec \zh{因为} correcte, adjectif précédé de \zh{很}.
\end{exemplebox}

Que remarque-t-on ? La copie B n'utilise pas de vocabulaire difficile : elle utilise \emph{correctement} des structures simples de l'année. C'est exactement ce que l'examen récompense. La stratégie gagnante n'est donc pas d'impressionner, mais de \cle{mobiliser sûrement} ce que l'on sait.

\begin{attention}[Le piège de l'adjectif « tout seul »]
En français, on dit « le chinois est intéressant » avec le verbe « être ». En chinois, un adjectif seul en fin de phrase sonne incomplet ou comparatif : \zh{汉语有意思。} sous-entend plutôt « le chinois, lui, est intéressant (par opposition à autre chose) ». Pour une simple description, placez \zh{很} devant l'adjectif : \zh{汉语很有意思。} (\emph{Hànyǔ hěn yǒu yìsi.}) « Le chinois est très intéressant. » Ici, \zh{很} ne marque presque plus l'intensité : il rend la phrase naturelle.
\end{attention}

\courssub{Le cercle de la stratégie d'épreuve}

Toute épreuve réussie suit le même cycle en cinq étapes. Mémorisez-le comme un réflexe.

\begin{popfigure}
\begin{tikzpicture}[scale=1]
\node[circle, draw=popblue, fill=popblueL, minimum size=2.1cm, align=center, font=\small] (a) at (90:2.6) {1. Lire les\\consignes};
\node[circle, draw=popgreen, fill=popgreenL, minimum size=2.1cm, align=center, font=\small] (b) at (18:2.6) {2. Comprendre\\le texte};
\node[circle, draw=poporange, fill=poporangeL, minimum size=2.1cm, align=center, font=\small] (c) at (-54:2.6) {3. Rédiger};
\node[circle, draw=poppurple, fill=poppurpleL, minimum size=2.1cm, align=center, font=\small] (d) at (-126:2.6) {4. Vérifier\\(ordre, tons)};
\node[circle, draw=poppink, fill=poppinkL, minimum size=2.1cm, align=center, font=\small] (e) at (162:2.6) {5. Relire\\la copie};
\draw[-{Stealth}, thick, popdark] (a) to[bend left=18] (b);
\draw[-{Stealth}, thick, popdark] (b) to[bend left=18] (c);
\draw[-{Stealth}, thick, popdark] (c) to[bend left=18] (d);
\draw[-{Stealth}, thick, popdark] (d) to[bend left=18] (e);
\draw[-{Stealth}, thick, popdark] (e) to[bend left=18] (a);
\node[align=center, font=\small\itshape, popmuted] at (0,0) {Stratégie\\d'épreuve};
\end{tikzpicture}
\legende{Le cycle en cinq étapes : lire, comprendre, rédiger, vérifier, relire.}
\end{popfigure}

\begin{methode}[Gérer le stress pendant l'épreuve]
\begin{enumerate}
\item \textbf{Commencez par les questions faciles.} Répondre vite à ce que vous maîtrisez vous met en confiance et sécurise des points.
\item \textbf{Laissez un blanc plutôt qu'une phrase inventée.} Une phrase au hasard avec des fautes coûte des points ; un blanc n'en coûte qu'un. Revenez-y ensuite.
\item \textbf{Revenez aux questions difficiles à la fin.} Souvent, un autre exercice vous aura rappelé le mot ou la structure oubliée.
\item \textbf{Respirez.} Si vous bloquez, fermez les yeux cinq secondes, puis relisez la consigne mot à mot : la réponse est souvent dans la question.
\end{enumerate}
\end{methode}

\courssub{Gestion du temps : la répartition 10/8/6/6 + 3}

Pour une épreuve blanche de 30 minutes, voici le plan de bataille recommandé :

\begin{center}
\begin{tabularx}{\linewidth}{|X|X|X|}\hline
\rowcolor{popblueL} \textbf{Partie} & \textbf{Temps} & \textbf{Objectif} \\ \hline
Compréhension de texte (5 questions) & 10 min & lire 2 fois le texte, répondre en phrases simples \\ \hline
Grammaire (5 QCM) & 8 min & éliminer d'abord les réponses fautives \\ \hline
Traduction (4 phrases) & 6 min & ordre des mots avant tout \\ \hline
Expression écrite (5 phrases min.) & 6 min & structures sûres uniquement \\ \hline
Relecture finale & 3 min & caractères lisibles, négations, verbes \\ \hline
\end{tabularx}
\end{center}

\begin{attention}[L'écriture des caractères compte !]
À l'examen, un caractère \cle{illisible} peut être compté faux, même si vous connaissiez la bonne réponse. Écrivez des traits \textbf{nets et espacés}, dans l'ordre appris (haut → bas, gauche → droite, horizontal avant vertical). Ne « dessinez » pas les caractères à la va-vite : \zh{人} (rén, personne) mal fermé devient \zh{入} (rù, entrer), et le sens change complètement. Mieux vaut écrire moins de caractères, mais correctement.
\end{attention}

\courssub{Stratégie QCM : éliminer avant de choisir}

Face à un QCM de grammaire, ne cherchez pas d'abord la bonne réponse : \cle{éliminez} les réponses impossibles. Deux fautes tuent presque toujours une option :

\begin{itemize}
\item une \textbf{faute d'ordre des mots} (ex. complément de temps placé après le verbe, comme en français) ;
\item une \textbf{faute de négation} (ex. \zh{不有} au lieu de \zh{没有} pour nier le verbe \zh{有}).
\end{itemize}

\begin{exemplebox}[QCM type corrigé]
\textbf{Question :} 我\_\_\_弟弟。 (\emph{Wǒ \_\_\_ dìdi.})\\
(a) \zh{不有} \quad (b) \zh{没有} \quad (c) \zh{不是}

\textbf{Analyse :} (a) est impossible : \zh{不} ne nie jamais \zh{有}. (c) est impossible : \zh{是} signifie « être », pas « avoir » ; la phrase voudrait dire « je ne suis pas un petit frère ». Reste (b).\\
\textbf{Réponse : (b)} \zh{我没有弟弟。} (\emph{Wǒ méiyǒu dìdi.}) « Je n'ai pas de petit frère. »
\end{exemplebox}

\begin{propriete}[Rappel : les deux négations de l'année]
\begin{itemize}
\item \zh{不} + verbe/adjectif : négation présente ou future, ou habitude. \zh{我不喝茶。} (\emph{Wǒ bù hē chá.}) « Je ne bois pas de thé. »
\item \zh{没(有)} + verbe : négation du passé ou de la possession. \zh{我昨天没去学校。} (\emph{Wǒ zuótiān méi qù xuéxiào.}) « Hier, je ne suis pas allé à l'école. »
\item \textbf{Exception absolue :} \zh{有} se nie toujours par \zh{没有}, jamais par \zh{不有}.
\end{itemize}
\end{propriete}

\begin{propriete}[Rappel express : 的, 得, 地 — trois mots-outils à ne plus confondre]
\begin{center}
\begin{tabularx}{\linewidth}{|X|X|X|}\hline
\rowcolor{popblueL} Structure & Exemple en caractères + pinyin & Traduction \\ \hline
nom/pronom + 的 + nom & \zh{我的书} \emph{wǒ de shū} & mon livre \\ \hline
verbe + 得 + adjectif & \zh{他说得很好} \emph{tā shuō de hěn hǎo} & il parle très bien \\ \hline
adjectif + 地 + verbe & \zh{他高兴地说} \emph{tā gāoxìng de shuō} & il dit joyeusement \\ \hline
\end{tabularx}
\end{center}
À l'examen, repérez d'abord ce qui suit le blanc : un nom appelle \zh{的}, un adjectif après un verbe appelle \zh{得}, un verbe après une manière appelle \zh{地}.
\end{propriete}

\courssub{Épreuve blanche chronométrée (30 minutes)}

\begin{textebox}[Texte de compréhension : 我的周末 (Mon week-end)]
我叫李雷，我是杜阿拉的学生。\\
\emph{Wǒ jiào Lǐ Léi, wǒ shì Dù'ālā de xuésheng.}\\
« Je m'appelle Li Lei, je suis élève à Douala. »\\
星期六早上，我和妈妈去市场买东西。\\
\emph{Xīngqīliù zǎoshang, wǒ hé māma qù shìchǎng mǎi dōngxi.}\\
« Samedi matin, je vais au marché avec maman faire des courses. »\\
下午，我和朋友们踢足球。我踢得不太好，但是我很高兴。\\
\emph{Xiàwǔ, wǒ hé péngyoumen tī zúqiú. Wǒ tī de bú tài hǎo, dànshì wǒ hěn gāoxìng.}\\
« L'après-midi, je joue au football avec mes amis. Je ne joue pas très bien, mais je suis content. »\\
星期天，我在家学习汉语。汉语很有意思，我也喜欢看中国的电影。\\
\emph{Xīngqītiān, wǒ zài jiā xuéxí Hànyǔ. Hànyǔ hěn yǒu yìsi, wǒ yě xǐhuan kàn Zhōngguó de diànyǐng.}\\
« Dimanche, j'étudie le chinois à la maison. Le chinois est très intéressant, et j'aime aussi regarder des films chinois. »\\
晚上，我们全家一起吃饭。\\
\emph{Wǎnshang, wǒmen quán jiā yìqǐ chīfàn.}\\
« Le soir, toute la famille mange ensemble. »
\end{textebox}

\begin{center}
\begin{tabularx}{\linewidth}{|X|X|X|X|}\hline
\rowcolor{popblueL} Caractères & Pinyin & Nature & Traduction \\ \hline
周末 & zhōumò & nom & week-end \\ \hline
市场 & shìchǎng & nom & marché \\ \hline
买东西 & mǎi dōngxi & expr. verb. & faire des courses \\ \hline
踢足球 & tī zúqiú & expr. verb. & jouer au football \\ \hline
全家 & quán jiā & nom & toute la famille \\ \hline
\end{tabularx}
\end{center}

\begin{exercice}[Épreuve blanche — 30 minutes]
\textbf{I. Compréhension de texte (10 min, 5 points).} Répondez en français ou en chinois simple.\\
1. Où habite Li Lei ? 2. Avec qui va-t-il au marché ? 3. Que fait-il samedi après-midi ? 4. Joue-t-il bien au football ? 5. Pourquoi aime-t-il le chinois ?

\textbf{II. Grammaire — QCM (8 min, 5 points).} Entourez la bonne réponse.\\
1. 我昨天\_\_\_去学校。(a) 不 (b) 没 (c) 不是\\
2. 他汉语说\_\_\_很好。(a) 的 (b) 得 (c) 地\\
3. 我\_\_\_一个哥哥。(a) 不有 (b) 没有 (c) 是没\\
4. 我们七点\_\_\_上课。(a) 就 (b) 也 (c) 都\\
5. 这本书\_\_\_那本书有意思。(a) 比 (b) 很 (c) 最

\textbf{III. Traduction (6 min, 4 points).} Traduisez en chinois (caractères + pinyin).\\
1. Je suis élève à Yaoundé. 2. Je n'ai pas de petit frère. 3. Hier, je ne suis pas allé au marché. 4. Le chinois est plus intéressant que l'anglais.

\textbf{IV. Expression écrite (6 min, 6 points).} Rédigez au moins 5 phrases : présentez-vous et parlez de votre week-end (nom, ville, activités, ce que vous aimez).
\end{exercice}

\begin{corrige}[Épreuve blanche]
\textbf{I. Compréhension (1 pt par réponse).} 1. À Douala (\zh{杜阿拉}). 2. Avec sa maman. 3. Il joue au football avec ses amis. 4. Non, il ne joue pas très bien (\zh{踢得不太好}). 5. Parce que le chinois est très intéressant (\zh{很有意思}).

\textbf{II. QCM (1 pt par réponse).} 1. (b) \zh{没} — passé nié. 2. (b) \zh{得} — complément de degré après le verbe. 3. (b) \zh{没有} — \zh{有} se nie par \zh{没}. 4. (a) \zh{就} — « dès sept heures ». 5. (a) \zh{比} — comparatif : A + 比 + B + adjectif.

\textbf{III. Traduction (1 pt par phrase).}\\
1. \zh{我是雅温得的学生。} \emph{Wǒ shì Yǎwēndé de xuésheng.}\\
2. \zh{我没有弟弟。} \emph{Wǒ méiyǒu dìdi.}\\
3. \zh{我昨天没去市场。} \emph{Wǒ zuótiān méi qù shìchǎng.} (temps avant le verbe !)\\
4. \zh{汉语比英语有意思。} \emph{Hànyǔ bǐ Yīngyǔ yǒu yìsi.}

\textbf{IV. Expression écrite (barème).} 5 phrases complètes avec verbe = 2 pts ; ordre des mots correct = 2 pts ; caractères lisibles et justes = 1 pt ; au moins une structure de l'année (\zh{因为}, \zh{比}, \zh{得}) = 1 pt.
\end{corrige}

\begin{exoresolu}[Auto-correction guidée : la grille d'erreurs types]
Énoncé : un élève a écrit en expression écrite : \zh{我去学校昨天。我没有不弟弟。} Corrigez sa copie avec la grille d'erreurs.
\tcblower
\textbf{Erreur 1 — ordre des mots :} \zh{我去学校昨天。} Le complément de temps \zh{昨天} doit se placer avant le verbe. Correction : \zh{我昨天去学校。} (\emph{Wǒ zuótiān qù xuéxiào.}) « Hier, je suis allé à l'école. »\\
\textbf{Erreur 2 — double négation fautive :} \zh{没有不} est impossible ; \zh{有} se nie par \zh{没} seul. Correction : \zh{我没有弟弟。} (\emph{Wǒ méiyǒu dìdi.}) « Je n'ai pas de petit frère. »\\
\textbf{Grille d'erreurs types à appliquer à votre copie :} (1) verbe présent dans chaque phrase ? (2) temps/lieu avant le verbe ? (3) négation correcte (\zh{不} / \zh{没}) ? (4) \zh{的}/\zh{得} au bon endroit ? (5) caractères lisibles ?
\end{exoresolu}

\courssub{Bilan personnel et complément utile}

Après l'auto-correction, chaque élève remplit son bilan : notez \textbf{3 points forts} (ex. « je maîtrise le comparatif avec 比 ») et \textbf{3 points à retravailler} (ex. « je confonds encore 不 et 没 »). Ce bilan est votre programme de révision personnel pour les derniers jours avant l'examen réel.

\begin{aretenir}[Complément utile à l'examen (hors-programme signalé)]
Ces deux formules ne sont pas exigées par le programme, mais elles font toujours bonne impression à l'écrit comme à l'oral :
\begin{itemize}
\item \zh{谢谢老师！} (\emph{Xièxie lǎoshī!}) « Merci, professeur ! » — formule de politesse écrite, à placer à la fin d'une lettre ou d'une production.
\item \zh{我叫……，我是……的学生。} (\emph{Wǒ jiào..., wǒ shì ... de xuésheng.}) « Je m'appelle..., je suis élève de/à... » — présentation éclair à personnaliser.
\end{itemize}
\end{aretenir}

\begin{textebox}[Modèle de présentation à mémoriser]
我叫……，我是雅温得的学生。\\
\emph{Wǒ jiào ..., wǒ shì Yǎwēndé de xuésheng.}\\
« Je m'appelle..., je suis élève à Yaoundé. »\\
我喜欢汉语，因为汉语很有意思。\\
\emph{Wǒ xǐhuan Hànyǔ, yīnwèi Hànyǔ hěn yǒu yìsi.}\\
« J'aime le chinois, parce que le chinois est très intéressant. »
\end{textebox}

\begin{savaistu}[La communication réelle est récompensée]
Au Cameroun, les épreuves de chinois valorisent la \cle{communication réelle} : un correcteur préfère toujours une phrase simple et vraie sur vous-même (\zh{我喜欢足球。}) à une phrase compliquée copiée et pleine de fautes. Dire quelque chose d'authentique — votre ville, votre famille, vos goûts — est la stratégie la plus sûre. En Chine aussi, les professeurs disent aux débutants : « 说真话 » (\emph{shuō zhēnhuà}), « dis des choses vraies ».
\end{savaistu}

\begin{experience}[Dernière répétition générale]
Par groupes de quatre : (1) chacun passe l'épreuve blanche en silence, chronomètre en main ; (2) échangez les copies et corrigez avec la grille d'erreurs types ; (3) chacun annonce à voix haute ses 3 points forts et ses 3 points à retravailler ; (4) terminez par une présentation orale éclair avec \zh{我叫……，我是……的学生。} devant le groupe.
\end{experience}

Vous êtes désormais prêts : vocabulaire révisé, structures synthétisées, caractères retravaillés, stratégie d'épreuve en place. Le jour de l'examen, respirez, lisez les consignes, commencez par le facile — et écrivez des choses vraies. \zh{加油！} (\emph{Jiāyóu!}) « Courage ! »

\newpage

\coursec{Activité d'intégration}

\courssub{Objectifs de l'activité}

À l'issue de cette activité d'intégration, l'élève sera capable de :
\begin{itemize}
\item mobiliser rapidement le vocabulaire de l'année (salutations, famille, école, temps, nourriture, transports, médias) à partir de définitions ;
\item repérer et corriger les erreurs grammaticales typiques des francophones (\zh{是} / \zh{在}, \zh{不} / \zh{没}, ordre des mots, place des compléments de temps et de lieu) ;
\item écrire des caractères fréquents en respectant l'ordre des traits et les radicaux ;
\item traduire des phrases courtes dans les deux sens (thème et version) ;
\item produire oralement une mini-présentation cohérente de cinq phrases en chinois ;
\item évaluer ses propres points forts et ses points à revoir avant l'examen.
\end{itemize}

\courssub{Situation}

C'est la fin de l'année scolaire au lycée de Yaoundé. Dans une semaine, la classe de Première A passera son évaluation de chinois. La professeure, Mme Ngo Bell, propose une grande révision sous forme de jeu : le « Grand défi des sinisants ». La classe est divisée en équipes de quatre élèves. Chaque équipe porte le nom d'une ville chinoise : équipe \zh{北京} (Běijīng), équipe \zh{上海} (Shànghǎi), équipe \zh{广州} (Guǎngzhōu), équipe \zh{西安} (Xī'ān).

Pour motiver les élèves, Mme Ngo Bell a affiché au tableau le règlement du jeu, rédigé en chinois :

\begin{textebox}[Le règlement du jeu — affiché au tableau]
同学们好！今天我们一起复习。\\
Tóngxuémen hǎo! Jīntiān wǒmen yìqǐ fùxí.\\
« Bonjour à tous ! Aujourd'hui, nous révisons ensemble. »\\[4pt]
我们班有四个队。每个队有四个学生。\\
Wǒmen bān yǒu sì ge duì. Měi ge duì yǒu sì ge xuésheng.\\
« Notre classe a quatre équipes. Chaque équipe a quatre élèves. »\\[4pt]
有五个活动：生词、语法、汉字、翻译、口语。\\
Yǒu wǔ ge huódòng: shēngcí, yǔfǎ, Hànzì, fānyì, kǒuyǔ.\\
« Il y a cinq ateliers : vocabulaire, grammaire, caractères, traduction, expression orale. »\\[4pt]
每个活动六分钟。加油！\\
Měi ge huódòng liù fēnzhōng. Jiāyóu!\\
« Chaque atelier dure six minutes. Bon courage ! »
\end{textebox}

Avant de commencer, deux élèves, Amina et Étienne, vérifient qu'ils ont bien compris :

\begin{textebox}[Dialogue entre Amina et Étienne]
Amina : 你准备好了吗？\\
Nǐ zhǔnbèi hǎo le ma?\\
« Es-tu prêt ? »\\[4pt]
Étienne : 我有点儿紧张，但是我很努力。\\
Wǒ yǒudiǎnr jǐnzhāng, dànshì wǒ hěn nǔlì.\\
« Je suis un peu stressé, mais je travaille beaucoup. »\\[4pt]
Amina : 别担心！我们一起复习吧。\\
Bié dānxīn! Wǒmen yìqǐ fùxí ba.\\
« Ne t'inquiète pas ! Révisons ensemble. »\\[4pt]
Étienne : 好！我们先复习生词，再复习语法。\\
Hǎo! Wǒmen xiān fùxí shēngcí, zài fùxí yǔfǎ.\\
« D'accord ! Révisons d'abord le vocabulaire, puis la grammaire. »
\end{textebox}

\courssub{Consignes}

Le parcours comporte cinq ateliers chronométrés de six minutes chacun. Chaque atelier rapporte jusqu'à 10 points. Travaillez en équipe, écrivez vos réponses sur la feuille de route, et parlez chinois autant que possible.

\begin{enumerate}
\item \textbf{Atelier 1 — « Mots éclair » (compréhension et lexique).} L'enseignant lit dix définitions en français (exemple : « le repas du midi », « le père de mon père », « regarder la télévision »). L'équipe écrit le mot chinois correspondant en caractères et en pinyin. Un point par mot correct (caractères + tons justes).
\item \textbf{Atelier 2 — « Phrases réparées » (grammaire).} Cinq phrases contenant chacune une erreur typique sont projetées au tableau (confusion \zh{是} / \zh{在}, \zh{不} / \zh{没}, ordre des mots, complément de lieu mal placé, classificateur manquant). L'équipe réécrit chaque phrase correctement et nomme la règle en français.
\item \textbf{Atelier 3 — « Caractères mystères » (écriture).} Un élève de l'équipe vient au tableau. L'enseignant dicte cinq caractères de l'année (exemples : \zh{好}, \zh{学}, \zh{饭}, \zh{朋}, \zh{看}). L'élève les écrit en respectant l'ordre des traits ; l'équipe doit ensuite donner le radical (la clé) de chaque caractère.
\item \textbf{Atelier 4 — « Traduction express » (thème et version).} Trois phrases à traduire du français vers le chinois, trois phrases du chinois vers le français. Attention aux tons, à l'ordre des mots et aux mots-outils \zh{的}, \zh{了}, \zh{吗}.
\item \textbf{Atelier 5 — « Mini-exposé » (production orale).} Chaque équipe invente un personnage (un élève camerounais ou chinois) et le présente oralement en cinq phrases minimum : identité, famille, journée, loisirs, projet. Un point par phrase correcte et prononcée avec des tons acceptables.
\end{enumerate}

À la fin du parcours, correction collective au tableau, calcul des points, puis chaque élève remplit individuellement sa fiche-bilan « Mes points forts / À revoir ».

\courssub{Ressources à mobiliser}

\begin{aretenir}[Structures-clés de l'année à réactiver]
\begin{itemize}
\item Phrase nominale avec \zh{是} : Sujet + \zh{是} + nom. \zh{我是学生。} (Wǒ shì xuésheng. « Je suis élève. ») — jamais \zh{是} devant un adjectif : on dit \zh{我很好} (Wǒ hěn hǎo).
\item Localisation avec \zh{在} : Sujet + \zh{在} + lieu (+ verbe). \zh{我在雅温得学习。} (Wǒ zài Yǎwēndé xuéxí. « J'étudie à Yaoundé. »)
\item Négation : \zh{不} pour le présent/futur et les états ; \zh{没(有)} pour le passé et la possession niée. \zh{我没吃饭。} (Wǒ méi chīfàn. « Je n'ai pas mangé. »)
\item Ordre des mots : Sujet + temps + lieu + verbe + complément. \zh{我明天在学校见朋友。} (Wǒ míngtiān zài xuéxiào jiàn péngyou.)
\item Question avec \zh{吗} : phrase affirmative + \zh{吗} + \zh{？}
\item Possession et description avec \zh{的} : \zh{我的老师} (wǒ de lǎoshī, « mon professeur »).
\item Classificateurs : \zh{一个人}, \zh{一本书}, \zh{一杯茶}.
\end{itemize}
\end{aretenir}

\begin{center}
\begin{tabularx}{\linewidth}{|X|X|X|X|}
\hline
\rowcolor{popblueL} Caractères & Pinyin & Nature & Traduction \\
\hline
复习 & fùxí & verbe & réviser \\
\hline
队 & duì & nom / classif. & équipe \\
\hline
活动 & huódòng & nom & activité, atelier \\
\hline
生词 & shēngcí & nom & mot nouveau, vocabulaire \\
\hline
语法 & yǔfǎ & nom & grammaire \\
\hline
汉字 & Hànzì & nom & caractère chinois \\
\hline
翻译 & fānyì & verbe / nom & traduire ; traduction \\
\hline
口语 & kǒuyǔ & nom & langue parlée, oral \\
\hline
紧张 & jǐnzhāng & adjectif & stressé, tendu \\
\hline
努力 & nǔlì & adjectif / verbe & travailleur ; faire des efforts \\
\hline
担心 & dānxīn & verbe & s'inquiéter \\
\hline
加油 & jiāyóu & expression & courage ! allez ! \\
\hline
\end{tabularx}
\end{center}

\begin{attention}[Pièges fréquents à surveiller pendant le jeu]
Ne confondez pas \zh{是} (être + nom) et \zh{在} (être situé, se trouver) : on dit \zh{我在学校} et non *\zh{我是学校}. N'utilisez pas \zh{不} pour nier une action passée : \zh{我昨天没上课} et non *\zh{我昨天不上课}. Placez le complément de temps avant le verbe : \zh{我八点起床} et non *\zh{我起床八点}. N'oubliez pas le classificateur entre le nombre et le nom : \zh{三个朋友}.
\end{attention}

\courssub{Proposition de production et corrigé commenté}

\begin{exoresolu}[Atelier 5 — Mini-exposé : production modèle de l'équipe « Běijīng »]
Énoncé : présentez en cinq phrases minimum un personnage fictif (identité, famille, journée, loisirs, projet).
\tcblower
\textbf{Production modèle :}\\[4pt]
\zh{他叫小马，是杜阿拉的中学生。}\\
Tā jiào Xiǎo Mǎ, shì Dù'ālā de zhōngxuésheng.\\
« Il s'appelle Xiao Ma, c'est un collégien de Douala. »\\[3pt]
\zh{他家有五口人：爸爸、妈妈、两个妹妹和他。}\\
Tā jiā yǒu wǔ kǒu rén: bàba, māma, liǎng ge mèimei hé tā.\\
« Sa famille compte cinq personnes : papa, maman, deux petites sœurs et lui. »\\[3pt]
\zh{他每天早上六点起床，七点去学校。}\\
Tā měitiān zǎoshang liù diǎn qǐchuáng, qī diǎn qù xuéxiào.\\
« Chaque matin, il se lève à six heures et va à l'école à sept heures. »\\[3pt]
\zh{他喜欢踢足球，也喜欢看手机上的视频。}\\
Tā xǐhuan tī zúqiú, yě xǐhuan kàn shǒujī shang de shìpín.\\
« Il aime jouer au football et regarder des vidéos sur son téléphone. »\\[3pt]
\zh{明年他想去中国学习汉语。}\\
Míngnián tā xiǎng qù Zhōngguó xuéxí Hànyǔ.\\
« L'année prochaine, il voudrait aller en Chine pour apprendre le chinois. »\\[6pt]
\textbf{Commentaire des choix de langue :}
\begin{itemize}
\item Phrase 1 : \zh{叫} + nom pour l'identité, puis \zh{是} + nom (phrase nominale correcte, sans adjectif après \zh{是}).
\item Phrase 2 : structure de dénombrement familial \zh{有} + nombre + \zh{口} + \zh{人}, classificateur \zh{个} devant \zh{妹妹}, et \zh{两} (et non \zh{二}) devant un classificateur.
\item Phrase 3 : ordre Sujet + temps (\zh{每天早上六点}) + verbe ; les heures utilisent \zh{点}.
\item Phrase 4 : coordination des goûts avec \zh{也} placé avant le verbe ; \zh{手机上} = « sur le téléphone » (localisation avec \zh{上}).
\item Phrase 5 : projet exprimé avec \zh{想} + verbe ; le complément de lieu \zh{去中国} précède le verbe \zh{学习}.
\end{itemize}
\end{exoresolu}

\begin{exoresolu}[Atelier 2 — Exemple de phrase réparée]
Énoncé : corrigez la phrase \zh{我昨天不去学校。} et expliquez la règle.
\tcblower
Correction : \zh{我昨天没去学校。} (Wǒ zuótiān méi qù xuéxiào. « Hier, je ne suis pas allé à l'école. »)\\
Règle : pour nier une action passée (marquée ici par \zh{昨天}), on utilise \zh{没(有)} et non \zh{不}. La négation \zh{不} s'emploie pour le présent, le futur et les états habituels.
\end{exoresolu}

\courssub{Bilan de l'activité}

Ce grand jeu a permis de réactiver, dans un cadre coopératif et chronométré proche des conditions d'examen, l'ensemble des compétences de l'année : compréhension orale des définitions, mobilisation du lexique, correction grammaticale argumentée, écriture des caractères avec leurs radicaux, traduction dans les deux sens et production orale guidée. La correction collective a mis en lumière les erreurs récurrentes de la classe (\zh{是} / \zh{在}, \zh{不} / \zh{没}, place des compléments de temps), que chaque élève a notées dans sa fiche-bilan.

\begin{methode}[Ma fiche-bilan « Prêt pour l'examen »]
\begin{enumerate}
\item J'écris mes trois points forts de la journée (exemple : « je connais mes classificateurs »).
\item J'écris mes deux points à revoir, avec la règle correspondante en français et un exemple en chinois correct.
\item Je choisis cinq caractères à réécrire dix fois chez moi, en respectant l'ordre des traits.
\item Je prépare une question à poser au professeur lors de la prochaine séance de révision.
\end{enumerate}
\end{methode}

\begin{experience}[Prolongement : défi inter-équipes à la maison]
Chaque équipe rédige cinq nouvelles questions de quiz (vocabulaire, grammaire, traduction) avec leurs corrigés. Lors de la séance suivante, les équipes s'interrogent mutuellement : une bonne question vaut un point, une bonne réponse en vaut deux. L'enseignant valide la correction grammaticale des questions avant le jeu.
\end{experience}

\newpage

\coursec{Méthodes et exercices résolus}

\courssub{Fiche méthode 1 — Mobiliser ses acquis : la stratégie de révision efficace}

\begin{methode}[Comment réviser tout le programme de l'année ?]
\begin{enumerate}
\item \textbf{Faire l'inventaire.} Liste sur une feuille les grands thèmes de l'année (salutations, famille, école, temps, nourriture, médias, téléphone, Internet) et les structures clés (ordre Sujet + Verbe + Complément, 了, 过, 比, classificateurs).
\item \textbf{Classer le vocabulaire par champs lexicaux.} Un mot s'apprend avec son classificateur et un exemple : \zh{一部手机} \emph{yí bù shǒujī} « un téléphone portable ».
\item \textbf{Fabriquer des fiches de structures.} Pour chaque règle : schéma + un exemple personnel + traduction. Exemple : A + 比 + B + adjectif → \zh{中国比喀麦隆大。} \emph{Zhōngguó bǐ Kāmàilóng dà.} « La Chine est plus grande que le Cameroun. »
\item \textbf{Réécrire les caractères pièges.} Refaire 5 fois les caractères confondus (人/入, 未/末, 已/己) en respectant l'ordre des traits : haut → bas, gauche → droite, horizontal avant vertical.
\item \textbf{S'auto-tester.} Couvrir la colonne pinyin ou traduction du tableau de vocabulaire et se tester ; refaire les exercices déjà corrigés sans regarder la solution.
\item \textbf{Travailler en temps limité.} Refaire un sujet type en 1 h 30, comme à l'examen.
\end{enumerate}
\end{methode}

\courssub{Bilan du vocabulaire de l'année : le champ lexical des médias}

Pour le module « Mass médias et télécommunication », voici le lexique à maîtriser absolument. Apprends chaque mot avec son pinyin exact et un exemple.

\begin{center}
\begin{tabularx}{\linewidth}{|X|X|X|X|}
\hline
\rowcolor{popblueL} Caractères & Pinyin & Nature & Traduction \\
\hline
手机 & shǒujī & nom & téléphone portable \\
\hline
电话 & diànhuà & nom & téléphone \\
\hline
上网 & shàng wǎng & verbe & aller sur Internet \\
\hline
网络 & wǎngluò & nom & Internet, réseau \\
\hline
新闻 & xīnwén & nom & informations, actualités \\
\hline
消息 & xiāoxi & nom & message, nouvelle \\
\hline
广播 & guǎngbō & nom & radio (diffusion) \\
\hline
电视 & diànshì & nom & télévision \\
\hline
电影 & diànyǐng & nom & film, cinéma \\
\hline
报纸 & bàozhǐ & nom & journal (presse écrite) \\
\hline
发 & fā & verbe & envoyer (un message) \\
\hline
收 & shōu & verbe & recevoir \\
\hline
用 & yòng & verbe & utiliser \\
\hline
媒体 & méitǐ & nom & médias \\
\hline
有用 & yǒuyòng & adjectif & utile \\
\hline
\end{tabularx}
\end{center}

\begin{aretenir}[Les classificateurs à connaître par cœur]
\begin{itemize}
\item \cle{个} \emph{ge} : classificateur général → \zh{一个人} \emph{yí ge rén} « une personne ».
\item \cle{本} \emph{běn} : livres, magazines → \zh{一本书} \emph{yì běn shū} « un livre ».
\item \cle{部} \emph{bù} : téléphones, films → \zh{一部电影} \emph{yí bù diànyǐng} « un film ».
\item \cle{杯} \emph{bēi} : verres, tasses → \zh{一杯茶} \emph{yì bēi chá} « une tasse de thé ».
\item \cle{张} \emph{zhāng} : objets plats (feuilles, tables, photos) → \zh{一张报纸} \emph{yì zhāng bàozhǐ} « un journal ».
\item \cle{次} \emph{cì} : fois (fréquence) → \zh{两次} \emph{liǎng cì} « deux fois ».
\end{itemize}
Rappel : devant un classificateur, « deux » se dit toujours \cle{两} \emph{liǎng}, jamais 二.
\end{aretenir}

\begin{exoresolu}[Quiz de mobilisation des acquis]
\textbf{Énoncé.} Réponds aux questions suivantes en chinois, avec pinyin et traduction : (a) Comment dit-on « J'ai déjà regardé ce film » ? (b) Construis une phrase comparative avec 比 sur Douala et Yaoundé. (c) Quel classificateur utilise-t-on avec 手机 ?
\tcblower
\textbf{Raisonnement.} (a) « Déjà fait » = action accomplie ou expérience → 了 (action ponctuelle achevée) ou 过 (expérience vécue). « Regarder » = 看. (b) Structure : A + 比 + B + adjectif. (c) Le classificateur des téléphones est 部.
\begin{itemize}
\item (a) \zh{我已经看过这个电影了。} \emph{Wǒ yǐjing kàn guo zhège diànyǐng le.} « J'ai déjà vu ce film. » — 已经 marque « déjà », 过 marque l'expérience vécue, le 了 final signale la situation acquise.
\item (b) \zh{杜阿拉比雅温得热。} \emph{Dù'ālā bǐ Yǎwēndé rè.} « Douala est plus chaude que Yaoundé. » — Après 比, l'adjectif seul suffit : pas de 很 (很 annulerait l'idée de comparaison).
\item (c) On utilise 部 : \zh{一部手机} \emph{yí bù shǒujī}. Attention au changement de ton de 一 : yī → yí devant un 4\textsuperscript{e} ton.
\end{itemize}
\textbf{Conclusion.} Une bonne réponse mobilise toujours trois éléments : la structure correcte, le mot-outil juste (了/过/比) et le classificateur approprié.
\end{exoresolu}

\courssub{Fiche méthode 2 — Réussir la compréhension écrite}

\begin{methode}[Lire et comprendre un texte chinois à l'examen]
\begin{enumerate}
\item \textbf{Lire d'abord les questions} pour savoir quoi chercher dans le texte.
\item \textbf{Première lecture globale} : identifier le thème (qui ? où ? quoi ?) sans s'arrêter sur les caractères inconnus.
\item \textbf{Repérer les mots-clés} : noms propres (北京, 杜阿拉), chiffres, mots de temps (昨天, 现在, 以后) et connecteurs (因为…所以…, 虽然…但是…).
\item \textbf{Deuxième lecture ciblée} : souligner la phrase qui contient la réponse à chaque question.
\item \textbf{Déduire les mots inconnus} par le contexte et les radicaux : 语 (clé de la parole 讠) → lien avec « langue » ; 饭 (clé de la nourriture 饣) → lien avec « manger ».
\item \textbf{Rédiger la réponse en chinois complet} (Sujet + Verbe + Complément), jamais un mot isolé, et vérifier que la réponse correspond bien à la question posée.
\end{enumerate}
\end{methode}

\begin{exoresolu}[Compréhension écrite : texte sur les médias au Cameroun]
\textbf{Texte.} \zh{现在，很多喀麦隆年轻人每天都用手机上网。他们看新闻，听音乐，也给朋友发消息。杜阿拉的马丽说：“我早上听广播，晚上看电视。网络让我知道世界上发生的事情。”}

\emph{Xiànzài, hěn duō Kāmàilóng niánqīngrén měi tiān dōu yòng shǒujī shàng wǎng. Tāmen kàn xīnwén, tīng yīnyuè, yě gěi péngyou fā xiāoxi. Dù'ālā de Mǎlì shuō : « Wǒ zǎoshang tīng guǎngbō, wǎnshang kàn diànshì. Wǎngluò ràng wǒ zhīdào shìjiè shang fāshēng de shìqing. »}

« Aujourd'hui, beaucoup de jeunes Camerounais utilisent leur téléphone pour aller sur Internet chaque jour. Ils lisent les informations, écoutent de la musique et envoient aussi des messages à leurs amis. Marie de Douala dit : "Le matin j'écoute la radio, le soir je regarde la télévision. Internet me permet de savoir ce qui se passe dans le monde." »

\textbf{Questions.} (1) Que font les jeunes Camerounais avec leur téléphone ? (2) Quand Marie écoute-t-elle la radio ? (3) Que permet Internet à Marie ?
\tcblower
\textbf{Raisonnement.} (1) La réponse est dans la deuxième phrase : trois actions coordonnées (看新闻, 听音乐, 发消息). (2) Chercher le mot de temps associé à 听广播 : 早上. (3) La structure 让我知道… donne directement la réponse.
\begin{itemize}
\item (1) \zh{他们用手机上网，看新闻，听音乐，也给朋友发消息。} \emph{Tāmen yòng shǒujī shàng wǎng, kàn xīnwén, tīng yīnyuè, yě gěi péngyou fā xiāoxi.} « Ils utilisent leur téléphone pour aller sur Internet, lire les informations, écouter de la musique et envoyer des messages à leurs amis. »
\item (2) \zh{她早上听广播。} \emph{Tā zǎoshang tīng guǎngbō.} « Elle écoute la radio le matin. » — Le complément de temps se place avant le verbe, jamais après.
\item (3) \zh{网络让她知道世界上发生的事情。} \emph{Wǎngluò ràng tā zhīdào shìjiè shang fāshēng de shìqing.} « Internet lui permet de savoir ce qui se passe dans le monde. »
\end{itemize}
\textbf{Conclusion.} Chaque réponse reprend les mots du texte dans une phrase complète et correcte : c'est ce que l'examinateur attend.
\end{exoresolu}

\courssub{Fiche méthode 3 — Maîtriser les structures et les mots-outils}

\begin{propriete}[Synthèse des mots-outils de l'année]
\begin{center}
\begin{tabularx}{\linewidth}{|X|X|X|}
\hline
\rowcolor{popblueL} Mot-outil et fonction & Exemple (caractères + pinyin) & Traduction \\
\hline
了 : action achevée & 我吃了饭。Wǒ chī le fàn. & J'ai mangé. \\
\hline
过 : expérience vécue & 我去过北京。Wǒ qù guo Běijīng. & Je suis déjà allé à Pékin. \\
\hline
的 : possession, description & 我的手机 wǒ de shǒujī & mon téléphone \\
\hline
得 : manière, degré & 他说得很好。Tā shuō de hěn hǎo. & Il parle très bien. \\
\hline
地 : adverbe de manière & 她高兴地说。Tā gāoxìng de shuō. & Elle dit joyeusement. \\
\hline
比 : comparaison & 汉语比法语难。Hànyǔ bǐ Fǎyǔ nán. & Le chinois est plus difficile que le français. \\
\hline
会 : futur, capacité & 我明天会去。Wǒ míngtiān huì qù. & J'irai demain. \\
\hline
\end{tabularx}
\end{center}
\end{propriete}

\begin{methode}[Choisir le bon mot-outil : 了, 过, 的, 得, 地]
\begin{enumerate}
\item \textbf{Identifier la fonction.} Action achevée → 了 après le verbe ; expérience vécue → 过 ; possession ou description → 的 ; degré de l'action (comment ?) → 得 après le verbe ; adverbe devant un verbe → 地.
\item \textbf{Vérifier la place.} 了 se met juste après le verbe : \zh{我吃了饭。} ; 得 se met après le verbe, avant le complément de manière : \zh{他说得很好。}
\item \textbf{Mémoriser les schémas.} Sujet + Verbe + 了 + Complément ; Sujet + Verbe + 得 + adjectif ; Nom + 的 + Nom.
\item \textbf{Tester par traduction.} Si le français dit « il parle bien », c'est un jugement sur la manière → 得, pas 的.
\item \textbf{Attention à la négation.} Le passé accompli se nie avec 没(有), et 了 disparaît : \zh{我没吃饭。} \emph{Wǒ méi chī fàn.} « Je n'ai pas mangé. »
\end{enumerate}
\end{methode}

\begin{popfigure}
\begin{tikzpicture}[node distance=0.6cm and 1.1cm,
  question/.style={rectangle, rounded corners, draw=popblue, fill=popblueL, align=center, font=\small, inner sep=5pt},
  reponse/.style={rectangle, rounded corners, draw=popgreen, fill=popgreenL, align=center, font=\small, inner sep=5pt},
  fleche/.style={-{Stealth[length=2.5mm]}, thick, popdark}]
\node[question] (q1) {Action achevée ?};
\node[reponse, right=of q1, align=center] (r1){了 après le verbe\\ 我吃了饭。};
\node[question, below=of q1] (q2) {Expérience vécue ?};
\node[reponse, right=of q2, align=center] (r2){过 après le verbe\\ 我去过北京。};
\node[question, below=of q2] (q3) {Comment ? (manière)};
\node[reponse, right=of q3, align=center] (r3){得 après le verbe\\ 他说得很好。};
\node[question, below=of q3] (q4) {Possession ?};
\node[reponse, right=of q4, align=center] (r4){的 entre les noms\\ 我的手机};
\draw[fleche] (q1) -- (r1);
\draw[fleche] (q2) -- (r2);
\draw[fleche] (q3) -- (r3);
\draw[fleche] (q4) -- (r4);
\end{tikzpicture}
\legende{Arbre de décision : quel mot-outil choisir ?}
\end{popfigure}

\begin{exoresolu}[Grammaire : compléter avec le bon mot-outil]
\textbf{Énoncé.} Complète avec 了, 过, 的, 得 ou 地 : (1) 我吃\_\_\_饭，现在不饿。(2) 他汉语说\_\_\_很好。(3) 这是我\_\_\_手机。(4) 我去\_\_\_北京两次。(5) 她高兴\_\_\_说：“太好了！”
\tcblower
\textbf{Raisonnement et solutions.}
\begin{itemize}
\item (1) 了 — action achevée (j'ai mangé, donc je n'ai plus faim) : \zh{我吃了饭，现在不饿。} \emph{Wǒ chī le fàn, xiànzài bù è.} « J'ai mangé, maintenant je n'ai pas faim. »
\item (2) 得 — on évalue la manière de parler : \zh{他汉语说得很好。} \emph{Tā Hànyǔ shuō de hěn hǎo.} « Il parle très bien chinois. »
\item (3) 的 — possession : \zh{这是我的手机。} \emph{Zhè shì wǒ de shǒujī.} « C'est mon téléphone. »
\item (4) 过 — expérience vécue (deux fois dans la vie) : \zh{我去过北京两次。} \emph{Wǒ qù guo Běijīng liǎng cì.} « Je suis allé deux fois à Pékin. »
\item (5) 地 — adverbe de manière devant le verbe dire : \zh{她高兴地说：“太好了！”} \emph{Tā gāoxìng de shuō : « Tài hǎo le ! »} « Elle dit joyeusement : "C'est génial !" »
\end{itemize}
\textbf{Conclusion.} Le bon réflexe : se demander « accompli ? expérience ? possession ? manière ? » avant d'écrire le mot-outil.
\end{exoresolu}

\courssub{Caractères : radicaux, traits et pièges}

\begin{propriete}[Les paires de caractères à ne plus confondre]
\begin{itemize}
\item \cle{人} \emph{rén} « personne » / \cle{入} \emph{rù} « entrer » : dans 人, le trait de gauche part en premier et les deux traits se rejoignent en haut ; dans 入, le trait de droite est plus long et dépasse.
\item \cle{未} \emph{wèi} « pas encore » / \cle{末} \emph{mò} « fin » : dans 未, la barre du haut est courte ; dans 末, c'est la barre du haut qui est longue.
\item \cle{已} \emph{yǐ} « déjà » / \cle{己} \emph{jǐ} « soi-même » : dans 已, la boucle est à moitié fermée ; dans 己, elle reste ouverte en bas.
\item \cle{买} \emph{mǎi} « acheter » / \cle{卖} \emph{mài} « vendre » : 卖 a un trait supplémentaire en haut (on « ajoute » quelque chose pour vendre).
\item \cle{问} \emph{wèn} « demander » / \cle{间} \emph{jiān} « entre, pièce » : 问 contient 口 (la bouche → demander), 间 contient 日 (le soleil dans la porte → l'intervalle).
\end{itemize}
Ordre général des traits : de haut en bas, de gauche à droite, l'horizontal avant le vertical, l'extérieur avant l'intérieur, et on ferme la boîte en dernier (ex. 国).
\end{propriete}

\begin{savaistu}[La clé de la parole 讠]
Beaucoup de caractères liés au langage portent la clé de la parole 讠: 说 \emph{shuō} « parler », 语 \emph{yǔ} « langue », 话 \emph{huà} « parole », 读 \emph{dú} « lire à voix haute », 认 \emph{rèn} « reconnaître ». Repérer cette clé dans un texte d'examen permet de deviner qu'un caractère inconnu a un rapport avec la parole ou la communication — très utile dans un module sur les médias !
\end{savaistu}

\courssub{Fiche méthode 4 — Réussir la traduction (thème et version)}

\begin{methode}[Traduire sans calquer le français]
\begin{enumerate}
\item \textbf{Identifier d'abord le squelette} : Sujet + Verbe + Complément. En chinois, l'ordre est rigide.
\item \textbf{Placer les compléments de temps et de lieu AVANT le verbe} : \zh{我明天在杜阿拉见他。} \emph{Wǒ míngtiān zài Dù'ālā jiàn tā.} « Je le vois demain à Douala. »
\item \textbf{Choisir le bon classificateur} : 个 (général), 本 (livres), 部 (téléphones, films), 杯 (verres, tasses), 张 (feuilles, tables, photos).
\item \textbf{Supprimer les mots français inutiles} : pas d'article (le/la/les), pas de conjugaison ; le temps est marqué par 了/过/会 ou par un mot de temps.
\item \textbf{Relire en vérifiant les tons} du pinyin et la cohérence des caractères (ne pas confondre 买 \emph{mǎi} « acheter » et 卖 \emph{mài} « vendre »).
\end{enumerate}
\end{methode}

\begin{exoresolu}[Thème : traduire en chinois]
\textbf{Énoncé.} Traduis en chinois (caractères + pinyin) : (1) « Hier soir, j'ai regardé la télévision avec mon frère. » (2) « Le chinois est plus difficile que le français, mais c'est intéressant. » (3) « Elle a acheté deux tasses de thé au marché de Yaoundé. »
\tcblower
\textbf{Raisonnement et solutions.}
\begin{itemize}
\item (1) Temps (昨天晚上) avant le verbe ; 和 + personne avant le verbe ; action achevée → 了. \zh{昨天晚上，我和哥哥一起看了电视。} \emph{Zuótiān wǎnshang, wǒ hé gēge yìqǐ kàn le diànshì.} « Hier soir, j'ai regardé la télévision avec mon grand frère. »
\item (2) Comparaison : A + 比 + B + adjectif ; contraste avec 但是. \zh{汉语比法语难，但是很有意思。} \emph{Hànyǔ bǐ Fǎyǔ nán, dànshì hěn yǒu yìsi.} « Le chinois est plus difficile que le français, mais c'est très intéressant. »
\item (3) Classificateur des tasses : 杯 ; « au marché » = complément de lieu avant le verbe. \zh{她在雅温得的市场买了两杯茶。} \emph{Tā zài Yǎwēndé de shìchǎng mǎi le liǎng bēi chá.} « Elle a acheté deux tasses de thé au marché de Yaoundé. » — Devant un classificateur, « deux » se dit 两, jamais 二.
\end{itemize}
\textbf{Conclusion.} Une traduction réussie = ordre des mots chinois + classificateur juste + mot de temps bien placé.
\end{exoresolu}

\courssub{Fiche méthode 5 — Expression écrite guidée et stratégies d'examen}

\begin{methode}[Rédiger un court texte et gérer l'épreuve]
\begin{enumerate}
\item \textbf{Lire la consigne deux fois} et souligner : le type de texte (message, lettre, paragraphe), le nombre de phrases, le thème.
\item \textbf{Faire un mini-plan} en 3 parties : introduction (qui, quand), développement (2-3 actions ou idées), conclusion (sentiment ou projet).
\item \textbf{Utiliser uniquement des structures maîtrisées} : mieux vaut des phrases simples correctes que des phrases longues fautives.
\item \textbf{Enrichir avec les connecteurs de l'année} : 因为…所以…, 虽然…但是…, 先…然后…, 除了…以外.
\item \textbf{Se relire avec la grille} : ordre des mots ? mots-outils ? classificateurs ? caractères bien formés ? tons du pinyin ?
\item \textbf{Gérer le temps} : 10 min de lecture, 60 min de travail, 10 min de relecture. Ne jamais laisser de question blanche.
\end{enumerate}
\end{methode}

\begin{exoresolu}[Expression écrite : un message à un ami]
\textbf{Énoncé.} Tu es André, élève en Première A à Yaoundé. Écris un message de 4 à 6 phrases à ton ami chinois 王明 pour lui parler de ton utilisation des médias (téléphone, télévision, Internet) et l'inviter à te répondre.
\tcblower
\textbf{Raisonnement.} Plan : salutation → mes habitudes (matin/soir, structures avec 先…然后… et 因为…所以…) → question à 王明 → invitation à répondre. Vocabulaire du module : 手机, 上网, 新闻, 消息, 广播, 电视.

\textbf{Modèle de rédaction.}

\zh{王明，你好！我是安德。我每天早上先听广播，然后用手机上网看新闻。晚上，我常常和家人一起看电视。因为网络很有用，所以我也用它学习汉语。你呢？你每天用什么媒体？请给我发消息。}

\emph{Wáng Míng, nǐ hǎo ! Wǒ shì Āndé. Wǒ měi tiān zǎoshang xiān tīng guǎngbō, ránhòu yòng shǒujī shàng wǎng kàn xīnwén. Wǎnshang, wǒ chángcháng hé jiārén yìqǐ kàn diànshì. Yīnwèi wǎngluò hěn yǒuyòng, suǒyǐ wǒ yě yòng tā xuéxí Hànyǔ. Nǐ ne ? Nǐ měi tiān yòng shénme méitǐ ? Qǐng gěi wǒ fā xiāoxi.}

« Wang Ming, bonjour ! Je suis André. Chaque matin, j'écoute d'abord la radio, puis j'utilise mon téléphone pour lire les informations sur Internet. Le soir, je regarde souvent la télévision avec ma famille. Parce qu'Internet est très utile, je l'utilise aussi pour apprendre le chinois. Et toi ? Quels médias utilises-tu chaque jour ? Envoie-moi un message, s'il te plaît. »

\textbf{Justifications.} 先…然后… enchaîne deux actions ; 因为…所以… exprime la cause ; 和…一起 « avec » se place avant le verbe ; 请 + verbe formule poliment l'invitation.

\textbf{Conclusion.} Six phrases simples, toutes correctes, avec deux connecteurs variés : c'est la copie idéale.
\end{exoresolu}

\begin{exercice}[Simulation d'examen (20 min)]
\begin{enumerate}
\item \textbf{Vocabulaire.} Donne le pinyin et la traduction de : 手机, 新闻, 广播, 上网, 消息.
\item \textbf{Grammaire.} Complète avec 了, 过, 得 ou 的 : (a) 他买\_\_\_一部手机。(b) 她唱歌唱\_\_\_很好。(c) 我去\_\_\_中国。(d) 这是我\_\_\_书。
\item \textbf{Traduction.} Traduis en chinois : « Mon petit frère regarde la télévision tous les soirs. »
\item \textbf{Expression écrite.} Écris 3 phrases sur ton média préféré avec 因为…所以….
\end{enumerate}
\end{exercice}

\begin{corrige}[Simulation d'examen]
\begin{enumerate}
\item 手机 \emph{shǒujī} « téléphone portable » ; 新闻 \emph{xīnwén} « informations » ; 广播 \emph{guǎngbō} « radio » ; 上网 \emph{shàng wǎng} « aller sur Internet » ; 消息 \emph{xiāoxi} « message ».
\item (a) 了 : \zh{他买了一部手机。} \emph{Tā mǎi le yí bù shǒujī.} « Il a acheté un téléphone portable. » (b) 得 : \zh{她唱歌唱得很好。} \emph{Tā chànggē chàng de hěn hǎo.} « Elle chante très bien. » (c) 过 : \zh{我去过中国。} \emph{Wǒ qù guo Zhōngguó.} « Je suis déjà allé en Chine. » (d) 的 : \zh{这是我的书。} \emph{Zhè shì wǒ de shū.} « C'est mon livre. »
\item \zh{我弟弟每天晚上看电视。} \emph{Wǒ dìdi měi tiān wǎnshang kàn diànshì.} « Mon petit frère regarde la télévision tous les soirs. » — « Tous les soirs » = complément de temps avant le verbe ; pas de 了 car c'est une habitude.
\item Exemple : \zh{因为我喜欢音乐，所以我每天听广播。} \emph{Yīnwèi wǒ xǐhuan yīnyuè, suǒyǐ wǒ měi tiān tīng guǎngbō.} « Parce que j'aime la musique, j'écoute la radio tous les jours. »
\end{enumerate}
\end{corrige}

\begin{attention}[Les 5 erreurs qui coûtent des points à l'examen]
\begin{enumerate}
\item \textbf{Les tons oubliés ou faux dans le pinyin.} Écrire \emph{ma} sans ton, ou confondre \emph{mǎi} (买, acheter) et \emph{mài} (卖, vendre), change le sens. Vérifie chaque voyelle : ā á ǎ à.
\item \textbf{Les caractères confondus.} 人/入, 未/末, 已/己, 买/卖, 问/间. Un trait de trop ou de moins rend le caractère faux : réécris-les en respectant l'ordre (haut → bas, gauche → droite).
\item \textbf{L'ordre des mots calqué sur le français.} Le complément de temps et de lieu se place AVANT le verbe : on écrit \zh{我昨天在杜阿拉看见他。} et jamais *\zh{我看见他昨天在杜阿拉。}
\item \textbf{Les mots-outils mal employés.} 了 avec la négation (*\zh{我没吃了饭} est FAUX → \zh{我没吃饭}) ; 的 au lieu de 得 pour la manière (*\zh{他说的很好} → \zh{他说得很好}) ; 二 au lieu de 两 devant un classificateur (\zh{两个}, pas *\zh{二个}).
\item \textbf{Les classificateurs oubliés ou faux.} Entre un nombre et un nom, le classificateur est obligatoire : \zh{一部手机} (téléphone), \zh{一本书} (livre), \zh{两杯茶} (deux tasses de thé). Écrire *\zh{一手机} est une faute grave.
\end{enumerate}
\end{attention}

\newpage

\coursec{Exercices}

\courssub{Je vérifie mes connaissances}

\begin{exercice}[QCM de grammaire : 10 items]
Pour chaque phrase, choisis la bonne réponse (A, B ou C) et entoure-la. Une seule réponse est correcte.
\begin{enumerate}
\item 我\underline{\hspace{1.5cm}}学生。 A. 是 \quad B. 在 \quad C. 有
\item 昨天我\underline{\hspace{1.5cm}}去学校。 A. 不 \quad B. 没 \quad C. 别
\item 他\underline{\hspace{1.5cm}}家\underline{\hspace{1.5cm}}三口人。 A. 有…了 \quad B. 家…有 \quad C. 有…在
\item 我喝\underline{\hspace{1.5cm}}茶。 A. 一个 \quad B. 一杯 \quad C. 一本
\item 她买\underline{\hspace{1.5cm}}两\underline{\hspace{1.5cm}}书。 A. 了…本 \quad B. 本…了 \quad C. 了…个
\item 我\underline{\hspace{1.5cm}}七点\underline{\hspace{1.5cm}}起床。 A. 在…都 \quad B. 每天…都 \quad C. 都…每天
\item 哥哥比我\underline{\hspace{1.5cm}}。 A. 很高 \quad B. 高 \quad C. 高得多一点儿
\item 请你说\underline{\hspace{1.5cm}}慢一点儿。 A. 的 \quad B. 得 \quad C. 地
\item 我\underline{\hspace{1.5cm}}去过北京。 A. 以前 \quad B. 已经 \quad C. 常常
\item 因为下雨，\underline{\hspace{1.5cm}}我们不去踢球。 A. 但是 \quad B. 所以 \quad C. 也
\end{enumerate}
\end{exercice}

\begin{exercice}[Vrai ou faux justifié]
Lis chaque phrase, dis si elle est \textbf{vraie} ou \textbf{fausse}, puis justifie ta réponse en citant la règle de grammaire concernée.
\begin{enumerate}
\item \zh{我不没吃早饭。} \emph{(Wǒ bù méi chī zǎofàn.)}
\item \zh{他是我的老师，也是我的朋友。} \emph{(Tā shì wǒ de lǎoshī, yě shì wǒ de péngyou.)}
\item \zh{我昨天去了商店买东西了。} (phrase correcte avec deux 了)
\item \zh{她汉语说得很好。} \emph{(Tā Hànyǔ shuō de hěn hǎo.)}
\item \zh{明天我有考试，所以今天晚上我想复习。} \emph{(Míngtiān wǒ yǒu kǎoshì, suǒyǐ jīntiān wǎnshang wǒ xiǎng fùxí.)}
\end{enumerate}
\end{exercice}

\begin{exercice}[Vocabulaire de l'année]
Complète chaque phrase avec le mot qui convient, choisi dans la liste : \zh{天气、爱好、医院、考试、电视、因为}.
\begin{enumerate}
\item 我头疼，要去\underline{\hspace{2cm}}看医生。
\item 我的\underline{\hspace{2cm}}是踢足球和听音乐。
\item 今天\underline{\hspace{2cm}}很好，我们去公园吧。
\item 明天有汉语\underline{\hspace{2cm}}，我要复习。
\item 晚上我喜欢看\underline{\hspace{2cm}}。
\item \underline{\hspace{2cm}}他病了，所以没来上课。
\end{enumerate}
\end{exercice}

\begin{exercice}[Pinyin et caractères]
Associe chaque caractère ou mot (colonne A) à son pinyin (colonne B). Attention aux tons !
\begin{center}
\begin{tabularx}{\linewidth}{|X|X|X|X|}
\hline
\rowcolor{popblueL} A — Caractères & & B — Pinyin & \\
\hline
1. 学 & & a. míngtiān & \\
\hline
2. 吃 & & b. xué & \\
\hline
3. 明天 & & c. chī & \\
\hline
4. 朋友 & & d. péngyou & \\
\hline
5. 喜欢 & & e. xǐhuan & \\
\hline
6. 中国 & & f. Zhōngguó & \\
\hline
\end{tabularx}
\end{center}
\end{exercice}

\courssub{Je m'entraîne}

\begin{exercice}[Traduction : du français au chinois]
Traduis les phrases suivantes en chinois (caractères obligatoires, pinyin en dessous).
\begin{enumerate}
\item Je m'appelle Paul. Ma famille a quatre personnes.
\item Hier je n'ai pas regardé la télévision.
\item J'aime le chinois parce qu'il est intéressant.
\item Mon grand frère est plus grand que moi.
\item Après les cours, je veux jouer au football avec mes amis.
\end{enumerate}
\end{exercice}

\begin{exercice}[Transformation de phrases]
Transforme chaque phrase selon la consigne entre parenthèses.
\begin{enumerate}
\item \zh{我喜欢喝茶。} (passer à la négation avec 不)
\item \zh{他昨天去了市场。} (passer à la négation du passé avec 没)
\item \zh{她学习汉语。} (ajouter 也 et 法语 pour dire « elle étudie aussi le français »)
\item \zh{今天很热。} (transformer en question avec 吗)
\item \zh{我去过北京。} (transformer en question avec 没有)
\end{enumerate}
\end{exercice}

\begin{exercice}[Texte à trous : mots-outils]
Complète le texte suivant avec les mots-outils qui conviennent : \zh{的、得、了、也、因为、在、比、过}. Chaque mot n'est utilisé qu'une seule fois.
\begin{textebox}[Texte à compléter]
我叫玛丽，我住\underline{\hspace{1cm}}雅温得。我家有五口人。我\underline{\hspace{1cm}}哥哥\underline{\hspace{1cm}}我大三岁，他\underline{\hspace{1cm}}学习汉语。去年他去\underline{\hspace{1cm}}中国。他说汉语说\underline{\hspace{1cm}}很好。\underline{\hspace{1cm}}汉语有意思，所以我每天学习两个小时。昨天我买\underline{\hspace{1cm}}一本新词典。\\
\emph{Wǒ jiào Mǎlì, wǒ zhù \_\_ Yǎwēndé. Wǒ jiā yǒu wǔ kǒu rén. Wǒ \_\_ gēge \_\_ wǒ dà sān suì, tā \_\_ xuéxí Hànyǔ. Qùnián tā qù \_\_ Zhōngguó. Tā shuō Hànyǔ shuō \_\_ hěn hǎo. \_\_ Hànyǔ yǒu yìsi, suǒyǐ wǒ měitiān xuéxí liǎng ge xiǎoshí. Zuótiān wǒ mǎi \_\_ yì běn xīn cídiǎn.}
\end{textebox}
\end{exercice}

\begin{exercice}[Remise en ordre et appariement]
\textbf{Partie A.} Remets les mots dans l'ordre pour faire une phrase correcte.
\begin{enumerate}
\item 我 / 每天 / 七点 / 起床 / 早上
\item 比 / 杜阿拉 / 雅温得 / 热
\item 了 / 他 / 两 / 杯 / 喝 / 茶
\item 因为 / 所以 / 下雨 / 我们 / 不 / 去 / 学校
\end{enumerate}
\textbf{Partie B.} Associe chaque début de phrase (1-4) à sa fin (a-d).
\begin{enumerate}
\item 我每天上课以前\underline{\hspace{2cm}} \quad a. 想去中国看看。
\item 下课以后\underline{\hspace{2cm}} \quad b. 先吃早饭。
\item 我学了汉语以后\underline{\hspace{2cm}} \quad c. 我和朋友们踢球。
\item 因为明天有考试，\underline{\hspace{2cm}} \quad d. 所以今天晚上我要复习。
\end{enumerate}
\end{exercice}

\courssub{Je me prépare au Bac}

\begin{exercice}[Compréhension écrite — type Baccalauréat]
Lis le texte suivant, puis réponds aux questions.

\begin{textebox}[Un élève camerounais apprend le chinois]
我叫恩东，我是雅温得一所中学的学生，今年十六岁。我们学校从去年开始有汉语课，我每个星期有三节汉语课。\\
我每天早上六点起床，七点去学校。汉语课很有意思，可是汉字不太容易，所以我每天晚上复习一个小时。我的汉语老师是从中国来的，她教我们说话、写字，也教我们唱中国歌。\\
学了一年以后，我现在可以说一点儿汉语，也会写一百多个汉字。明年我想参加汉语考试，以后我想去中国学习。\\
\emph{Wǒ jiào Ēndōng, wǒ shì Yǎwēndé yì suǒ zhōngxué de xuésheng, jīnnián shíliù suì. Wǒmen xuéxiào cóng qùnián kāishǐ yǒu Hànyǔ kè, wǒ měi ge xīngqī yǒu sān jié Hànyǔ kè.}\\
\emph{Wǒ měitiān zǎoshang liù diǎn qǐchuáng, qī diǎn qù xuéxiào. Hànyǔ kè hěn yǒu yìsi, kěshì Hànzì bù tài róngyì, suǒyǐ wǒ měitiān wǎnshang fùxí yí ge xiǎoshí. Wǒ de Hànyǔ lǎoshī shì cóng Zhōngguó lái de, tā jiāo wǒmen shuōhuà, xiězì, yě jiāo wǒmen chàng Zhōngguó gē.}\\
\emph{Xué le yì nián yǐhòu, wǒ xiànzài kěyǐ shuō yìdiǎnr Hànyǔ, yě huì xiě yìbǎi duō ge Hànzì. Míngnián wǒ xiǎng cānjiā Hànyǔ kǎoshì, yǐhòu wǒ xiǎng qù Zhōngguó xuéxí.}
\end{textebox}

\textbf{Questions de compréhension :}
\begin{enumerate}
\item Vrai ou faux ? Justifie avec une phrase du texte : 恩东的汉语老师是喀麦隆人。
\item Vrai ou faux ? Justifie avec une phrase du texte : 恩东觉得汉字很容易。
\item 恩东每天什么时候起床？什么时候去学校？
\item 他的汉语老师教他们做什么？（说出三件事）
\item 恩东明年想做什么？以后想做什么？
\end{enumerate}
\textbf{Questions de langue :}
\begin{enumerate}
\item[6.] Trouve dans le texte une phrase avec la structure 因为…所以… et recopie-la.
\item[7.] Trouve dans le texte un exemple du marqueur 了 et explique son emploi.
\end{enumerate}
\end{exercice}

\begin{exercice}[Thème et version — type Baccalauréat]
\textbf{A. Thème.} Traduis en chinois (caractères obligatoires) :
\begin{enumerate}
\item Je m'appelle Paul. Ma famille a quatre personnes.
\item Hier je n'ai pas regardé la télévision.
\item J'aime le chinois parce qu'il est intéressant.
\item Ma sœur est plus jeune que moi, elle a douze ans.
\end{enumerate}
\textbf{B. Version.} Traduis en français le texte suivant :
\begin{textebox}[Texte à traduire]
我姐姐每天都很忙。她早上六点起床，七点去医院工作。她很喜欢她的工作，因为工作很有意思。下班以后，她想学习汉语，因为她以后想去中国。\\
\emph{Wǒ jiějie měitiān dōu hěn máng. Tā zǎoshang liù diǎn qǐchuáng, qī diǎn qù yīyuàn gōngzuò. Tā hěn xǐhuan tā de gōngzuò, yīnwèi gōngzuò hěn yǒu yìsi. Xiàbān yǐhòu, tā xiǎng xuéxí Hànyǔ, yīnwèi tā yǐhòu xiǎng qù Zhōngguó.}
\end{textebox}
\end{exercice}

\begin{exercice}[Expression écrite et dictée — type Baccalauréat]
\textbf{A. Expression écrite.} Rédige un texte de \textbf{6 à 8 phrases} (minimum 60 caractères chinois) sur le sujet : « Mon école et mes loisirs » (\zh{我的学校和我的爱好}).
Consignes obligatoires :
\begin{itemize}
\item utilise au moins deux connecteurs : \zh{也} et \zh{因为…所以…} ;
\item utilise au moins un marqueur de temps : \zh{每天、以后、以前} ou \zh{的时候} ;
\item utilise au moins une structure de l'année : \zh{比}, \zh{得} (complément de degré) ou \zh{了} (passé) ;
\item écris en caractères chinois, avec soin ; le pinyin n'est pas exigé.
\end{itemize}
\textbf{B. Dictée de caractères.} Ton professeur dicte 10 mots de l'année. Pour chaque mot, écris les caractères et complète le pinyin avec les tons. Barème : 1 point par mot (0,5 pour les caractères, 0,5 pour les tons).
\begin{center}
\begin{tabularx}{\linewidth}{|X|X|X|}
\hline
\rowcolor{popblueL} N° & Caractères & Pinyin avec tons \\
\hline
1. 学校 & & \\
\hline
2. 朋友 & & \\
\hline
3. 喜欢 & & \\
\hline
4. 明天 & & \\
\hline
5. 医院 & & \\
\hline
\end{tabularx}
\end{center}
\textbf{C. Repérage d'erreurs.} Les 5 phrases suivantes contiennent chacune une erreur typique de francophone. Corrige-les et justifie chaque correction en citant la règle.
\begin{enumerate}
\item \zh{我没不昨天去学校。}
\item \zh{我喝一个茶。}
\item \zh{他比我很高。}
\item \zh{她说汉语说的很好。}
\item \zh{我昨天去商店了买东西。}
\end{enumerate}
\end{exercice}

\newpage

\coursec{Corrigés des exercices}

\begin{corrige}[Exercice 1 — QCM de grammaire]
\begin{enumerate}
\item \textbf{A. 是} : \zh{我是学生。} \emph{(Wǒ shì xuésheng.)} « Je suis élève. » Le verbe d'identité est \zh{是} ; \zh{在} indique la localisation, \zh{有} la possession.
\item \textbf{B. 没} : \zh{昨天我没去学校。} \emph{(Zuótiān wǒ méi qù xuéxiào.)} La négation du passé (action non accomplie) se fait avec \zh{没(有)}, jamais avec \zh{不}.
\item \textbf{B. 家…有} : \zh{他家有三口人。} \emph{(Tā jiā yǒu sān kǒu rén.)} « Sa famille compte trois personnes. » Structure : famille + \zh{有} + nombre + \zh{口} + \zh{人}.
\item \textbf{B. 一杯} : \zh{我喝一杯茶。} \emph{(Wǒ hē yì bēi chá.)} Le classificateur des boissons servies est \zh{杯} (verre/tasse).
\item \textbf{A. 了…本} : \zh{她买了两本书。} \emph{(Tā mǎi le liǎng běn shū.)} \zh{了} suit le verbe ; le classificateur des livres est \zh{本}.
\item \textbf{B. 每天…都} : \zh{我每天都七点起床。} \emph{(Wǒ měitiān dōu qī diǎn qǐchuáng.)} « Tous les jours » : \zh{每天…都}, \zh{都} se place avant le verbe.
\item \textbf{B. 高} : \zh{哥哥比我高。} \emph{(Gēge bǐ wǒ gāo.)} Dans la comparaison avec \zh{比}, l'adjectif seul suffit ; on ne met pas \zh{很} (et « 高得多一点儿 » mélange deux degrés).
\item \textbf{B. 得} : \zh{请你说得慢一点儿。} \emph{(Qǐng nǐ shuō de màn yìdiǎnr.)} Le complément de degré après le verbe s'introduit avec \zh{得}.
\item \textbf{A. 以前} : \zh{我以前去过北京。} \emph{(Wǒ yǐqián qùguo Běijīng.)} « Avant, autrefois » = \zh{以前} ; \zh{已经} exigerait un autre contexte, \zh{常常} marque l'habitude.
\item \textbf{B. 所以} : \zh{因为下雨，所以我们不去踢球。} \emph{(Yīnwèi xià yǔ, suǒyǐ wǒmen bù qù tī qiú.)} Paire de connecteurs \zh{因为…所以…}.
\end{enumerate}
\textbf{Barème indicatif :} 1 point par item, total 10 points.
\end{corrige}

\begin{corrige}[Exercice 2 — Vrai ou faux justifié]
\begin{enumerate}
\item \textbf{Faux.} \zh{我不没吃早饭。} est incorrect : on ne cumule jamais \zh{不} et \zh{没}. Pour un fait passé : \zh{我没吃早饭。} \emph{(Wǒ méi chī zǎofàn.)} « Je n'ai pas pris de petit-déjeuner. »
\item \textbf{Vrai.} Phrase correcte : \zh{也} « aussi » se place avant le verbe \zh{是}. « Il est mon professeur, et aussi mon ami. »
\item \textbf{Faux.} On n'emploie qu'un seul \zh{了} dans une phrase simple au passé : \zh{我昨天去商店买东西了。} \emph{(Wǒ zuótiān qù shāngdiàn mǎi dōngxi le.)} « Hier, je suis allé au magasin faire des courses. »
\item \textbf{Vrai.} Structure correcte du complément de degré avec objet : verbe + objet + verbe répété + \zh{得} + adjectif. « Elle parle très bien chinois. »
\item \textbf{Vrai.} Paire \zh{所以} bien employée après la cause ; \zh{想} + verbe exprime le projet. « Demain j'ai un examen, donc ce soir je veux réviser. »
\end{enumerate}
\textbf{Barème indicatif :} 2 points par item (1 pour vrai/faux, 1 pour la justification), total 10 points.
\end{corrige}

\begin{corrige}[Exercice 3 — Vocabulaire de l'année]
\begin{enumerate}
\item \zh{我头疼，要去看医生。} → \textbf{医院} : \zh{我头疼，要去医院看医生。} \emph{(Wǒ tóu téng, yào qù yīyuàn kàn yīshēng.)} « J'ai mal à la tête, je dois aller à l'hôpital voir le médecin. »
\item \textbf{爱好} : \zh{我的爱好是踢足球和听音乐。} \emph{(Wǒ de àihào shì tī zúqiú hé tīng yīnyuè.)} « Mes loisirs sont le football et la musique. »
\item \textbf{天气} : \zh{今天天气很好，我们去公园吧。} \emph{(Jīntiān tiānqì hěn hǎo, wǒmen qù gōngyuán ba.)} « Il fait beau aujourd'hui, allons au parc. »
\item \textbf{考试} : \zh{明天有汉语考试，我要复习。} \emph{(Míngtiān yǒu Hànyǔ kǎoshì, wǒ yào fùxí.)} « Demain il y a un examen de chinois, je dois réviser. »
\item \textbf{电视} : \zh{晚上我喜欢看电视。} \emph{(Wǎnshang wǒ xǐhuan kàn diànshì.)} « Le soir, j'aime regarder la télévision. »
\item \textbf{因为} : \zh{因为他病了，所以没来上课。} \emph{(Yīnwèi tā bìng le, suǒyǐ méi lái shàngkè.)} « Parce qu'il était malade, il n'est pas venu en cours. »
\end{enumerate}
\textbf{Barème indicatif :} 1 point par phrase, total 6 points.
\end{corrige}

\begin{corrige}[Exercice 4 — Pinyin et caractères]
Correspondances : \textbf{1-b, 2-c, 3-a, 4-d, 5-e, 6-f}.
\begin{center}
\begin{tabularx}{\linewidth}{|X|X|X|}
\hline
\rowcolor{popblueL} Caractères & Pinyin & Traduction \\
\hline
学 & xué & étudier \\
\hline
吃 & chī & manger \\
\hline
明天 & míngtiān & demain \\
\hline
朋友 & péngyou & ami \\
\hline
喜欢 & xǐhuan & aimer \\
\hline
中国 & Zhōngguó & la Chine \\
\hline
\end{tabularx}
\end{center}
\textbf{Points de vigilance :} \zh{学} porte le 2\textsuperscript{e} ton (xué), pas le 1\textsuperscript{er} ; \zh{吃} est au 1\textsuperscript{er} ton ; dans \zh{喜欢}, la deuxième syllabe est au ton neutre (huan) ; \zh{中国} prend la majuscule car c'est un nom propre.
\textbf{Barème indicatif :} 1 point par association, total 6 points.
\end{corrige}

\begin{corrige}[Exercice 5 — Traduction : du français au chinois]
\begin{enumerate}
\item \zh{我叫保罗。我家有四口人。} \emph{(Wǒ jiào Bǎoluó. Wǒ jiā yǒu sì kǒu rén.)} — Structure « famille + \zh{有} + nombre + \zh{口} + \zh{人} ».
\item \zh{昨天我没看电视。} \emph{(Zuótiān wǒ méi kàn diànshì.)} — Négation du passé avec \zh{没}, sans \zh{了}.
\item \zh{我喜欢汉语，因为汉语很有意思。} \emph{(Wǒ xǐhuan Hànyǔ, yīnwèi Hànyǔ hěn yǒu yìsi.)} — \zh{因为} introduit la cause ; adjectif prédicat précédé de \zh{很}.
\item \zh{我哥哥比我高。} \emph{(Wǒ gēge bǐ wǒ gāo.)} — Comparaison : A + \zh{比} + B + adjectif, sans \zh{很}.
\item \zh{下课以后，我想和朋友们踢足球。} \emph{(Xiàkè yǐhòu, wǒ xiǎng hé péngyoumen tī zúqiú.)} — \zh{以后} « après » se place après le mot ou la proposition ; \zh{和} « avec ».
\end{enumerate}
\textbf{Barème indicatif :} 2 points par phrase (1 pour la structure, 1 pour les caractères), total 10 points.
\end{corrige}

\begin{corrige}[Exercice 6 — Transformation de phrases]
\begin{enumerate}
\item \zh{我不喜欢喝茶。} \emph{(Wǒ bù xǐhuan hē chá.)} « Je n'aime pas boire le thé. » — \zh{不} devant le verbe pour un état ou une habitude.
\item \zh{他昨天没去市场。} \emph{(Tā zuótiān méi qù shìchǎng.)} « Hier, il n'est pas allé au marché. » — \zh{没} remplace la négation et \zh{了} disparaît.
\item \zh{她也学习法语。} \emph{(Tā yě xuéxí Fǎyǔ.)} « Elle étudie aussi le français. » — \zh{也} se place avant le verbe.
\item \zh{今天很热吗？} \emph{(Jīntiān hěn rè ma?)} « Fait-il chaud aujourd'hui ? » — Question fermée par ajout de \zh{吗} en fin de phrase.
\item \zh{你去过北京没有？} \emph{(Nǐ qùguo Běijīng méiyǒu?)} « Es-tu déjà allé à Pékin ? » — Question affirmative-négative : verbe + \zh{过} … \zh{没有}.
\end{enumerate}
\textbf{Barème indicatif :} 2 points par transformation, total 10 points.
\end{corrige}

\begin{corrige}[Exercice 7 — Texte à trous : mots-outils]
Texte complété :
\begin{textebox}[Texte corrigé]
我叫玛丽，我住\textbf{在}雅温得。我家有五口人。我\textbf{的}哥哥\textbf{比}我大三岁，他\textbf{也}学习汉语。去年他去\textbf{过}中国。他说汉语说\textbf{得}很好。\textbf{因为}汉语有意思，所以我每天学习两个小时。昨天我买\textbf{了}一本新词典。\\
\emph{Wǒ jiào Mǎlì, wǒ zhù zài Yǎwēndé. Wǒ jiā yǒu wǔ kǒu rén. Wǒ de gēge bǐ wǒ dà sān suì, tā yě xuéxí Hànyǔ. Qùnián tā qùguo Zhōngguó. Tā shuō Hànyǔ shuō de hěn hǎo. Yīnwèi Hànyǔ yǒu yìsi, suǒyǐ wǒ měitiān xuéxí liǎng ge xiǎoshí. Zuótiān wǒ mǎi le yì běn xīn cídiǎn.}
\end{textebox}
Justifications :
\begin{itemize}
\item \zh{在} : localisation avec le verbe \zh{住} « habiter ».
\item \zh{的} : possession « mon grand frère ».
\item \zh{比} : comparaison « trois ans de plus que moi ».
\item \zh{也} : « lui aussi étudie le chinois ».
\item \zh{过} : expérience passée « il est déjà allé en Chine ».
\item \zh{得} : complément de degré « il parle très bien ».
\item \zh{因为} : cause, en paire avec \zh{所以}.
\item \zh{了} : action accomplie « j'ai acheté ».
\end{itemize}
\textbf{Barème indicatif :} 1 point par mot-outil bien placé, total 8 points.
\end{corrige}

\begin{corrige}[Exercice 8 — Remise en ordre et appariement]
\textbf{Partie A.}
\begin{enumerate}
\item \zh{我每天早上七点起床。} \emph{(Wǒ měitiān zǎoshang qī diǎn qǐchuáng.)} « Je me lève tous les jours à sept heures du matin. » — Ordre : sujet + temps + heure + verbe.
\item \zh{杜阿拉比雅温得热。} \emph{(Dù'ālā bǐ Yǎwēndé rè.)} « Douala est plus chaude que Yaoundé. » — A + \zh{比} + B + adjectif.
\item \zh{他喝了两杯茶。} \emph{(Tā hē le liǎng bēi chá.)} « Il a bu deux tasses de thé. » — \zh{了} après le verbe, classificateur \zh{杯}.
\item \zh{因为下雨，所以我们不去学校。} \emph{(Yīnwèi xià yǔ, suǒyǐ wǒmen bù qù xuéxiào.)} « Parce qu'il pleut, nous n'allons pas à l'école. »
\end{enumerate}
\textbf{Partie B.} Correspondances : \textbf{1-b, 2-c, 3-a, 4-d}.
\begin{enumerate}
\item \zh{我每天上课以前先吃早饭。} « Avant les cours, je prends d'abord mon petit-déjeuner. »
\item \zh{下课以后，我和朋友们踢球。} « Après les cours, je joue au ballon avec mes amis. »
\item \zh{我学了汉语以后，想去中国看看。} « Après avoir appris le chinois, je veux aller voir la Chine. »
\item \zh{因为明天有考试，所以今天晚上我要复习。} « Parce que demain il y a examen, ce soir je dois réviser. »
\end{enumerate}
\textbf{Barème indicatif :} 2 points par phrase (partie A), 1 point par appariement (partie B), total 12 points.
\end{corrige}

\begin{corrige}[Exercice 9 — Compréhension écrite — type Baccalauréat]
\textbf{Questions de compréhension :}
\begin{enumerate}
\item \textbf{Faux.} Le texte dit : \zh{我的汉语老师是从中国来的。} \emph{(Wǒ de Hànyǔ lǎoshī shì cóng Zhōngguó lái de.)} « Mon professeur de chinois vient de Chine » : elle est donc chinoise, pas camerounaise.
\item \textbf{Faux.} Le texte dit : \zh{可是汉字不太容易} \emph{(kěshì Hànzì bù tài róngyì)} « mais les caractères ne sont pas très faciles ».
\item \zh{他每天早上六点起床，七点去学校。} \emph{(Tā měitiān zǎoshang liù diǎn qǐchuáng, qī diǎn qù xuéxiào.)} « Il se lève à six heures et va à l'école à sept heures. »
\item \zh{她教他们说话、写字，也教他们唱中国歌。} \emph{(Tā jiāo tāmen shuōhuà, xiězì, yě jiāo tāmen chàng Zhōngguó gē.)} « Elle leur apprend à parler, à écrire, et aussi à chanter des chansons chinoises. »
\item \zh{明年他想参加汉语考试，以后他想去中国学习。} \emph{(Míngnián tā xiǎng cānjiā Hànyǔ kǎoshì, yǐhòu tā xiǎng qù Zhōngguó xuéxí.)} « L'année prochaine, il veut passer l'examen de chinois ; plus tard, il veut aller étudier en Chine. »
\end{enumerate}
\textbf{Questions de langue :}
\begin{enumerate}
\item[6.] \zh{汉语课很有意思，可是汉字不太容易，所以我每天晚上复习一个小时。} (structure cause-conséquence avec \zh{所以} ; on accepte aussi toute phrase du texte contenant la relation cause-conséquence correctement recopiée).
\item[7.] Exemple : \zh{学了一年以后} \emph{(xué le yì nián yǐhòu)} : \zh{了} marque ici une action accomplie d'une durée déterminée (« après avoir étudié pendant un an »). Autre exemple accepté : \zh{我们学校从去年开始有汉语课} ne contient pas \zh{了} ; on peut aussi citer \zh{我学了汉语以后} si repris du texte de l'élève.
\end{enumerate}
\textbf{Barème indicatif :} questions 1-2 : 2 points chacune (1 pour vrai/faux, 1 pour la citation) ; questions 3-5 : 2 points chacune ; questions 6-7 : 2 points chacune ; total 14 points.
\textbf{Points de vigilance :} toujours justifier un « vrai/faux » par une citation exacte du texte ; recopier les caractères sans les modifier ; répondre par des phrases complètes.
\end{corrige}

\begin{corrige}[Exercice 10 — Thème et version — type Baccalauréat]
\textbf{A. Thème.}
\begin{enumerate}
\item \zh{我叫保罗。我家有四口人。} \emph{(Wǒ jiào Bǎoluó. Wǒ jiā yǒu sì kǒu rén.)}
\item \zh{昨天我没看电视。} \emph{(Zuótiān wǒ méi kàn diànshì.)}
\item \zh{我喜欢汉语，因为汉语很有意思。} \emph{(Wǒ xǐhuan Hànyǔ, yīnwèi Hànyǔ hěn yǒu yìsi.)}
\item \zh{我妹妹比我小，她十二岁。} \emph{(Wǒ mèimei bǐ wǒ xiǎo, tā shí'èr suì.)} — « Plus jeune » se dit \zh{小} dans la comparaison d'âge ; l'âge se donne sans verbe : nombre + \zh{岁}.
\end{enumerate}
\textbf{B. Version.} Traduction proposée :
« Ma grande sœur est très occupée tous les jours. Elle se lève à six heures du matin et va travailler à l'hôpital à sept heures. Elle aime beaucoup son travail, parce que son travail est très intéressant. Après le travail, elle veut apprendre le chinois, parce que plus tard elle veut aller en Chine. »
\textbf{Barème indicatif :} thème : 2 points par phrase (8 points) ; version : 6 points (fidélité 4, qualité du français 2) ; total 14 points.
\textbf{Points de vigilance :} en thème, ne pas oublier \zh{口} pour les membres de la famille et ne pas mettre \zh{很} dans la comparaison avec \zh{比} ; en version, \zh{下班以后} signifie « après le travail », \zh{以后} renvoie ici à « plus tard, dans le futur ».
\end{corrige}

\begin{corrige}[Exercice 11 — Expression écrite et dictée — type Baccalauréat]
\textbf{A. Expression écrite — proposition de rédaction modèle :}
\begin{textebox}[我的学校和我的爱好 — modèle]
我在雅温得的一所中学学习。我的学校不太大，可是很漂亮。我每天都七点去学校。我最喜欢汉语课，因为汉语很有意思。我的汉语老师很好，她教我们说话，也教我们写汉字。下课以后，我常常和朋友们踢足球。足球比篮球有意思，所以我每个星期踢三次。晚上我复习一个小时，然后看电视。\\
\emph{Wǒ zài Yǎwēndé de yì suǒ zhōngxué xuéxí. Wǒ de xuéxiào bù tài dà, kěshì hěn piàoliang. Wǒ měitiān dōu qī diǎn qù xuéxiào. Wǒ zuì xǐhuan Hànyǔ kè, yīnwèi Hànyǔ hěn yǒu yìsi. Wǒ de Hànyǔ lǎoshī hěn hǎo, tā jiāo wǒmen shuōhuà, yě jiāo wǒmen xiě Hànzì. Xiàkè yǐhòu, wǒ chángcháng hé péngyoumen tī zúqiú. Zúqiú bǐ lánqiú yǒu yìsi, suǒyǐ wǒ měi ge xīngqī tī sān cì. Wǎnshang wǒ fùxí yí ge xiǎoshí, ránhòu kàn diànshì.}
\end{textebox}
Ce modèle contient : \zh{也} (deux fois), \zh{因为…所以…}, les marqueurs de temps \zh{每天} et \zh{以后}, la comparaison \zh{比} et le complément de fréquence. Environ 110 caractères.
\textbf{Barème indicatif :} 10 points — respect des consignes (connecteurs, marqueur de temps, structure de l'année) : 4 points ; correction grammaticale : 3 points ; caractères et présentation : 2 points ; richesse et cohérence : 1 point.
\textbf{Points de vigilance :} ne pas traduire mot à mot le français ; placer les compléments de temps avant le verbe ; un seul \zh{了} par phrase au passé ; vérifier les tons si le pinyin est ajouté.

\textbf{B. Dictée de caractères — corrigé :}
\begin{center}
\begin{tabularx}{\linewidth}{|X|X|X|}
\hline
\rowcolor{popblueL} N° & Caractères & Pinyin avec tons \\
\hline
1. 学校 & 学校 & xuéxiào \\
\hline
2. 朋友 & 朋友 & péngyou \\
\hline
3. 喜欢 & 喜欢 & xǐhuan \\
\hline
4. 明天 & 明天 & míngtiān \\
\hline
5. 医院 & 医院 & yīyuàn \\
\hline
\end{tabularx}
\end{center}
Points de vigilance : \zh{校} a la clé du bois \zh{木} à gauche ; \zh{朋} est composé de deux \zh{月} ; \zh{喜} s'écrit de haut en bas ; dans \zh{喜欢} et \zh{朋友}, la deuxième syllabe est au ton neutre.

\textbf{C. Repérage d'erreurs :}
\begin{enumerate}
\item \zh{我没不昨天去学校。} → \zh{我昨天没去学校。} \emph{(Wǒ zuótiān méi qù xuéxiào.)} Règle : on ne cumule pas \zh{没} et \zh{不} ; le complément de temps se place après le sujet, avant le verbe.
\item \zh{我喝一个茶。} → \zh{我喝一杯茶。} \emph{(Wǒ hē yì bēi chá.)} Règle : chaque nom exige son classificateur ; pour une boisson servie, c'est \zh{杯}, pas \zh{个}.
\item \zh{他比我很高。} → \zh{他比我高。} \emph{(Tā bǐ wǒ gāo.)} Règle : dans la structure A + \zh{比} + B + adjectif, l'adjectif n'est pas précédé de \zh{很} (la comparaison suffit).
\item \zh{她说汉语说的很好。} → \zh{她汉语说得很好。} ou \zh{她说汉语说得很好。} \emph{(Tā shuō Hànyǔ shuō de hěn hǎo.)} Règle : le complément de degré s'introduit avec \zh{得}, pas \zh{的}.
\item \zh{我昨天去商店了买东西。} → \zh{我昨天去商店买东西了。} \emph{(Wǒ zuótiān qù shāngdiàn mǎi dōngxi le.)} Règle : dans une phrase à verbes en série, \zh{了} se place en fin de phrase (ou après le dernier verbe), pas après le premier.
\end{enumerate}
\textbf{Barème indicatif :} 2 points par correction (1 pour la phrase corrigée, 1 pour la règle citée), total 10 points.
\end{corrige}

\newpage

\coursec{Fiche bilan}

\begin{popfigure}
\begin{tikzpicture}[mindmap, grow cyclic, every node/.style={concept, text=white, font=\small}, concept color=popblue,
level 1/.append style={level distance=3.2cm, sibling angle=60},
level 2/.append style={level distance=1.8cm}]
\node[concept, font=\small\bfseries, align=center]{Révisions générales\\ 总复习}
  child[concept color=poporange] { node[align=center]{Vocabulaire\\ 词汇} }
  child[concept color=popgreen] { node[align=center]{Grammaire\\ 语法} }
  child[concept color=poppurple] { node[align=center]{Caractères\\ 汉字} }
  child[concept color=poppink] { node[align=center]{Compréhension\\ 阅读理解} }
  child[concept color=popgold] { node[align=center]{Traduction\\ 翻译} }
  child[concept color=popblue!70] { node[align=center]{Expression écrite\\ 写作} };
\end{tikzpicture}
\legende{Carte mentale de la leçon : les six compétences à mobiliser pour l'examen.}
\end{popfigure}

\begin{bilan}
\textbf{1. Lexique clé de l'année (à revoir par thème)}
\begin{itemize}
  \item \textbf{Salutations et identité :} \zh{你好} nǐ hǎo « bonjour », \zh{名字} míngzi « nom », \zh{国籍} guójí « nationalité », \zh{喀麦隆人} Kāmàilóngrén « Camerounais ».
  \item \textbf{Famille et école :} \zh{家人} jiārén « famille », \zh{同学} tóngxué « camarade de classe », \zh{老师} lǎoshī « professeur », \zh{学习} xuéxí « étudier ».
  \item \textbf{Vie quotidienne :} \zh{吃饭} chīfàn « manger », \zh{买东西} mǎi dōngxi « faire des achats », \zh{多少钱} duōshao qián « combien (prix) », \zh{市场} shìchǎng « marché ».
  \item \textbf{Temps et activités :} \zh{今天} jīntiān « aujourd'hui », \zh{周末} zhōumò « week-end », \zh{踢足球} tī zúqiú « jouer au football », \zh{看电影} kàn diànyǐng « regarder un film ».
  \item \textbf{Médias et communication (module 5) :} \zh{手机} shǒujī « téléphone portable », \zh{上网} shàngwǎng « aller sur Internet », \zh{发消息} fā xiāoxi « envoyer un message », \zh{打电话} dǎ diànhuà « téléphoner », \zh{新闻} xīnwén « informations ».
\end{itemize}

\textbf{2. Grammaire à connaître par cœur}
\begin{center}
\begin{tabularx}{\linewidth}{|X|X|X|}
\hline
\rowcolor{popblueL} Structure & Exemple (caractères + pinyin) & Traduction \\
\hline
Sujet + 很 + adjectif & \zh{雅温得很大。} Yǎwēndé hěn dà. & Yaoundé est grande. \\
\hline
Sujet + 是 + nom & \zh{我是喀麦隆人。} Wǒ shì Kāmàilóngrén. & Je suis Camerounais. \\
\hline
Sujet + 在 + lieu + verbe & \zh{我在学校学习。} Wǒ zài xuéxiào xuéxí. & J'étudie à l'école. \\
\hline
Verbe + 了 (action accomplie) & \zh{我吃了饭。} Wǒ chī le fàn. & J'ai mangé. \\
\hline
Verbe + 过 (expérience) & \zh{我去过中国。} Wǒ qùguo Zhōngguó. & Je suis déjà allé en Chine. \\
\hline
A + 比 + B + adjectif & \zh{汉语比法语难。} Hànyǔ bǐ Fǎyǔ nán. & Le chinois est plus difficile que le français. \\
\hline
Numéro + classificateur + nom & \zh{三个学生} sān ge xuésheng & trois élèves \\
\hline
Sujet + 会 / 想 / 要 + verbe & \zh{我想学汉语。} Wǒ xiǎng xué Hànyǔ. & Je veux apprendre le chinois. \\
\hline
Question avec 吗 / 呢 / 什么 & \zh{你好吗？} Nǐ hǎo ma? & Comment vas-tu ? \\
\hline
\end{tabularx}
\end{center}

\textbf{3. Mots-outils essentiels}
\begin{itemize}
  \item \zh{的} de : possession et qualification (\zh{我的书} wǒ de shū « mon livre »).
  \item \zh{了} le : action accomplie, changement d'état — jamais un simple « passé » mécanique.
  \item \zh{不} bù (négation présent/futur) ≠ \zh{没} méi (négation de \zh{有} et du passé).
  \item Classificateurs : \zh{个} ge (général), \zh{本} běn (livres), \zh{张} zhāng (objets plats), \zh{杯} bēi (verres).
\end{itemize}

\textbf{4. Caractères : pièges classiques}
\begin{itemize}
  \item Ne pas confondre : \zh{人} rén « personne » / \zh{入} rù « entrer » ; \zh{大} dà / \zh{太} tài ; \zh{未} wèi / \zh{末} mò.
  \item Ordre des traits : haut → bas, gauche → droite, horizontal avant vertical, extérieur avant intérieur, fermer la boîte en dernier.
  \item Reconnaître les clés : \zh{口} bouche, \zh{女} femme, \zh{氵} eau, \zh{心} cœur, \zh{讠} parole.
\end{itemize}

\textbf{5. Méthodes express pour l'examen}
\begin{itemize}
  \item \textbf{Compréhension :} lire les questions d'abord, repérer les mots connus, déduire le reste par le contexte.
  \item \textbf{Traduction (version) :} identifier sujet + verbe, traduire les compléments de temps et de lieu avant le verbe en français.
  \item \textbf{Thème :} respecter l'ordre chinois : Sujet + temps/lieu + verbe + complément ; vérifier les tons du pinyin.
  \item \textbf{Expression écrite :} 5 à 8 phrases simples et correctes valent mieux qu'une phrase longue et fautive.
\end{itemize}
\end{bilan}

\begin{aretenir}[Checklist avant l'examen]
\begin{itemize}
  \item $\square$ Je connais les salutations, nombres, dates et heures sans hésiter.
  \item $\square$ Je maîtrise l'ordre des mots : Sujet + (temps/lieu) + verbe + complément.
  \item $\square$ Je sais employer \zh{的}, \zh{了}, \zh{过}, \zh{不} et \zh{没} correctement.
  \item $\square$ Je connais les classificateurs courants (\zh{个}, \zh{本}, \zh{张}, \zh{杯}).
  \item $\square$ Je sais construire une comparaison avec \zh{比}.
  \item $\square$ Je forme les questions avec \zh{吗}, \zh{呢} et les mots interrogatifs.
  \item $\square$ J'écris les 80 caractères du programme dans le bon ordre des traits.
  \item $\square$ Je distingue les caractères proches (\zh{人}/\zh{入}, \zh{大}/\zh{太}).
  \item $\square$ Je mets les tons corrects dans le pinyin (ā á ǎ à).
  \item $\square$ Je relis ma production : tons, ponctuation chinoise （，。？）, caractères complets.
  \item $\square$ Je gère mon temps : compréhension, grammaire, traduction, rédaction.
  \item $\square$ Je reste calme : je commence par les questions que je maîtrise.
\end{itemize}
\end{aretenir}

\findecours

\end{document}
